\documentclass[11pt]{article}
\usepackage[T1]{fontenc}
\usepackage[utf8]{inputenc}
\usepackage{lmodern,amsmath,amssymb,amsthm,mathtools,mathrsfs}
\usepackage[margin=1.1in]{geometry}
\usepackage{microtype,enumitem}
\setlength{\emergencystretch}{2em}
\usepackage[unicode,colorlinks=true,linkcolor=blue,citecolor=blue,urlcolor=blue]{hyperref}
\usepackage{xurl}
\newtheorem{theorem}{Theorem}[section]
\newtheorem{lemma}[theorem]{Lemma}
\newtheorem{proposition}[theorem]{Proposition}
\newtheorem{corollary}[theorem]{Corollary}
\theoremstyle{definition}
\newtheorem{definition}[theorem]{Definition}
\newtheorem{remark}[theorem]{Remark}
\numberwithin{equation}{section}
\DeclareMathOperator{\Supp}{Supp}
\DeclareMathOperator{\rk}{rk}
\DeclareMathOperator{\ord}{ord}
\DeclareMathOperator{\Pic}{Pic}
\DeclareMathOperator{\Spec}{Spec}
\DeclareMathOperator{\Proj}{Proj}
\DeclareMathOperator{\Frac}{Frac}
\DeclareMathOperator{\vol}{vol}
\title{Metric descent and rank-preserving contractions}
\author{OpenAI}
\date{September 26, 2026}
\hypersetup{pdftitle={Metric descent and rank-preserving contractions},pdfauthor={OpenAI},pdfsubject={Metric-descent-and-rank-preserving-contractions-September-26-2026}}
\begin{document}
\maketitle
\begin{abstract}
From a smooth semipositive anticanonical metric on a smooth projective
complex variety, we construct an
unweighted integrable semipositive metric on a corrected anticanonical
line of a smooth base, using a normal equidimensional toroidal model
and full generic adjoint rank one at exponent zero. The explicit
rational boundary correction pulls back to an exceptional divisor.
On varieties without positive-degree holomorphic forms, two-metric
transfer gives sections with prescribed boundary poles from a
finite-volume log-anticanonical metric and a target line with bounded
semipositive weights. In particular, a smooth projective complex
variety with a smoothly semipositive anticanonical bundle and no such
forms has a nonzero invariant section of a positive anticanonical
multiple for every torus action with its natural linearization.
On this class of varieties, generic section and adjoint-rank hypotheses
give contractions of invariant fibrations that preserve both conditions
unless an invariant global section already exists.
\end{abstract}
\tableofcontents
\section{Introduction}\label{g:introduction}

Anticanonical nonvanishing asks when positivity of $-K_X$ produces a
nonzero section of $-mK_X$ for some integer $m>0$. We study metric
descent and dimension reduction for this problem, including sections
with prescribed poles and sections fixed by a torus.

Throughout the paper varieties are over $\mathbb C$. We write line
bundles additively. On a smooth variety $X$, the anticanonical bundle is
$-K_X=\det T_X$, with its natural linearization whenever a group acts.
A singular Hermitian metric is semipositive if its local squared frame
norm is $e^{-\varphi}$ for a plurisubharmonic weight $\varphi$.
A metric with locally bounded weights satisfies two-sided local bounds
on $\varphi$. Such a condition is stronger than pseudoeffectivity and
is preserved by the normalized maximum construction below.

Two section theorems are applications of the constructions developed
here. Let $X$ be smooth projective with no positive-degree holomorphic
forms, and let $D$ be an effective integral divisor. View its canonical
section $s_D$ as a $(-K_X+D)$-valued top form. If a semipositive metric
$h$ on $-K_X+D$ makes the measure $|s_D|_h^2$ finite, then for every
line bundle $L$ with locally bounded semipositive weights,
Theorem~\ref{g:allowed} gives integers $a>0$ and $c\ge0$ such that
\[
 H^0(X,aL+cD)\ne0.
\]
If $-K_X$ is smoothly semipositive and an algebraic torus $T$ acts on
$X$, Theorem~\ref{rb:invariant} gives
$H^0(X,-mK_X)^T\ne0$ for some $m>0$, for the natural action.
These conclusions are also established in the companion
\cite[Theorems~1.2 and~8.1]{HFiniteVolume}, which uses invariant sections
and finite-cover descent to obtain anticanonical nonvanishing without
the no-forms hypothesis \cite[Theorem~1.1]{HFiniteVolume}.
Here the emphasis is on reusable transfer statements, rank-preserving
contractions, and alternative constructions by orbit integrals and
finite-flat norms.

A section constructed on the base of a rational fibration need not
return to the original variety. Multiplying its pullback by a rational
relative section can leave poles over degenerate fibers, and a second
contraction can destroy the one-dimensional space that made an
integral metric positive. We study these two obstructions together:
the boundary corrections must disappear under birational pushdown,
and the contraction must preserve both a generic section and its
adjoint section ranks.

\subsection{From a volume on the total space to a metric on the base}

Let $X$ be smooth projective and suppose $L=-K_X$ has a smooth
semipositive metric. The metric determines a smooth positive volume
form on $X$. Resolve a rational fibration to obtain
\[
 X\xleftarrow{\ \pi\ } Z\xrightarrow{\ g\ }Y,
 \qquad E=K_{Z/X},
\]
with $X,Z,Y$ smooth projective. The adjoint identity is
\[K_{Z/Y}+\pi^*(-K_X)=E+g^*(-K_Y).\]
Over a smooth fiber, the Jacobian section
$s_E$ is an adjoint section after a base canonical frame has been
chosen. If it spans the generic adjoint space, integration along the
fibers turns the pushforward volume into the squared norm of a line.
Relative direct-image positivity then supplies the curvature inequality
on the smooth locus. The question is what happens at its boundary.

The first metric result, Theorem~\ref{g:toroidal-volume}, answers this question
on a normal equidimensional toroidal intermediate model
$W\to Y$ dominating $X$. For a base prime $Q$, let $e_D$ be its pullback
multiplicity at a prime $D$ of $W$ above $Q$, and let $a_D$ be the
Jacobian order over $X$. The correction is the explicit rational divisor
\[
 A_0=\sum_Q
 \max\left\{0,\min_{D\mapsto Q}
          \left(\frac{1+a_D}{e_D}-1\right)\right\}Q.
\]
Its pullback to $W$ is exceptional over $X$. The resulting metric on
$-K_Y+A_0$ has locally integrable squared frame norms without inserting
a vanishing section. If, away from the Jacobian divisor, the coefficient
curvature dominates a positive pulled-back base form on one patch,
then this metric is smooth and strictly positive on a nonempty ordinary
open set. The proof checks every divisorial valuation, including those
centered in intersections of the boundary. This construction is proved
before any section-existence argument. The companion
\cite[Sections~5.3 and~11.2]{EControlledBoundary} uses precisely this
unweighted integrability, strictness, and exceptional support for its
fixed-error effectivity argument and its Fano-type model of an auxiliary
base.

The distinction between a normal intermediate model and its smooth
resolution is useful. Equidimensionality controls which divisors can
carry poles. Smoothness permits fiber integration. A smooth resolution
need not be equidimensional, but its additional divisors over subsets of
base codimension at least two are exceptional over the original variety.
Keeping both models makes the final support statement exact.

\subsection{A contraction that keeps its generic section}

The second main result is Theorem~\ref{g:cf:contraction-theorem}.
Let $X$ have no positive-degree holomorphic forms, let $-K_X$ be smoothly
semipositive, and let a torus $T$ act on $X$. Consider an invariant
rational fibration to a trivially acted-on base $Y$ whose function field is relatively
algebraically closed in $\mathbb C(X)$. Assume that a positive
anticanonical multiple has a nonzero invariant section on the generic
fiber, and put $\eta=\operatorname{Spec}\mathbb C(Y)$. Suppose that the invariant generic adjoint spaces
\[
 H^0\bigl(Z_\eta,(K_{Z/X}+m\pi^*(-K_X))|_{Z_\eta}\bigr)^T
\]
have dimension one for every positive multiple of a fixed integer.
Then either a positive anticanonical multiple has a global invariant
section, or the fibration factors rationally through a strictly
lower-dimensional base satisfying the same two conditions.

The two assumptions have different roles. The generic section supplies
the line whose high moments yield a bounded base metric. Multiplication
by its powers also compares degree zero with the positive-degree
adjoint spaces and forces rank one at zero; this supplies the zeroth
density as well as the high moments. The rank condition makes ratios of sections into base
functions. After contraction, the proof must recover that statement
for the larger generic fiber, as well as producing its section.

Here is the mechanism. Vertical minima normalize the section to a
regular map from a base line. A toroidal density calculation produces
both a bounded metric on that line and a nef anti-log-canonical class
with a signed boundary. The negative boundary is supported on
those base primes where section poles can be absorbed upstairs.
If a suitable small boundary correction is effective, its section
lifts to $X$. Otherwise a movable class annihilates the relevant
divisors. Twisted differential forms give an unbounded sequence of
effective divisors of degree zero. Their ratios define a smaller
field. Affine relations among their horizontal multiplicities give the
new generic section, and the controlled base poles restore all the
adjoint ranks. This last step is what permits iteration.

\subsection{Methods, predecessors, and other metric hypotheses}

The distinction between numerical and linear effectivity is important
in positioning these results. Lazi\'c--Matsumura--Peternell--Tsakanikas--Xie
prove numerical effectivity of nef anti-log-canonical divisors for projective
log canonical threefold pairs under $\mathbb Q$-factoriality or rational
singularities \cite[Theorem~A]{LMPTX2023}. A numerical equivalence with
an effective divisor need not supply a section of a multiple of the
given line. M\"uller proves actual nonvanishing for projective klt
threefold pairs with nef anti-log-canonical divisor, and in arbitrary
dimension when that divisor is semiample on a general
maximal-rationally-connected fiber
\cite[Theorem~A and Corollary~B]{Muller2025}.
His equivariant theorem gives invariant plurisections for commutative
linear algebraic groups on projective sub-log-canonical pairs with
semiample anti-log-canonical divisor \cite[Theorem~C]{Muller2025}.
Our invariant conclusion instead uses smooth semipositivity on a
smooth variety without positive-degree forms and treats algebraic
tori. It does not assert the same generality of pairs or groups.
M\"uller's moment-polytope argument also precedes the toric initial-section
construction in Section~\ref{g:cf:fields}
\cite[Proposition~3.13 and Corollary~3.14]{Muller2025}.

The birational preparations use flattening by Raynaud--Gruson
\cite{RaynaudGruson1971}, the toroidalization of
Abramovich--Denef--Karu \cite{ADK2013}, and the cone subdivisions and
weak semistable reduction of Abramovich--Karu \cite{AK2000}.
Birational toroidal preparation and reduced fibers after alteration are
different constructions here. We use the former for valuation and
support estimates; the latter gives a separate norm construction in
Section~\ref{g:reduced-flat}.

The curvature mechanism builds on Berndtsson's smooth direct-image
positivity \cite[Theorem~1.2]{Berndtsson2009} and the relative Bergman
theorem of Berndtsson--P\u aun \cite[Theorem~0.1]{BerndtssonPaun2008}.
The singular direct-image formulation is due to P\u aun--Takayama
\cite[Theorem~3.3.5]{PT2018}; the actual-multiplier-ideal formulation is
also presented in \cite[Theorem~21.1]{HPS}. In the ordinary singular case
we verify integrability of the generating adjoint section, including
the actual multiplier ideal. In the invariant case smooth coefficients
allow compact averaging to give an orthogonal holomorphic quotient.
This explains why the singular and invariant forms of the transfer
statement have separate hypotheses. The normalized high-moment limit
then produces a bounded metric, while the zeroth moment retains a
finite-volume donor metric on a different line bundle.

Movable-class slopes connect those metrics to dimension reduction.
We use divisor--curve duality \cite{BDPP2013}, movable slope theory
\cite{GKP2016}, and the algebraicity criterion of Campana--P\u aun
\cite{CampanaPaun2019}. The foliation argument follows the strategy
used by Ou for generic nefness \cite{Ou2023}; the boundary and
integrability conditions required here are retained in each local
version. Proposition~\ref{g:bt:weighted-slope} proves the general finite-volume
cotangent estimate by approximating a rational boundary while preserving
its strict discrepancy inequalities. For integral boundaries, the companion
\cite[Theorem ``Finite-volume cotangent bound'']{HFiniteVolume} gives a
different proof through ramification estimates and relative Bergman
metrics. Hard Lefschetz with multiplier ideals
\cite[Theorem~0.1]{DPS2001} supplies twisted forms. Their determinant
extraction follows Lazi\'c--Peternell \cite[Lemma~4.1]{LP2018}, as adapted
in \cite[Lemma~5.1]{LMPTX2023}. The horizontal divisor comparison is
a relative version of the three-index argument in
\cite[Theorem~3.1, Step~3]{LP2018}: ratios in the smaller base field
give equality of horizontal parts instead of the absolute divisor
equality furnished there by Iitaka dimension zero. The subsequent
control of vertical poles and restoration of adjoint rank are separate
steps.

Several useful variants require different input data. An integrable
metric after multiplication by a boundary section need not have
integrable frame norms. A signed sub-klt divisor is also different
from an effective klt boundary. We therefore distinguish weighted,
unweighted, and signed-boundary transfer statements. The large-moment
construction uses only eventual adjoint ranks in its curvature step.
Together with its regular distinguished section those ranks imply
rank one also in the earlier degrees, but the construction does not
use zeroth-moment curvature. Each variant is carried
through its support and rank verification before it meets a common
contraction or induction.

Two further constructions illuminate the metric step. On a dense torus
orbit, logarithmic coordinates turn the volume density into an integral
of a log-concave function. This is the setting of complex Pr\'ekopa
positivity and Kiselman's minimum principle
\cite{Kiselman1978,Berndtsson1998,Berndtsson2006}.
The Brascamp--Lieb variance calculation
\cite{BrascampLieb1976,KolesnikovMilman2017} proves that the negative logarithms of its integral and its supremum
are plurisubharmonic in the base variables; Section~\ref{g:orbits} gives the complete
argument and its invariant-section application. On a reduced flat
alteration, a generic section instead determines a line nonzero on
every fiber. Its finite-flat determinant norm returns that line to a
birational model of the original base. The norm multiplies the section
exponent by its degree and preserves the required polar support.

\subsection{Reading the proofs}

Section~\ref{g:sec-toroidal-volume} proves the exceptional unweighted
metric construction. Section~\ref{g:analytic} gives two-metric transfer
and its rank-one metric hypotheses; Section~\ref{g:cf:contraction} then
proves the rank-preserving contraction. Section~\ref{g:bt} establishes
the weighted, unweighted, and signed-boundary variants.

For the ordinary allowed-pole theorem, the direct route is from the
two-metric transfer in Section~\ref{g:analytic}, through the weighted
cotangent estimate, to the induction in Section~\ref{g:reductions}.
For invariant sections, Section~\ref{g:orbits} applies that theorem
using either rank-one quotient integrals or explicit logarithmic orbit
integration. These proofs do not require the rank-preserving
contraction theorem. Sections~\ref{g:cf:fields}
and~\ref{g:field-variants} then identify the ratio fields, extract their
generic sections, and vary the base construction. The boundary
applications in Section~\ref{g:ba} and multiplicative systems in
Section~\ref{g:multiplicative} give other contractions with their own
rank hypotheses.
Section~\ref{g:reduced-flat} replaces the metric-transfer step by an
alteration and a finite-flat norm, then proves the rationally connected
cases of the ordinary and invariant conclusions over the original base
field.

\section{An integrable anticanonical metric on a toroidal base}
\label{g:sec-toroidal-volume}

A smooth metric on the anticanonical bundle of a smooth variety determines
a smooth positive volume form.  Integrating that form along a fibration
produces a density on the base.  The question here is whether this density
is the squared norm of a positively curved anticanonical metric, and
whether its singularities can be corrected without changing divisors on
the original variety.  Adjoint rank one answers the first question.
The comparison between ramification and Jacobian orders answers the
second.

We use additive notation for line bundles.  If a local frame has squared
norm $e^{-\varphi}$, its curvature is $i\partial\bar\partial\varphi$.
For a rational line bundle, choose a positive Cartier multiple and a
nonvanishing local holomorphic frame of that multiple, and divide its
weight by the chosen integer.  Integrability below means
$e^{-\varphi}\in L^1_{\mathrm{loc}}$ for these root weights.  Changing
the holomorphic frame multiplies the squared norm by a locally bounded
positive factor with locally bounded inverse, and therefore does not
change this condition.

A toroidal embedding is a normal variety with an open subset locally
modeled on a toric variety and its dense torus.  Its boundary is the
complement of that open subset.  We use strict toroidal embeddings,
whose irreducible boundary components are normal, and morphisms locally
modeled by toric morphisms.  The full reduced boundary will let us
control discrepancies even when the source variety is singular.

\begin{theorem}[Toroidal volume metric]\label{g:toroidal-volume}
Let $X$ and $Y$ be smooth connected projective complex varieties, with
$n=\dim X$ and $d=\dim Y>0$, and let $L=-K_X$ carry a smooth Hermitian metric $h$ of
semipositive curvature $\alpha$.  Suppose
\[
 p:W\longrightarrow X,\qquad f:W\longrightarrow Y
\]
are projective morphisms, where $W$ is normal and integral, $p$ is
birational, and $f$ is surjective, equidimensional, and toroidal for
strict toroidal embeddings.  Choose a smooth projective resolution
$\rho:Z\to W$, and put $\pi=p\rho$, $g=f\rho$, and $E_Z=K_{Z/X}$.
Assume
\begin{equation}\label{g:volume-rank-zero}
 \operatorname{rk}g_*\mathcal O_Z(E_Z)=1.
\end{equation}
For each prime divisor $Q\subset Y$ and each prime $E\subset W$
dominating $Q$, set
\[
 a_E=\operatorname{ord}_E K_{W/X},\qquad
 e_E=\operatorname{ord}_E f^*Q,
\]
where the Jacobian divisor $K_{W/X}=K_W-p^*K_X$ is defined in
codimension one.  Define
\begin{equation}\label{g:volume-correction}
 a_Q=\max\left\{0,\min_{f(E)=Q}
             \left(\frac{1+a_E}{e_E}-1\right)\right\},
 \qquad A_0=\sum_Q a_Q Q.
\end{equation}
Then $A_0$ is an effective rational divisor with finite support, and
$f^*A_0$ is $p$-exceptional.  The rational line bundle $-K_Y+A_0$
has a singular Hermitian metric of semipositive curvature whose squared
frame norms are locally integrable without a divisor weight.  This
metric is smooth on a nonempty Zariski open subset of $Y$.

Moreover, let $U\subset Y$ be a nonempty Zariski open set where $g$
is smooth and formation of the adjoint direct image
$g_*(K_{Z/Y}+\pi^*L)$ commutes with base change.  If there are a
nonempty ordinary open patch
$\mathcal V\subset g^{-1}(U)\setminus\operatorname{Supp}E_Z$,
a K\"ahler form $\omega_Y$, and a constant $c>0$ such that
\begin{equation}\label{g:volume-horizontal-patch}
 \pi^*\alpha\geq c\,g^*\omega_Y\qquad\text{on }\mathcal V,
\end{equation}
then the constructed metric is strictly positively curved on a
nonempty ordinary open subset of $Y$.
\end{theorem}

The rank in \eqref{g:volume-rank-zero} is the full generic adjoint rank
at exponent zero.  It identifies the entire direct image with a line
on a suitable open set.  In particular, the theorem does not obtain
positivity by restricting a direct-image metric to an arbitrary
subline.  Likewise, \eqref{g:volume-horizontal-patch} is a hypothesis
on one patch; its conclusion is strict curvature on an ordinary open,
not a global K\"ahler lower bound.

The theorem assumes a toroidal equidimensional model rather than
constructing one.  Birational weak toroidalization and projective
subdivisions provide such models in the applications; these are the
constructions of Abramovich--Denef--Karu and Abramovich--Karu
\cite{ADK2013,AK2000}.  No cover making the fibers reduced is part of
the theorem.  We retain the multiplicities $e_E$ in the correction.

\subsection{The correction and its discrepancies}

We first isolate the algebraic reason that the correction is disposable
on $X$ and integrable on every higher model.  The transverse map $y_1=u_1^{e_E}$ sends a density of order
$|u_1|^{2a_E}$ to one of order
$|y_1|^{2(a_E-e_E+1)/e_E}$. This explains the coefficient in
\eqref{g:volume-correction}. After the correction, on a normal-crossing
resolution $Z^\prime\to W$, with composite $g^\prime:Z^\prime\to Y$,
the upstairs density will have factors
$|z_j|^{2\gamma_j}$ with
$\gamma_j=\operatorname{coeff}_{\{z_j=0\}}
(K_{Z'/X}-(g')^*A_0)$. Such factors are integrable exactly when
$\gamma_j>-1$. New exceptional divisors can be centered in boundary
intersections, so checking only the primes already on $W$ is insufficient.
The following lemma controls all the new exponents by comparing with the
full toroidal boundary. The canonical divisors in
this discussion are chosen compatibly with a rational top form on
$X$; all Jacobian differences are consequently unambiguous.

For a rational Weil divisor $\Theta$ with $K_W+\Theta$ rationally
Cartier, its crepant pullback to a smooth model $\tau:W'\to W$ is
defined by $K_{W'}+\Theta_{W'}=\tau^*(K_W+\Theta)$.
The pair is \emph{sub-klt} if every coefficient of $\Theta_{W'}$
is less than one on every such model; negative coefficients are
allowed.  Equivalently, every divisorial log discrepancy
$A_{W,\Theta}(\operatorname{ord}_F)=1-\operatorname{coeff}_F\Theta_{W'}$
is positive.  Replacing ``positive'' by ``nonnegative'' defines
sub-log-canonicity, or log canonicity when the boundary is effective.

\begin{lemma}\label{g:volume-subklt}
Let $X,Y$ be smooth connected projective varieties with $\dim Y>0$,
let $W$ be normal and integral, and let $p:W\to X$ be projective
birational and $f:W\to Y$ projective, surjective, equidimensional,
and toroidal for strict toroidal embeddings. Define $A_0$ by
\eqref{g:volume-correction}, and put
\[
 \Theta=-K_{W/X}+f^*A_0.
\]
Then $\Theta\leq0$, the pair $(W,\Theta)$ is sub-klt, and
$f^*A_0$ is $p$-exceptional.  In particular, for every smooth model
$\tau:Z'\to W$,
\begin{equation}\label{g:volume-all-valuations}
 \operatorname{coeff}_F\bigl(K_{Z'/X}-(f\tau)^*A_0\bigr)>-1
\end{equation}
at every prime $F\subset Z'$.
\end{lemma}

This assertion uses only the morphisms and their divisor orders.
Neither a metric nor an adjoint-rank assumption is needed; this will
allow us to use the same discrepancy estimate with invariant direct
images later.

\begin{proof}
At a general point of a prime of $W$, the variety is smooth and the
Jacobian of the morphism to $X$ is holomorphic.  Hence $a_E\geq0$.
At a prime that is not exceptional over $X$, the morphism is an
isomorphism at its general point, so $a_E=0$.

Every vertical prime of $W$ dominates a prime of $Y$.  Indeed, an
equidimensional map of relative dimension $\dim W-\dim Y$ cannot map a
divisor into a set of codimension at least two: its dimension would
then be at most $\dim W-2$.  If $a_Q>0$, the minimum in
\eqref{g:volume-correction} is positive, so every prime over $Q$ has
$a_E>0$ and is $p$-exceptional.  This proves the asserted support of
$f^*A_0$.  It also proves finite support, since there are only finitely
many $p$-exceptional prime divisors.  Rationality follows from the
integer orders in \eqref{g:volume-correction}.

If $E$ dominates $Q$ with $a_Q>0$, its coefficient in $\Theta$
satisfies
\[
 -a_E+e_Ea_Q\leq 1-e_E\leq0.
\]
At a prime above a $Q$ with $a_Q=0$, and at any horizontal prime,
the coefficient is $-a_E\leq0$.  Thus $\Theta\leq0$.

Let $D_W$ be the full reduced toroidal boundary.  The toroidal pair
$(W,D_W)$ is log canonical, $K_W+D_W$ is rationally Cartier, and
$W\setminus D_W$ is smooth.  These assertions can be checked on a
regular subdivision of each toric chart: the logarithmic top form
gives
$K_{\widetilde W}+D_{\widetilde W}=\tau^*(K_W+D_W)$,
and the reduced boundary upstairs has simple normal crossings.
We also have
\[
 K_W+\Theta=p^*K_X+f^*A_0.
\]
Consequently $B=D_W-\Theta$ is an effective rational Cartier divisor
whose support contains $D_W$.  For any divisorial valuation $v$ over
$W$, the log discrepancies satisfy
\begin{equation}\label{g:volume-discrepancy-comparison}
 A_{W,\Theta}(v)=A_{W,D_W}(v)+v(B).
\end{equation}
The full toroidal boundary gives nonnegative log discrepancies;
the effective rational Cartier difference adds a strictly positive term precisely
where that discrepancy could vanish. Both terms on the right are nonnegative.  If the first is zero, the
center of $v$ lies in $D_W$, since the complement is smooth with zero
boundary.  Its center is then contained in the support of the
effective rational Cartier divisor $B$, so $v(B)>0$.  If the first
term is positive there is nothing more to check.  This proves
$A_{W,\Theta}(v)>0$ for every divisorial valuation, which is the
sub-klt assertion.  This argument does not require the support of
$K_{W/X}$ to lie in the toroidal boundary.

Finally the crepant pullback of $\Theta$ to $Z'$ is
\[
 \Theta_{Z'}=(f\tau)^*A_0-K_{Z'/X}.
\]
Its coefficients are less than one by the sub-klt assertion, giving
\eqref{g:volume-all-valuations}.
\end{proof}

\subsection{Positivity of the integrated density}

We next explain why the volume density gives a semipositive metric,
and why the patch in \eqref{g:volume-horizontal-patch} gives strict
curvature.  The analytic input is the smooth direct-image theorem
of Berndtsson \cite[Theorem~1.2]{Berndtsson2009}.  We use its
primitive-representative curvature formula before the additional
product and fiberwise-positivity assumptions considered later in that
paper.

For clarity, consider a proper smooth K\"ahler family
$f:\mathcal Z\to\mathbb D$ over a disc with coordinate $t$,
relative dimension $r>0$, and a smooth semipositive
coefficient line $(\mathcal L,h_{\mathcal L})$.  Assume formation of
$\mathcal E=f_*(K_{\mathcal Z/\mathbb D}+\mathcal L)$ commutes with
base change.  Choose a holomorphic section $u$ of $\mathcal E$
that is Chern-normal at $0$.  Berndtsson's representative construction
\cite[Proposition~4.2]{Berndtsson2009} gives an absolute
$\mathcal L$-valued $(r,0)$-form $\widehat u$ representing $u$, with
\[
 \bar\partial\widehat u=dt\wedge\eta,\qquad
 \partial^{h_{\mathcal L}}\widehat u=dt\wedge\mu,
 \qquad \mu|_{\mathcal Z_0}=0,
 \qquad (\omega\wedge\eta)|_{\mathcal Z_0}=0.
\]
Here $\omega$ is a K\"ahler form on the total space.  Differentiating
the fiber integral with this representative gives
\begin{equation}\label{g:volume-curvature-bound}
 -i\partial\bar\partial\|u\|^2\big|_0
 \ \geq\
 f_*\!\left(i^{r^2}\widehat u\wedge\overline{\widehat u}
       h_{\mathcal L}\wedge i\Theta_{\mathcal L,h_{\mathcal L}}\right)
       \big|_0.
\end{equation}
The omitted term is the nonnegative squared norm of the primitive
$(r-1,1)$-form $\eta$; this is
\cite[formula~(4.8)]{Berndtsson2009}.
The displayed integrand is nonnegative because an $(r,0)$-form in
dimension $r+1$ is decomposable.  If the coefficient curvature
dominates $c\,f^*(i\,dt\wedge d\bar t)$ on a patch where $u(0)$
is nonzero, the right side is strictly positive: its integral on that
patch is bounded below by the positive integral of the squared
relative form times $c\,i\,dt\wedge d\bar t$.
Applying this argument on discs in every base tangent direction proves
horizontal strictness.  When $r=0$ and the full direct image has rank
one, the smooth map has one-point fibers, and its line metric is the
coefficient metric itself.

We now apply this calculation to the theorem.  Write $\Omega_h$ for
the smooth positive volume form on $X$ determined by $h$.  In a
coordinate frame $\partial/\partial x_1\wedge\cdots\wedge
\partial/\partial x_n$ of $-K_X$, it has density equal to that
frame's squared norm.  Its pullback to $Z$ contains the squared
Jacobian factor and vanishes along $E_Z$ to twice its order.

Choose a nonempty Zariski open $U$ over which $g$ is smooth and
base change holds.  After shrinking $U$, the canonical section
$s_{E_Z}$ induces an
isomorphism
\begin{equation}\label{g:volume-adjoint-identification}
 g_*(K_{Z/Y}+\pi^*L)|_U
 =g_*\mathcal O_Z(E_Z)|_U\otimes(-K_Y)|_U
 \simeq(-K_Y)|_U.
\end{equation}
Indeed, the first direct image on the right is generically a line
by \eqref{g:volume-rank-zero}, and $s_{E_Z}$ is a nonzero generic
section of that line.  Local freeness, base change, and its
nonvanishing as a frame hold on a smaller open set.  When
\eqref{g:volume-horizontal-patch} is assumed, the chosen ordinary
patch still meets its inverse image.  We also remove
$\operatorname{Supp}A_0$ from $U$.

In coordinates $y=(y_1,\ldots,y_d)$ on $U$, let $dV_y$ denote the
Euclidean volume with the same normalization used for top forms, and
write
\begin{equation}\label{g:volume-density}
 g_*\pi^*\Omega_h=H(y)\,dV_y.
\end{equation}
By \eqref{g:volume-adjoint-identification}, $H$ is exactly the squared
$L^2$ norm of the corresponding frame of $-K_Y$.  It is smooth and
positive on $U$, and $-\log H$ is plurisubharmonic by smooth
direct-image positivity.  If \eqref{g:volume-horizontal-patch} holds,
the section corresponding to that frame is
nonzero wherever the Jacobian section is nonzero.  Therefore
\eqref{g:volume-curvature-bound}, applied with coefficient
$\mathcal L=\pi^*L$, and
\eqref{g:volume-horizontal-patch} make $-\log H$ strictly
plurisubharmonic at an image point of the patch.  Smoothness then gives
strict positivity on an ordinary neighborhood of that point.

\subsection{Extension and unweighted integrability}

The remaining task is to cross the omitted base locus.  Extension of a
plurisubharmonic weight requires an upper bound; integrability requires
an upper bound on its exponential with the opposite sign.  The
ramification calculation provides the first, while the discrepancy
calculation provides the second.

On a coordinate neighborhood in $Y$, choose local equations $t_Q$ for
the finitely many components of $A_0$.  On its intersection with $U$
define
\begin{equation}\label{g:volume-weight}
 \Phi=-\log H+\sum_Q a_Q\log|t_Q|^2.
\end{equation}
These are the weights of a metric on $-K_Y+A_0$: its squared frame
norm is $H\prod_Q|t_Q|^{-2a_Q}$.  The assertion follows either by
the frame transition rule or after clearing the denominators of
$A_0$.  On $U$ the added terms are pluriharmonic, so the metric has
the curvature already established.

Let $Q$ be a prime component of $Y\setminus U$.  Choose a prime $E$
over $Q$ attaining the minimum in
\eqref{g:volume-correction}.  At general points of $E$ and $Q$, both
are smooth and $E\to Q$ is a submersion.  Choose the resolution to be
an isomorphism there.  After absorbing a nonvanishing unit into a
local coordinate, the map has the form
\[
 (u_1,\ldots,u_n)\longmapsto
 (u_1^{e_E},u_2,\ldots,u_d),
 \qquad Q=\{y_1=0\}.
\]
On a smaller coordinate polydisc the density of $\pi^*\Omega_h$ is
bounded below by a positive constant times $|u_1|^{2a_E}$.
Integrate over one fixed polydisc in the remaining fiber coordinates.
The change of variables $y_1=u_1^{e_E}$, whose real Jacobian is
$e_E^2|u_1|^{2(e_E-1)}$, yields
\begin{equation}\label{g:volume-density-lower-bound}
 H(y)\geq c_Q\,|y_1|^{2b_Q},
 \qquad b_Q=\min_{f(E)=Q}
                  \left(\frac{1+a_E}{e_E}-1\right).
\end{equation}
All other local factors in \eqref{g:volume-weight} are bounded above
and below near the chosen general point.  Since
$b_Q-a_Q=b_Q-\max\{0,b_Q\}\leq0$, the lower bound in
\eqref{g:volume-density-lower-bound} makes $\Phi$ bounded above
there.  The removable-singularity theorem for plurisubharmonic
functions extends it across a general open subset of each omitted
divisor.

The remaining omitted analytic set has codimension at least two in
the smooth base.  The usual Hartogs argument supplies the upper
bound needed there as well.  In a small polydisc, choose one
coordinate direction with circle boundary disjoint from that set.
For transverse parameters outside its projection the entire disc
misses the set, so the submean inequality bounds $\Phi$ by its bound
on the circle.  The exceptional parameter set has measure zero;
applying the submean inequality on a small ball about any remaining
point where $\Phi$ is defined gives the same bound there.  Taking the upper
semicontinuous extension proves removability.  The extensions agree
in overlapping frames by uniqueness, and define a semipositive
metric on $-K_Y+A_0$.

To prove integrability, use the projection formula on the smooth
locus, rather than pulling back the already integrated density:
\begin{equation}\label{g:volume-projection-formula}
 e^{-\Phi}\,dV_y
 =g_*\left(\prod_Q|t_Q\circ g|^{-2a_Q}\,
                      \pi^*\Omega_h\right).
\end{equation}
Take a further resolution $Z'\to W$ on which the Jacobian divisor
and the pullback of $A_0$ have simple normal crossings, and denote
its maps by $\pi':Z'\to X$ and $g':Z'\to Y$.  On a normal-crossing
coordinate chart the upstairs measure on the right of
\eqref{g:volume-projection-formula} is a positive bounded smooth
factor times
\[
 \prod_{j=1}^s |z_j|^{2\gamma_j}\,dV_z,
 \qquad
 \gamma_j=\operatorname{coeff}_{\{z_j=0\}}
                 \bigl(K_{Z'/X}-(g')^*A_0\bigr)>-1,
\]
where the strict inequalities are
Lemma~\ref{g:volume-subklt}.  Each one-variable integral
$\int_0^\varepsilon r^{2\gamma_j+1}\,dr$ is finite, so this measure
is locally integrable on $Z'$.  Properness makes its pushforward
locally finite on $Y$.  Formula~\eqref{g:volume-projection-formula}
then proves $e^{-\Phi}\in L^1$ off the omitted analytic set, with
finite integral on every relatively compact coordinate neighborhood.
The extension changes the weight only on a set of Lebesgue measure
zero and therefore has the same integrability.  This proves
Theorem~\ref{g:toroidal-volume}.

\begin{corollary}\label{g:correction-iitaka-zero}
For the correction in Theorem~\ref{g:toroidal-volume},
$\kappa(Y,A_0)=0$.  More precisely, for every positive integer $m$
for which $mA_0$ is integral,
\[
 H^0(Y,\mathcal O_Y(mA_0))
   =\mathbb C\,s_{mA_0}.
\]
\end{corollary}

\begin{proof}
A section divided by $s_{mA_0}$ is a rational function $q$ on $Y$
with poles bounded by $mA_0$.  The rational function $f^*q$ on
$W$ has only $p$-exceptional poles.  Viewed in
$\mathbb C(W)=\mathbb C(X)$, it consequently has no pole along any
prime of the smooth variety $X$.  Normality extends it regularly
across the remaining subset of codimension at least two.  It is
constant because $X$ is connected and projective.  Dominance of
$f$ makes $q$ constant as well.  The canonical section exists since
$A_0\geq0$, giving the asserted one-dimensional space and Iitaka
dimension.
\end{proof}

\section{Two metrics and controlled pole support}\label{g:analytic}\label{rb:two-metrics}

The section problem involves two line bundles with different metric
requirements.  The first supplies a finite volume measure and controls
cotangent tensors.  The second is the bundle in which we seek a section;
its bounded metric is preserved when the problem is transferred to a
lower-dimensional variety.

All varieties in this section are smooth connected projective varieties
 over $\mathbb C$, unless a normal intermediate model is specified.
We use additive notation for line bundles.  A semipositive singular
Hermitian metric has local squared norm $e^{-\phi}$, where $\phi$ is
plurisubharmonic and not identically $-\infty$.  Bounded weights are
locally bounded from both sides.  Let $D\geq0$ be an integral divisor and
put $A=-K_V+D$.  Its canonical section $s_D$ is an $A$-valued top form.
For a metric $h_A$ on $A$, the notation
\[
                 \mu=|s_D|_{h_A}^{2}
\]
denotes the associated positive measure.  The finite-volume condition
$\int_V\mu<\infty$ is a separate assumption; it does not follow from
semipositivity of $h_A$.  This measure has positive finite density almost
everywhere and gives no mass to proper algebraic subsets.

\subsection{Models and adjoint measures}\label{g:models}

For a dominant rational map with connected general fibres we use a
normal projective intermediate model $\bar V$ and a smooth resolution
$W$ fitting into
\[
 W\longrightarrow\bar V,\qquad
 \bar p:\bar V\longrightarrow V,\quad
 \bar f:\bar V\longrightarrow Y,
 \qquad p:W\longrightarrow V,\quad f:W\longrightarrow Y.
\]
Here $Y$ is smooth projective, $\bar p$ is birational, and $\bar f$ is
surjective and equidimensional.  Such models follow by applying
flattening to the graph, resolving the modified base, taking its
dominating transform, and normalizing \cite{RaynaudGruson1971}.
The dimension bound from the flat model persists under base change for
the dominating transform and under finite normalization.  Every
vertical prime of $\bar V$ therefore maps onto a prime divisor of $Y$.
A prime of $W$ mapping into a subset of $Y$ of codimension at least two
is exceptional over $\bar V$ and hence over $V$.  The smooth resolution
itself need not be equidimensional.

Write
\[
 J=K_{W/V}\geq0,\qquad S=J+p^*D.
\]
The section $s_S=s_Jp^*s_D$ is a section of $K_W+p^*A$.  Its squared
norm is the birational pullback of $\mu$, including the squared Jacobian.
Change of variables off the exceptional locus gives a finite measure
on $W$.  This statement concerns densities almost everywhere and does
not require a pointwise definition of a singular metric on its polar
set.

\subsection{Transferring the two metrics}

The next statement separates normalization of divisor orders from the
analytic input which gives curvature.  This permits both a rank-one
adjoint argument and a direct calculation along torus orbits.

\begin{lemma}[Transfer with controlled pole support]\label{g:transfer}\label{g:transfer-recalled}\label{rb:transfer}
Let $V$ be smooth connected projective, let $D\geq0$ be integral,
and suppose $A=-K_V+D$ has a semipositive singular metric $h_A$ with
$\int_V|s_D|_{h_A}^2<\infty$. Let $L$ have a semipositive metric $h_L$
with bounded weights.  Use the model diagram above with
$0<\dim Y\leq\dim V$.  Suppose a nonzero rational section $\rho$ of $kL$,
$k>0$, has no horizontal poles on $\bar V$.  On a dense base open set
put
\[
 g_0(y)\,d\lambda_y=f_*p^*\mu,
 \qquad
 S_\rho(y)=\operatorname*{ess\,sup}_{W_y}|p^*\rho|_{h_L^k}^2.
\]
Assume $-\log g_0$ in coordinate frames and $-\log S_\rho$ have psh
representatives on this open set.  Then, after clearing a positive
integer denominator, there are a line bundle $B$ with bounded
semipositive weights, an effective integral divisor $D_Y$, and a
regular nonzero map
\[
                   f^*B\longrightarrow bk\,p^*L
\]
for some $b>0$.  The line bundle $-K_Y+D_Y$ has a semipositive metric
whose distinguished adjoint measure has finite mass.  Every nonzero section
of $aB+cD_Y$, where $a>0$ and $c\geq0$ are integers, produces a
nonzero rational section of $abkL$ with poles only on $D$.
If the data and $\rho$ are invariant for a torus acting trivially on
$Y$, then the produced section is invariant.
\end{lemma}

\begin{proof}
We first normalize the section, obtaining both the regular map and its
bounded base metric. We then correct the pushforward measure along base
primes whose poles disappear on return to $V$ or remain within $D$.

For a prime $T$ of $Y$ set
\[
 d_T=\min_{Q\mapsto T}
       \frac{\operatorname{ord}_Q(\bar p^*\rho)}
            {\operatorname{ord}_Q(\bar f^*T)},
 \qquad B_0=\sum_Td_TT.
\]
The orders in this expression are taken on $\bar V$.  Only finitely
many $d_T$ are nonzero.  Equidimensionality and the absence of
horizontal poles give
\[
             \operatorname{div}_{\bar V}(\rho)-\bar f^*B_0\geq0.
\]
This is an effective $\mathbb Q$-Cartier divisor.  Choose $b>0$ clearing
the coefficients of $B_0$ and set $B=bB_0$.  Normality extends the
normalized section in codimension one to a regular morphism
$\bar f^*B\to bk\bar p^*L$, and pullback gives the displayed map.
Over each prime $T$ a component realizing the minimum is absent from
its vanishing divisor.

For a local meromorphic equation $l$ of $B$, define the weight
\[
                     -b\log S_\rho+\log|l|^2.
\]
It is psh on the initial open set.  Its negative exponential is the
fibre supremum of the squared norm of the regular normalized section
$(p^*\rho)^b/f^*l$.  Properness and bounded weights of $h_L$ bound this
norm above locally over every base point, giving a lower bound for the
weight even near the omitted locus. At a general point of each
omitted base prime $T$ choose a minimizing component. The normalized section
is nonzero there. At general points of this component and $T$, the
transverse map has coordinates $y_1=z_1^e$, $y_i=z_i$ for
$2\leq i\leq\dim Y$. Thus every nearby fibre meets a patch where
the normalized section norm is bounded below. The fibre
measure has positive density almost everywhere on such a patch, so
the essential supremum has a fixed positive lower bound.  This gives
an upper bound for the weight near a general point of that prime.
Psh removable singularities extend the weight there and then across
the remaining analytic set of codimension at least two.  The lower
bound persists, while the extended psh weight is locally bounded
above.  These weights glue and give the bounded metric on $B$.

The bounded metric is now constructed. For the finite-volume metric,
call a base prime $T$ invisible if every component of $\bar f^*T$ is
exceptional over $V$ or maps into $D$.  There are finitely many such
primes, since they are images of the finitely many exceptional or
boundary primes. Choose a component above $T$ on $\bar V$ and denote
its strict transform on $W$ by $Q$. At its general point write
\[
 e=\operatorname{ord}_Q(f^*T),\qquad
 a=\operatorname{ord}_Q(J+p^*D).
\]
The local psh coefficient weight of $h_A$ is bounded above, so its
negative exponential is bounded below by a positive constant. The
Jacobian and boundary section therefore bound the pulled-back measure
below by a positive constant times $|z_1|^{2a}$ times coordinate volume.
Integrating over
a fixed disc in the free variables and using $y_1=z_1^e$ gives
\begin{equation}\label{rb:density-lower}
              g_0(y)\geq C|y_1|^{2(a-e+1)/e}
\end{equation}
almost everywhere near the general base point.  If $T$ is visible,
choose a component with $a=0$; then $-\log g_0$ is bounded above there.
For the invisible primes choose sufficiently large integral
coefficients of $D_Y$.  Equation~\eqref{rb:density-lower} makes
\[
                  -\log g_0+\log|s_{D_Y}|^2
\]
locally bounded above at their general points.  It extends psh in
codimension one and then across codimension two, defining a metric on
$-K_Y+D_Y$.  Pairing it with $s_{D_Y}$ cancels the added factor.  Its
adjoint measure is $g_0d\lambda_y=f_*p^*\mu$ almost everywhere and is
therefore finite.

A section of $aB+cD_Y$ is a rational section of $aB$ with poles only on
$D_Y$.  Pull it back and compose with the $a$-th power of the normalized
map.  On $\bar V$ all possible new poles lie above invisible primes.
They are thus exceptional over $V$ or map into $D$.  The normal
intermediate model excludes nonexceptional divisors over base subsets
of codimension at least two.  Birational pushdown therefore leaves
poles only on $D$, and a sufficiently large multiple of $D$ clears
them.  The rational section remains nonzero.  The normalization uses
only base factors, so it preserves invariance when present.
\end{proof}

\begin{lemma}[Rank-one sufficient condition]\label{g:moments}\label{rb:rank-one}
In the setting of Lemma~\ref{rb:transfer}, suppose on the generic fibre
of $f$ that
\[
 h^0\bigl(W_\eta,(J+p^*D+mkp^*L)|_{W_\eta}\bigr)=1
                       \qquad(m\geq0).
\]
Then the two psh assumptions in that lemma hold.  There is also an
invariant version: take $D=0$, smooth input metrics invariant under a
compact form of an acting torus, and an invariant $\rho$; it suffices
that the invariant dimensions in the display equal one.
\end{lemma}

\begin{proof}
On a base open set where $\rho$ is regular, define
\[
 g_m(y)\,d\lambda_y
       =f_*\bigl(|p^*\rho|_{h_L^k}^{2m}p^*\mu\bigr).
\]
These are finite positive densities almost everywhere, since the
section norms are bounded on compact preimages and $p^*\mu$ has finite
mass. Give $p^*(A+mkL)$ the coefficient metric
$h_m=p^*(h_Ah_L^{mk})$, and consider the integrable direct image
\[
 \mathcal F_m=f_*\bigl(\mathcal O_W(K_{W/Y}+p^*(A+mkL))
                         \otimes\mathcal I(h_m)\bigr).
\]
For a local canonical frame $\eta_Y$ on the base, the section
\[
             u_m=\frac{s_{J+p^*D}(p^*\rho)^m}{f^*\eta_Y}
\]
is locally integrable on the total space and therefore belongs to
$\mathcal F_m$. It is nonzero and generates
the ordinary direct image generically by the rank-one hypothesis.
Consequently $\mathcal F_m$ also has rank one. On a dense open set
where $u_m$ is its frame, the squared direct-image norm is $g_m$.

Singular direct-image positivity gives $-\log g_m$ a psh
representative there. We use the formulation with the actual
multiplier ideal in \cite[Theorem~21.1]{HPS}, based on
\cite[Theorem~3.3.5]{PT2018}. Equivalently, in the relative Bergman
description, the weight is the logarithm of the generator's squared
holomorphic coefficient minus $\log g_m$ pulled back from the base
\cite[Theorem~0.1]{BerndtssonPaun2008}.
In the invariant case, use smooth direct-image positivity
\cite[Theorem~1.2]{Berndtsson2009}.  Averaging over the compact torus gives
a holomorphic projection onto the invariant line.  Distinct weight
spaces are orthogonal, so its subspace metric equals the quotient
metric and is semipositive.

There is a fixed dense open set where $f$ is a submersion and each
fibre has a patch avoiding the zeros of $s_{J+p^*D}$ and $\rho$.
Integrating over such a patch gives
$g_m\geq Cc^m$ locally, with $C,c>0$.  Hence the psh representatives of
$-m^{-1}\log g_m$ have uniform local upper bounds, independent of
possible extra base-change exceptions for individual $m$; psh
extension fills those exceptions on the fixed open set.  Finite
fibre measure and its positive density almost everywhere imply
\[
  -\frac1m\log g_m(y)\longrightarrow
  -\log\operatorname*{ess\,sup}_{W_y}|p^*\rho|_{h_L^k}^2
\]
almost everywhere. The limit is finite there by the bounded metric
and the regularity of $\rho$. Local compactness for psh functions
now gives the required representative. Indeed, a subsequence tending
locally uniformly to $-\infty$ would contradict this finite limit.
Every $L^1_{\mathrm{loc}}$-convergent subsequence has a further
subsequence converging almost everywhere, so its psh limit agrees
almost everywhere with the displayed power-mean limit. The direct-image
argument at $m=0$ supplies $-\log g_0$.
\end{proof}


\subsection{Divisors from twisted forms}

\begin{lemma}\label{g:determinant}\label{rb:determinant}
Let $L$ have bounded semipositive weights on a smooth projective
$n$-fold $V$ with $H^0(V,\Omega_V^j)=0$ for $j>0$.  There are a line
bundle $M$ mapping nontrivially into a tensor power of $\Omega_V^1$ and
effective integral divisors
\[
                       N_i\sim M+m_iL,
                       \qquad m_i\longrightarrow+\infty.
\]
The integers $m_i$ can be chosen strictly increasing and divisible by any prescribed positive integer.
\end{lemma}

\begin{proof}
Replacing $L$ at the outset by a prescribed positive multiple gives the divisibility assertion; it is therefore enough to prove the statement without that prescription. Hodge symmetry gives $\chi(\mathcal O_V)=1$, so Serre duality gives
$\chi(K_V)=(-1)^n$.  Riemann--Roch makes $\chi(K_V+mL)$ a polynomial
in $m$ with nonzero constant term.  It is nonzero for all sufficiently
large integers.  A fixed cohomological degree $q$ therefore has
$H^q(V,K_V+mL)\ne0$ for unbounded $m>0$.  The metric on $mL$ has trivial
multiplier ideal.  Hard Lefschetz with pseudoeffective coefficients
\cite[Theorem~0.1]{DPS2001} gives nonzero sections of
$\Omega_V^{n-q}\otimes mL$ at these exponents.

Take the saturated subsheaf $\mathcal E\subset\Omega_V^{n-q}$ defined
by the span of their untwisted generic images.  Its determinant
$M=(\bigwedge^{\operatorname{rank}\mathcal E}\mathcal E)^{**}$ is a line
bundle and maps to a cotangent tensor power.  This map is first defined
outside codimension two and extends because the ambient tensor bundle
is locally free.  Every nonzero vector from the family can be completed
to a generic basis using vectors of the same family.  Wedging that
basis gives a nonzero section of $M$ twisted by the sum of its positive
exponents.  These sums are unbounded because one basis vector can
have an arbitrarily large exponent.  Saturation ensures that the
maps factor through $\mathcal E$, and extension across codimension two
gives global determinant sections.  Their zero divisors are the
required $N_i$, after taking a strictly increasing subsequence.  This
is the determinant method of \cite[Lemma~4.1]{LP2018}, used for
anticanonical nonvanishing in \cite[Lemma~5.1]{LMPTX2023}.
\end{proof}

\section{Contracting the base while preserving adjoint rank}
\label{g:cf:contraction}

A relative section is useful only if it survives the next change of base.
The theorem below makes that requirement explicit. The section records
existence on a complete generic fiber; the one-dimensional adjoint space
controls the ratios that arise when a larger fiber is formed by contraction.
The proof constructs two different metrics on the base. The first is
bounded and controls the relative section. The second comes from the
pushforward volume and controls cotangent tensors. Their common numerical
zero directions determine the contraction.

\begin{theorem}[Rank-preserving contraction]\label{g:cf:contraction-theorem}
Let $X$ be a smooth connected projective complex variety with
$H^0(X,\Omega_X^q)=0$ for $q>0$, and suppose that $L=-K_X$ has a smooth
semipositive Hermitian metric. Let an algebraic torus $T$ act on $X$,
with its natural action on $L$. Let $f:X\dashrightarrow Z$ be a dominant
$T$-invariant rational map to a smooth connected projective variety,
where $\mathbb C(Z)$ is relatively algebraically closed in $\mathbb C(X)$.
On a smooth projective resolution $\pi:Y\to X$, $f_Y:Y\to Z$, put
$L_Y=\pi^*L$, $E_Y=K_{Y/X}$, and
$\eta=\operatorname{Spec}\mathbb C(Z)$. Suppose that:
\begin{enumerate}
\item $H^0(Y_\eta,pL_Y|_{Y_\eta})^T\ne0$ for some integer $p>0$;
\item there is an integer $a>0$ such that
\[
\dim_{\mathbb C(Z)}H^0(Y_\eta,(E_Y+mL_Y)|_{Y_\eta})^T=1
\quad\text{for every positive multiple }m\text{ of }a.
\]
\end{enumerate}
Then either $H^0(X,rL)^T\ne0$ for some $r>0$, or there is a dominant
rational map $h:Z\dashrightarrow W$ to a smooth connected projective
variety with $\dim W<\dim Z$ such that $h\circ f$ satisfies the same two
conditions, possibly with different positive integers. The field
$\mathbb C(W)$ is relatively algebraically closed in $\mathbb C(X)$.
\end{theorem}

We prove the theorem in four steps. Vertical orders normalize the given
section and identify the base primes at which poles disappear on $X$.
Fiber integration then supplies the bounded metric and a nef signed
log-anticanonical class. A supporting movable class produces effective
divisors in unbounded twists. Finally, interpolation of their horizontal
orders gives the new section, and the same field of section ratios restores
adjoint rank one.

\subsection{The two generic-fiber conditions}
\label{g:cf:generic-conditions}

Average the weights of the given metric over a maximal compact subgroup of
$T$.  The resulting smooth semipositive metric $h$ is invariant under that
compact subgroup.  Via the identity section of $K_X+L$, it defines a smooth
positive volume form $d\mu_X$.

Fix the rational map and resolution in Theorem~\ref{g:cf:contraction-theorem},
writing $f:Y\to Z$ for the resolved map. We call its two hypotheses the
\emph{section condition} and the \emph{rank condition}. The action on the
base and its bundles is trivial; $E_Y=K_{Y/X}$ has its natural invariant
Jacobian section.

Both conditions are unchanged by resolving further or by replacing the
base birationally.  Indeed, if $\nu:Y'\to Y$ is a further smooth
modification, then
\[
 E_{Y'}=\nu^*E_Y+K_{Y'/Y},
 \qquad \nu_*\mathcal O_{Y'}(K_{Y'/Y})=\mathcal O_Y.
\]
The comparison of the adjoint spaces uses multiplication by the relative
Jacobian section.  The same comparison on a common resolution defines the
invariant rational sections even on a model on which the torus action is
not regular.  The degree-zero rank is a consequence of both hypotheses: multiplication by a
suitable power of the section in the first condition gives an injection
into a positive-degree invariant space of dimension one (choose a common
multiple of its degree and the integer in the rank condition), while the
Jacobian section gives a nonzero element at degree zero.

If $\dim Z=0$, the section condition gives a global section on $Y$, and
then on $X$.  We may therefore assume $d=\dim Z>0$.  Every smooth base
encountered below has no positive-degree holomorphic forms: pullback
under a dominant map is injective, and holomorphic forms are invariant
under smooth projective birational modification.  Consequently
\begin{equation}\label{g:cf:euler}
 \chi(Z,\mathcal O_Z)=1.
\end{equation}
Our task is to obtain a global invariant section or construct another
map satisfying both conditions with a smaller base.

\subsection{Equidimensional preparation and the exceptional pole divisor}
\label{g:cf:preparation}

Apply weak toroidalization to a resolution of the rational map,
prescribing the complement of an isomorphism open over $X$ as part of the
source boundary \cite[Theorem~1.1]{ADK2013}.  Then apply the equidimensional
modification by projective subdivisions
\cite[Proposition~4.4, arXiv:alg-geom/9707012v1]{AK2000}.  After renaming the base, we have
\[
 \pi_0:Y_0\longrightarrow X,
 \qquad f_0:Y_0\longrightarrow Z,
\]
where $Y_0$ is normal and toroidal, $Z$ is smooth projective, $\pi_0$ is
projective birational, and $f_0$ is projective, surjective, equidimensional,
and toroidal.  The reduced source boundary $\Sigma$ contains every
$\pi_0$-exceptional divisor.  The base boundary has simple normal
crossings.  This step is birational: it does not assert that the source
is smooth or that the fibers are reduced.  We may subsequently use
smooth toroidal resolutions by subdivisions \cite{KKMSD}; these remain
toroidal over $Z$ but need not remain equidimensional.

Let $s$ be a rational section supplied by the section condition.  On
$Y_0$, its horizontal divisor is effective.  Write
\begin{equation}\label{g:cf:normalization}
 \frac1p\operatorname{div}(s)=H+f_0^*V,
\end{equation}
where the coefficient of $V$ at a prime $P\subset Z$ is the minimum of
\[
 \frac{\operatorname{coeff}_D(\operatorname{div}(s)/p)}
      {\operatorname{mult}_D(f_0^*P)}
 \quad(D\text{ a component of }f_0^*P).
\]
Equidimensionality ensures that every vertical prime maps onto a prime
of $Z$.  Thus $V$ is a rational divisor and $H$ is effective and
$\mathbb Q$-Cartier.  Choose $r>0$ sufficiently divisible.  There is an
invariant rational section
\[
 \sigma\in H^0\bigl(Y_0,r\pi_0^*L-f_0^*(rV)\bigr),
 \qquad \operatorname{div}(\sigma)=rH.
\]
Its pullback to every dominating smooth model is regular. We take $r$
also divisible by the integer in the rank condition.  For every
prime of the base there is at least one component above it with
coefficient zero in $H$.

Let $G$ be the reduced sum of the primes $P\subset Z$ for which every
component of $f_0^*P$ having coefficient zero in $H$ is exceptional over
$X$.  This is a finite sum, since each such prime is the image of an
exceptional divisor of $\pi_0$.  Every nonexceptional component above
$G$ has strictly positive coefficient in $H$.  Hence, for every
sufficiently small rational $\delta>0$,
\begin{equation}\label{g:cf:exceptional-negative}
 (H-\delta f_0^*G)^-
 \text{ is exceptional over }X.
\end{equation}
If $V+\delta G$ is rationally effective for such a $\delta$, multiply its
section by a suitable power of $\sigma$ and divide by the corresponding
canonical section of $\delta G$.  Equation
\eqref{g:cf:exceptional-negative} shows that the resulting invariant
rational section of a positive multiple of $L$ has no pole along any
prime of $X$.  Smoothness of $X$ then extends it globally.  Thus it
suffices either to prove this effectivity or to contract the base.

\subsection{The bounded metric and the pushforward density}
\label{g:cf:metrics}

Write
\begin{equation}\label{g:cf:rho}
 \rho(z)=\frac{f_*(\pi^*d\mu_X)}{d\lambda(z)}
\end{equation}
on the smooth locus, using any smooth resolution of the prepared
toroidal model. The next estimate uses only the smoothness and
positivity of $d\mu_X$; its curvature and the adjoint ranks enter
afterwards, when the density is made into a metric.

Write the base boundary as $\sum_iP_i$.  For a component $D$
above $P_i$, let
\[
 b_{iD}=\operatorname{mult}_D(f_0^*P_i),\qquad
 e_D=\operatorname{coeff}_D(K_{Y_0}-\pi_0^*K_X),
\]
where canonical divisors are chosen compatibly using a rational top
form on $X$.  At the regular generic point of $D$, $e_D$ is its
Jacobian order; it is nonnegative, and zero when $D$ is
nonexceptional.  Define
\begin{equation}\label{g:cf:thresholds}
 c_i=\min_{D\mapsto P_i}\frac{1+e_D}{b_{iD}}>0,
 \qquad B=\sum_i(1-c_i)P_i=B^+-B^-.
\end{equation}
If $c_i>1$, every component above $P_i$ is exceptional, so $P_i$ is
a component of $G$.  Thus $B^-$ is supported on the primes all of whose components
are exceptional over $X$, and in particular on $G$.  The coefficients of $B^+$ are strictly less than one.


\begin{lemma}[Toroidal density bounds]\label{g:cf:density-bounds}
For the prepared toroidal diagram and smooth positive volume above,
near an arbitrary base point,
with $y_i=0$ the local boundary branches, for every $\eta>0$,
\begin{equation}\label{g:cf:upper-density}
 \rho(y)\leq C_\eta
       \prod_i|y_i|^{-2(1-c_i+\eta)}
\end{equation}
on the general smooth locus.  Near a general point of $P_i$,
\begin{equation}\label{g:cf:lower-density}
 \rho(y)\geq C|y_i|^{-2(1-c_i)},\qquad C>0.
\end{equation}
At a general point of a prime outside the base boundary, $\rho$ also
has a positive constant lower bound.

\end{lemma}

\begin{proof}
The divisor
\begin{equation}\label{g:cf:effective-log-jacobian}
 K_{Y_0}+\Sigma-\pi_0^*K_X-\sum_i c_i f_0^*P_i
\end{equation}
is effective and $\mathbb Q$-Cartier.  Indeed, no prime contributes to
the pullbacks of two distinct $P_i$, by equidimensionality.  At a
component above $P_i$, its coefficient is $1+e_D-c_i b_{iD}\geq0$.
At the other boundary components it is $1+e_D$, and away from the
boundary it is $e_D$.  The divisor $K_{Y_0}+\Sigma$ is Cartier and
pulls back log crepantly under toroidal subdivisions: the local top
logarithmic form of a toric model trivializes it and pulls back to the
top logarithmic form on the subdivision.  Consequently, on a smooth
toroidal resolution, if $x_j=0$ are the local source boundary branches,
then
\begin{equation}\label{g:cf:orders}
 a_j:=1+e_j>0,\qquad
 a_j\geq\sum_i c_i b_{ij},
\end{equation}
where $e_j$ is the Jacobian order over $X$ and $b_{ij}$ is the order
of $y_i$ along $x_j=0$.  These inequalities include every new ray on every toroidal subdivision,
in particular the exceptional rays over intersections of boundary
components. They follow from pullback of an effective log-Cartier divisor,
so this is more than an estimate at the original boundary primes.

For use in analytic integration we make explicit why formal toroidal
charts can be replaced by analytic monomial charts.  In a formal chart
write
\begin{equation}\label{g:cf:formal-units}
 f^*y_i=\prod_jx_j^{A_{ij}},\qquad
 x_j=X_jU_j,\quad y_i=Y_iV_i.
\end{equation}
In \eqref{g:cf:formal-units}, the indices include all source and
target toric coordinates, including unit coordinates; thus
$A$ is a $d\times\dim X$ integer matrix. Here $X_j$ and $Y_i$ are
analytic boundary equations, or the constant
$1$ for unit coordinates; $U_j,V_i$ are formal units.  The coordinates
minus their residues are regular parameters.  Dominance implies
that $A$ has full row rank, since the logarithmic monomial map has
generic rank equal to the target dimension.  Factoring the analytic
boundary equations gives analytic units $R_i$ with
\[
 f^*Y_i=\Bigl(\prod_jX_j^{A_{ij}}\Bigr)R_i.
\]
Choose analytic units $V_i'$ having the same first jets as $V_i$.
We seek analytic units $U_j'$ satisfying
\begin{equation}\label{g:cf:unit-equations}
 \prod_j(U_j')^{A_{ij}}=R_i f^*V_i'.
\end{equation}
Use the constants of the formal solution, divide by them, and take
logarithms near $1$.  The resulting equations are linear equations
with matrix $A$ and holomorphic right hand sides.  A linear right
inverse gives a holomorphic solution.  Its first jet differs from
the prescribed formal first jet by a vector in $\ker A$; a
kernel-valued affine correction removes that difference without
changing the equations.  Exponentiating gives the desired $U_j'$.
Thus $X_jU_j'$ and $Y_iV_i'$ have the original first jets and remain
coordinate systems, while the monomial identities hold exactly
analytically.  Negative exponents can occur only on unit
coordinates and cause no difficulty.

We can now estimate by coarea in logarithmic coordinates.  Cut
arguments into bounded sectors and work on relatively compact
pieces of these charts, away from the boundary itself.  Relative
to Lebesgue measure in the logarithmic variables, the pullback
of $d\mu_X$ is bounded above by a constant times
\[
 \prod_j|x_j|^{2a_j}
\]
over the vanishing coordinates; all unit moduli are bounded above
and away from zero.  The factor $|x_j|^2$ in this formula is the
Jacobian of the logarithmic change of variables.  The analogous
change on the base introduces $\prod_i|y_i|^2$, up to bounded
positive factors from its unit coordinates.  The logarithmic
map is the fixed full-rank complex linear map $A$, up to boundedly
many argument translates.  Its pushforward density is therefore
an integral over affine real planes with fixed coarea factors.

In the following products, $i$ ranges only over vanishing target
coordinates and $j$ only over vanishing source coordinates; $b_{ij}$
is the corresponding boundary-order block of $A$.
To estimate $\rho\prod_i|y_i|^{2(1-c_i+\eta)}$, divide the
logarithmic pushforward density by
$\prod_i|y_i|^{2(c_i-\eta)}$.  On an affine fiber, this factor is
constant and can be brought into the integrand.  The resulting
exponent on a vanishing source coordinate is
\begin{equation}\label{g:cf:residual-exponents}
 2\Bigl(a_j-\sum_i(c_i-\eta)b_{ij}\Bigr)>0.
\end{equation}
If some $b_{ij}>0$, strict positivity follows from
\eqref{g:cf:orders} and $\eta>0$; otherwise it follows from $a_j>0$.
The integrand thus decays exponentially as the real logarithms
of vanishing coordinates tend to $-\infty$.  Their upper bounds,
and the bounds on all remaining logarithmic coordinates, allow
us to dominate it by $C\exp(-\varepsilon\|u\|)$ on the whole
chart piece in real logarithmic variables.  Integrals of this
bound over affine planes of a fixed direction are uniformly
bounded independently of the translation.  For example, write
such a plane as $v+\ker A$ with $v$ orthogonal to $\ker A$ and
use $\|v+w\|\geq\|w\|$.  Properness supplies finitely many
chart pieces over a sufficiently small base neighborhood.  The
coordinate equations of each boundary prime differ only by
bounded units on overlaps.  Summing the bounds proves
\eqref{g:cf:upper-density} almost everywhere, and hence on the
general locus where the density is continuous.

For \eqref{g:cf:lower-density}, take a general point of a component
which realizes the minimum in \eqref{g:cf:thresholds}.  At that point take coordinates
\begin{equation}\label{g:cf:power-chart}
 y_1=x_1^{b_{iD}},\qquad y_j=x_j\quad(2\leq j\leq d).
\end{equation}  The pulled-back volume is
comparable there to $|x_1|^{2e_D}$ times ordinary coordinate
volume.  The change $y_1=x_1^{b_{iD}}$, followed by integration
over a fixed patch in the remaining fiber variables, contributes
\[
 |y_1|^{2((1+e_D)/b_{iD}-1)}
   =|y_1|^{-2(1-c_i)}.
\]
This is a lower bound for the full pushforward density.

At a general point of an excluded prime outside the base boundary,
the density also has a positive constant lower bound.  In fact
there is a nonexceptional component over that prime.  An
exceptional component would lie in the source boundary, whereas
a boundary component meeting the inverse image of the base
toroidal open cannot be vertical there: all target coordinates
are units, so the full-rank monomial matrix uses only source
unit coordinates and the map restricted to the component is
dominant.  At a nonexceptional component the Jacobian order is
zero. If $b$ is its pullback multiplicity over the base prime,
the same power-coordinate calculation gives a lower bound,
since its exponent is $2(1/b-1)\leq0$.
\end{proof}

\begin{lemma}\label{g:cf:metric-conclusions}
In the preceding setup, $V$ has a semipositive metric with locally bounded
weights. For the effective rational divisors $B^+$ and $B^-$ defined
in \eqref{g:cf:thresholds}, there is a nef rational divisor $M$ such that
\begin{equation}\label{g:cf:log-relation}
 K_Z+B^++M\sim_{\mathbb Q}B^-.
\end{equation}
The support of $B^+$ has simple normal crossings and all its coefficients
are less than one.  The divisor $B^-$ is supported on those primes of $Z$ all of whose
components on $Y_0$ are exceptional over $X$; in particular it is
supported on $G$.
\end{lemma}

\begin{proof}
We first match the data with the smooth invariant version of
Lemmas~\ref{g:moments} and~\ref{g:transfer}. Their original variety is
$X$, their allowed-pole divisor is zero, and both input bundles are
$A=L=-K_X$ with the compact-invariant smooth metric $h$. The finite
measure is $d\mu_X$. Take exponent $r$ and the rational section
$\widetilde s$ of $rL$ whose pullback has divisor $rH+f_0^*(rV)$.
It is the power of the original generic section used to construct
$\sigma$. It is invariant and has no horizontal poles. On an equivariant
resolution the invariant adjoint ranks equal one for every multiple of
$r$, including zero by the argument above. All hypotheses of
Lemma~\ref{g:moments} are therefore satisfied.

The zeroth density is $\rho$ from \eqref{g:cf:rho}; hence
$-\log\rho$ is psh on the smooth base-change locus. The normalized
base divisor in Lemma~\ref{g:transfer} is $rV$, by the full-component
minimum in \eqref{g:cf:normalization}. Its bounded metric, divided by
$r$, gives the asserted metric on $V$. More concretely, for a local frame
$v$ of $rV$, it is determined by
\[
 |v|^2=\left(\max_{f_0^{-1}(z)}
      |\sigma f_0^*v|_{(\pi_0^*h)^r}^{2/r}\right)^r.
\]
The essential supremum in the transfer lemma is this maximum: on a
smooth connected general fiber the smooth measure has full support.
The invariance step is the orthogonal quotient in the smooth
Berndtsson theorem, as proved in Lemma~\ref{g:moments}; no singular
invariant-subbundle positivity is invoked.

We now use Lemma~\ref{g:cf:density-bounds} to construct the signed
log-anticanonical metric. Put $M=-K_Z-B$.  On the general open its local metric weight is
\begin{equation}\label{g:cf:corrected-weight}
 \varphi_M=-\log\rho-\sum_i(1-c_i)\log|y_i|^2.
\end{equation}
This has the transition law of $M$: a frame of $-K_Z$ is
multiplied by the boundary factor $\prod_i y_i^{1-c_i}$,
after clearing denominators.  It is plurisubharmonic away
from the boundary.  The lower density bounds give the upper
bound needed for extension across a dense open of every
excluded divisor.  Hartogs extension treats the remaining
codimension-two set.  The upper density estimate implies
\begin{equation}\label{g:cf:arbitrarily-small-poles}
 \varphi_M\geq\eta\sum_i\log|y_i|^2-O_\eta(1)
 \qquad\text{for every }\eta>0.
\end{equation}
In particular this metric has zero Lelong numbers: the lower estimate
holds with arbitrarily small logarithmic coefficients. For every fixed
divisible $k>0$, choose $\eta<1/k$.
Normal-crossing integrability shows that
$\mathcal I(k\varphi_M)=\mathcal O_Z$.  This assertion is
about every such power, not just the first multiplier ideal.

Choose a very ample divisor $A$ on $Z$.  Nadel vanishing
\cite{Nadel}, applied after adding the smooth positive metric
on $(d+1-i)A$, gives
\[
 H^i\bigl(Z,K_Z+kM+(d+1-i)A\bigr)=0
 \qquad(1\leq i\leq d).
\]
Castelnuovo--Mumford regularity makes
$K_Z+kM+(d+1)A$ globally generated.  Its class divided by $k$
is nef, and the limit as $k$ tends to infinity proves that
$M$ is nef.  The identical argument for the bounded metric
on $V$, whose every power has trivial multiplier ideal,
also shows that $V$ is nef.  Equation
\eqref{g:cf:log-relation} follows from the definition of $M$.
\end{proof}

\subsection{Cotangent tensors on a supporting movable class}
\label{g:cf:slopes}

If $V+G$ is big, then for $0<\delta\leq1$,
\[
 V+\delta G=\delta(V+G)+(1-\delta)V
\]
is big, because $V$ is pseudoeffective.  Choosing a sufficiently
small rational $\delta$ and using
\eqref{g:cf:exceptional-negative} finishes the proof.  We may
therefore assume that $V+G$ is not big.  Duality between the
pseudoeffective divisor cone and the closed movable curve cone
\cite{BDPP2013} gives a nonzero movable curve class $\beta$ with
\[
 (V+G)\cdot\beta=0.
\]
Since $V$ and $G$ are pseudoeffective and $B^-$ is supported
on $G$, we obtain
\begin{equation}\label{g:cf:null-degrees}
 V\cdot\beta=G\cdot\beta=B^-\cdot\beta=0.
\end{equation}

\begin{lemma}\label{g:cf:cotangent-slope}
Every line subsheaf $N\subseteq(\Omega_Z^1)^{\otimes s}$, for
$s\geq0$, satisfies $c_1(N)\cdot\beta\leq0$.
\end{lemma}

\begin{proof}
For $s=0$, the inclusion into $\mathcal O_Z$ gives the
assertion directly.  Let $s>0$ and let $t>0$ be rational.
Since $M$ is nef, $M+tA$ is ample.  Choose a sufficiently
divisible very ample multiple and a general member; divide
that member by the multiple and add it to $B^+$ to obtain a
rational simple normal crossings boundary $\Delta_t$, all
of whose coefficients lie in $[0,1)$.  Then
\[
 K_Z+\Delta_t\sim_{\mathbb Q}B^-+tA
\]
is pseudoeffective.

Apply the orbifold cotangent theorem
\cite[Theorem~1.3 and Section~5]{CampanaPaun2019} to this smooth
projective log-canonical pair.  On a finite adapted cover
$\tau:\widetilde Z\to Z$, the adapted orbifold cotangent
bundle $\mathcal E_t$ has determinant class
$\tau^*(K_Z+\Delta_t)$.  Every torsion-free quotient of each
$\mathcal E_t^{\otimes s}$ has nonnegative degree against
$\tau^*\beta$.  This holds for classes on the boundary of
the movable cone as well.

There is an inclusion $\tau^*\Omega_Z^1\subseteq\mathcal E_t$.
In codimension one, at a branch $z=w^b$ corresponding to
boundary coefficient $\gamma$, the orbifold generator is
$w^{-b\gamma}dz$; hence the ordinary generator $dz$ belongs
to the adapted bundle.  Additional branch divisors have
$\gamma=0$.  Let $\mathcal N$ be the saturation of the
pulled-back line $\tau^*N$ in $\mathcal E_t^{\otimes s}$.
Saturation adds an effective divisorial class, so
\[
 \deg(\tau)\,c_1(N)\cdot\beta
   \leq c_1(\mathcal N)\cdot\tau^*\beta.
\]
The quotient by $\mathcal N$ is torsion-free.  Its determinant
has nonnegative degree, and the tensor determinant identity
for a bundle of rank $d$ is
\[
 c_1(\mathcal E_t^{\otimes s})
      =s d^{s-1}c_1(\mathcal E_t).
\]
If the quotient has rank zero, it is zero and the same
inequality holds with determinant degree zero.  Dividing
by the cover degree and using \eqref{g:cf:null-degrees}
gives
\begin{equation}\label{g:cf:tensor-bound}
 c_1(N)\cdot\beta
   \leq s d^{s-1}(K_Z+\Delta_t)\cdot\beta
   =s d^{s-1}tA\cdot\beta.
\end{equation}
Letting $t$ decrease to zero proves the assertion.
\end{proof}

\begin{lemma}\label{g:cf:effective-sequence}
There are a fixed line bundle $N$ and effective divisors $D_j$
such that
\begin{equation}\label{g:cf:divisor-sequence}
 D_j\sim m_jV+N,\qquad
 m_j\longrightarrow+\infty,\qquad
 D_j\cdot\beta=0,
\end{equation}
where the positive integers $m_j$ may be chosen strictly increasing
and sufficiently divisible.
\end{lemma}

\begin{proof}
Choose $b>0$ with $bV$ Cartier. Lemma~\ref{g:cf:metric-conclusions}
gives $bV$ bounded semipositive weights, and the base has no
positive-degree forms by the argument preceding \eqref{g:cf:euler}.
Apply Lemma~\ref{g:determinant} to $bV$, prescribing any further
divisibility required for the construction. It gives a cotangent line
$N$ and effective divisors $D_j\sim N+n_j(bV)$ with $n_j\to\infty$.
Set $m_j=bn_j$. Lemma~\ref{g:cf:cotangent-slope} gives
$N\cdot\beta\leq0$, while $V\cdot\beta=0$. Effectivity gives the
opposite inequality for $D_j\cdot\beta$, so every degree is zero.
\end{proof}

\subsection{Maximal zero-degree systems and horizontal interpolation}
\label{g:cf:zero-systems}

Choose a linear system of divisor class $D$ with
$D\cdot\beta=0$ whose rational map has maximal possible
image dimension among all such systems.  The constant
map is permitted.  This maximal dimension is smaller
than $d$: a generically finite system makes its divisor
big, whereas a nonzero supporting functional cannot
vanish on a big class.  Resolve this rational map and
take its Stein factorization, followed by a smooth
projective resolution of the Stein base.  We obtain
\[
 h:\widehat Z\longrightarrow Z,\qquad
 \widehat g:\widehat Z\longrightarrow W,
 \qquad \dim W<d,
\]
where the general fibers of $\widehat g$ are connected.

\begin{lemma}\label{g:cf:maximal-field}
Every ratio of two nonzero sections in any one linear
system whose divisor class has degree zero against
$\beta$ belongs to $\mathbb C(W)\subseteq\mathbb C(Z)$.
\end{lemma}

\begin{proof}
Multiply the sections of that system by those defining
the chosen maximal system.  Their products form a
system of degree zero, and its map is the product of
the two maps followed by the Segre embedding.  By
maximality it has the same image dimension as the
chosen map.  The added section ratios are therefore
algebraic over the field of that image.  The Stein
factorization identifies $\mathbb C(W)$ with the
relative algebraic closure of the image field in
$\mathbb C(Z)$.  This proves the assertion.
\end{proof}

The interpolation below is a relative form of the divisor argument in
\cite[proof of Theorem~3.1, Step~3]{LP2018}. There, zero Iitaka
dimension makes linearly equivalent effective divisors equal. Here
the maximal-field property gives equality of their horizontal parts.
The subsequent rank-restoration argument still has to control their
vertical poles.

\begin{lemma}\label{g:cf:horizontal-slope}
The rational divisor $h^*V$ is rationally effective
on the complete generic fiber of $\widehat g$.
\end{lemma}

\begin{proof}
Use the sequence in \eqref{g:cf:divisor-sequence}.
For $j<k<l$, linear equivalence gives
\begin{equation}\label{g:cf:interpolation}
 (m_l-m_j)D_k
   \sim (m_l-m_k)D_j+(m_k-m_j)D_l.
\end{equation}
Both sides are effective divisors in the same
zero-degree system.  By Lemma~\ref{g:cf:maximal-field},
their ratio is a rational function from $W$.
Its pullback has no horizontal divisor over $W$.
Thus \eqref{g:cf:interpolation}, after pullback by
$h$, is an equality of actual horizontal divisors.

Fix the first two indices.  For every horizontal
prime $Q$ on $\widehat Z$, the coefficients of
$h^*D_j$ have the form $a_Q+m_jb_Q$.  The preceding
identity proves this simultaneously for every
index.  Since all $D_j$ are effective and
$m_j\to\infty$, every $b_Q$ is nonnegative.
These slopes have finite support: they are the
coefficients of the horizontal part of
\[
 \frac{h^*(D_2-D_1)}{m_2-m_1}
     \sim_{\mathbb Q}h^*V.
\]
This representative of $h^*V$ has effective
horizontal part.  Its vertical components disappear
on the generic fiber, yielding the stated linear
effectivity there.  The argument uses actual
divisors and rational functions, rather than only
numerical equivalence.
\end{proof}

\subsection{Preserving the section and rank conditions}
\label{g:cf:rank-preservation}

Compose $X\dashrightarrow Z$ with $Z\dashrightarrow W$.
The composition is invariant and has connected general
fibers: the inclusions of function fields remain
relatively algebraically closed.  Choose a smooth
projective resolution $\pi':Y'\to X$ dominating $Y_0$
and the map to $\widehat Z$, and write
\[
 f':Y'\longrightarrow\widehat Z,\qquad
 u=\widehat g\circ f'.
\]
Multiplication of the pullback of $\sigma$ by the
generic section of $h^*V$ from
Lemma~\ref{g:cf:horizontal-slope} gives an invariant
section of a positive multiple of $L_{Y'}$ on the
generic fiber of $u$.  Thus the section condition
is preserved.

To prove the rank condition, fix a sufficiently
divisible $m>0$ and take any two nonzero invariant
sections of $E_{Y'}+mL_{Y'}$ over the generic fiber
of $u$.  They can be extended simultaneously to
global sections of
\begin{equation}\label{g:cf:spreading-twist}
 E_{Y'}+mL_{Y'}+ku^*A_W
\end{equation}
for a sufficiently large integer $k$, where $A_W$
is very ample.  For completeness, every pole
divisor has image in a proper closed subset of
$W$.  A high multiple of a very ample bundle
has a section in any prescribed power of the
ideal of this finite union of images.  Pullback
then vanishes to at least the required orders
on all the finitely many pole divisors.  The
same section clears the poles of both rational
sections.  This also handles images of
codimension greater than one.  It is a base
section, so it preserves invariance and does
not change their ratio.

Divide the two sections of
\eqref{g:cf:spreading-twist} by the Jacobian
section and $\sigma^{m/r}$.  The old rank
condition says that the results come from
rational sections on $\widehat Z$ of
\[
 mh^*V+k\widehat g^*A_W.
\]
Push their divisors down to $Z$.  A pole outside
$G$ is impossible: for a prime $P\not\subseteq G$,
choose a component above $P$ on $Y_0$ with
coefficient zero in $H$ and nonexceptional over
$X$.  Its strict transform on $Y'$ has
coefficient zero in both the pulled-back $H$
and the Jacobian divisor $E_{Y'}$.  Since $h$
is an isomorphism at the generic point of $P$,
a pole downstairs would remain a pole on this
component, contradicting regularity of the
original global section.

After bounding the poles of both sections
along $G$ by a common integer $\ell$, their
ratio is a ratio of sections in the single
linear system
\begin{equation}\label{g:cf:rank-system}
 mV+k h_*\widehat g^*A_W+\ell G.
\end{equation}
We verify that its degree against $\beta$ is
zero.  Only the middle term needs proof.
If $W$ is a point it is zero.  Otherwise,
let $A_I$ be the pullback to $W$ of the
hyperplane class of the image of the
defining maximal system.  The map from
$W$ to that image is generically finite,
so $A_I$ is big.  On a sufficiently high
resolution of the system,
\[
 h^*D\sim\widehat g^*A_I+F,
 \qquad F\geq0.
\]
Some sufficiently large multiple of $A_I$
minus $A_W$ is rationally effective, by
bigness.  Pulling back and pushing forward
shows that a positive multiple of $D$
minus $h_*\widehat g^*A_W$ is rationally
effective.  The latter divisor is itself
effective after choosing an effective
representative of $A_W$.  Since
$D\cdot\beta=0$ and movable classes are
nonnegative on effective divisors, both
inequalities force
\[
 (h_*\widehat g^*A_W)\cdot\beta=0.
\]
Together with \eqref{g:cf:null-degrees},
this proves the asserted zero degree of
\eqref{g:cf:rank-system}.

Lemma~\ref{g:cf:maximal-field} now puts
the ratio of the two original sections
in $\mathbb C(W)$.  Thus the new
invariant generic adjoint space has
dimension at most one for every positive multiple of a common integer.  The new section
condition and the Jacobian supply a
nonzero element in every sufficiently
divisible degree, giving dimension
exactly one.  Both conditions have
therefore been preserved while the
base dimension has strictly decreased.

This proves Theorem~\ref{g:cf:contraction-theorem}. The construction is
conditional on the two generic-fiber hypotheses, and their preservation is
part of its conclusion. In particular the argument can be iterated without
introducing a strict-curvature assumption on the base.

\section{Boundary corrections and integrable metrics}
\label{g:bt}

Integration along a fibration can lose regularity along its discriminant.
The correction needed to extend the resulting metric also determines which
sections can be returned to the original variety.  We keep these two roles
together.  Three outputs will be useful: a metric integrable after
multiplication by a boundary section; an unweighted integrable metric; and
a nef anti-log-canonical class for a signed boundary.  Their slope
consequences are proved under the hypotheses actually available.

A metric weight is normalized by $|e|^2=\exp(-\psi)$ in a holomorphic
frame.  Integrability at exponent one therefore means integrability of
$\exp(-\psi)$.  For a rational line bundle we clear denominators and
divide the resulting weight by the denominator.  A signed rational
divisor $B$ on a smooth variety defines a \emph{sub-klt pair} if every
coefficient of its crepant boundary on a log resolution is less than one.
Negative coefficients are allowed; a sub-klt divisor need not be an
effective klt boundary.

\subsection{Approximation without a strict curvature hypothesis}
\label{g:bt:approximation}

The approximation below uses a fixed ample auxiliary line to produce
sections.  The metric being approximated is only semipositive.  We give
both the ideal-theoretic construction and the Gaussian construction,
since the latter directly preserves the weighted integral.

\begin{lemma}[Weighted sub-klt approximation]
\label{g:bt:weighted-approximation}
Let $Y$ be smooth projective, of dimension $d>0$, let $J$ be a rational
line bundle with a semipositive metric of weights $\psi$, and let $P\geq0$
be a rational divisor.  Suppose that, in local divisor frames,
\[
                 |s_P|^2e^{-\psi}\in L^1_{\mathrm{loc}}.
\]
There is a fixed ample Cartier divisor $A$ such that, for every sufficiently
large divisible integer $m$, there is an effective rational divisor
$\Delta_m\sim_{\mathbb Q}J+A/m$ for which $(Y,\Delta_m-P)$ is sub-klt.
When $P=0$, the pair $(Y,\Delta_m)$ is klt.
\end{lemma}

\begin{proof}
Choose a very ample $H$ and an ample $A$ such that $A-K_Y-iH$ is ample for
$1\leq i\leq d$.  Nadel vanishing and Castelnuovo--Mumford regularity give
global generation of
\[
       \mathcal O_Y(mJ+A)\otimes\mathcal I(m\psi).
\]
Indeed its twist by $-iH$ is an adjoint sheaf with coefficient metric of
strictly positive curvature, supplied by $A-K_Y-iH$, and multiplier ideal
$\mathcal I(m\psi)$.  The ideal is coherent and nonzero; local point
extension where $\psi$ is finite gives a nonzero integrable germ
\cite{Nadel,OhsawaTakegoshi,raDemailly}.

Let $u:Y'\to Y$ resolve this ideal and $P$, and write
$\mathcal I(m\psi)\mathcal O_{Y'}=\mathcal O_{Y'}(-F_m)$.
For a prime $E$ on $Y'$ put
\[
 b_E=\operatorname{ord}_E F_m,\qquad
 p_E=\operatorname{ord}_E u^*P,\qquad
 a_E=\operatorname{ord}_E K_{Y'/Y}.
\]
We claim that
\begin{equation}\label{g:bt:strict-approximation}
                         b_E/m-p_E<1+a_E.
\end{equation}
The local point-extension estimate for $m\psi$ supplies holomorphic
functions $v_x$ with $|v_x(x)|^2=e^{m\psi(x)}$ and uniformly bounded
$L^2(e^{-m\psi})$ norm on fixed nested coordinate balls.  The local upper
bound for $\psi$ bounds their ordinary norms, hence their suprema on a
smaller ball.  Each $v_x$ belongs to $\mathcal I(m\psi)$.  At a general
point of $E=(z_1=0)$, Cauchy estimates after division by $z_1^{b_E}$ give
\[
               e^{m\psi\circ u}\leq C_m|z_1|^{2b_E}.
\]
The change-of-variables formula bounds the pulled-back integrable density
below by a positive constant times $|z_1|^{2(p_E+a_E-b_E/m)}$.
Its transverse integral is finite only if
\eqref{g:bt:strict-approximation} holds.

Global generation identifies the fixed divisor of the pulled-back
system with $F_m$.  Choose a general member of the free residual system
and divide the corresponding divisor downstairs by $m$.  The crepant
boundary of $\Delta_m-P$ has coefficient $b_E/m-p_E-a_E<1$ on each fixed
component, and coefficient $1/m<1$ on new moving components.  Bertini,
applied also to the strata of the resolved divisor, makes their union
simple normal crossings.  Decreasing coefficients from the positive
part preserves positive log discrepancies, proving sub-klt on every
higher model.

There is a second way to choose the member.  The global Bergman lower
estimate for the same sufficiently ample $A$ gives local expressions
$s_{i,m}$ of an orthonormal basis with
\[
                   \sum_i|s_{i,m}|^2\geq C_m e^{m\psi}.
\]
This globalization by a fixed ample twist is the point-extension
construction in \cite[Section~9, proof of Proposition~9.1]{Demailly1993}.
One obtains the estimate by extending from a point, cutting off, and
solving $\bar\partial$ globally with a logarithmic pole of order $d$ at
that point.  The fixed positivity of $A-K_Y$ pays the uniformly bounded
negative curvature of the cutoff pole; integrability forces the
correction to vanish at the point.  Singular weights are handled by
decreasing regularization and weak limits in the $L^2$ estimate
\cite[Corollary~13.9, Remark~13.10, and Theorem~14.2]{raDemailly}.
For independent standard complex Gaussians $c_i$ and $m>1$,
\[
 \mathbb E\left|\sum_i c_i s_{i,m}\right|^{-2/m}
   =C'_m\left(\sum_i|s_{i,m}|^2\right)^{-1/m},
\]
where $C'_m<\infty$ because $\mathbb E|c_1|^{-2/m}<\infty$.
Multiplication by $|s_P|^2$ and Fubini show that almost every such member
has the weighted integrability required for
\eqref{g:bt:strict-approximation}.
\end{proof}

The approximation retains the negative part $-P$.  Its immediate use is
a cotangent estimate.  We first isolate the common slope argument.

\begin{lemma}[From foliation determinants to tensor bounds]
\label{g:bt:slope-principle}
Let $Y$ be smooth projective and $P\geq0$ rational.  Suppose that
$-K_Y+P$ is pseudoeffective and that $K_{\mathcal F}-K_Y+P$ is
pseudoeffective for every algebraically integrable saturated foliation
$\mathcal F\subset T_Y$ of positive rank and corank.  Then
$hP-c_1(M)$ is pseudoeffective for every line subsheaf
$M\subseteq(\Omega_Y^1)^{\otimes h}$, $h\geq0$.
\end{lemma}

\begin{proof}
The assertion on a point is immediate. For positive dimension,
fix a nonzero movable class $\alpha$ and put $t=\alpha\cdot P\geq0$.
The total degree of $T_Y$ is at least $-t$.  If its minimum
Harder--Narasimhan slope were less than $-t$, its positive-slope part
$\mathcal F$ would be nonzero and proper: the last negative graded piece
already has degree less than $-t$, whereas the whole sheaf does not.
The quotient by the positive part has only nonpositive slopes and degree
less than $-t$.  In particular
\[
 \mu_{\alpha,\min}(\mathcal F)>0,\qquad
 2\mu_{\alpha,\min}(\mathcal F)>
       \mu_{\alpha,\max}(T_Y/\mathcal F).
\]
The bracket map vanishes by the slope inequality, and the
Campana--P\u aun algebraicity theorem makes $\mathcal F$ algebraically
integrable \cite[Proposition~2.2]{Ou2023}.  But
$c_1(T_Y/\mathcal F)=K_{\mathcal F}-K_Y$ has degree less than $-t$,
contradicting the hypothesis.  Thus
$\mu_{\alpha,\min}(T_Y)\geq-t$.
Duality and the tensor-product rule for movable slopes give
$c_1(M)\cdot\alpha\leq ht$ \cite{GKP2016,CampanaPaun2019}.
Movable-cone duality proves the assertion \cite{BDPP2013}.
For $h=0$, the inclusion in $\mathcal O_Y$ proves it directly.
\end{proof}

\begin{proposition}[Finite-volume cotangent estimate]
\label{g:bt:weighted-slope}\label{g:slope-recalled}\label{rb:cotangent}
Let $Y$ be smooth projective and $P\geq0$ rational.  If $J=-K_Y+P$
has a semipositive metric with $|s_P|^2e^{-\psi}$ locally integrable, then
$hP-c_1(M)$ is pseudoeffective for every line subsheaf
$M\subseteq(\Omega_Y^1)^{\otimes h}$.
If $\dim Y>0$ and $\alpha\ne0$ is movable, then
$\mu_{\alpha,\min}(T_Y)\geq-P\cdot\alpha$. In particular,
$P\cdot\alpha=0$ gives $\mu_{\alpha,\min}(T_Y)\geq0$ and
$c_1(M)\cdot\alpha\leq0$.
\end{proposition}

\begin{proof}
The assertion on a point is immediate. Suppose $\dim Y>0$.
Lemma~\ref{g:bt:weighted-approximation} gives effective
$\Delta_m\sim_{\mathbb Q}-K_Y+P+A/m$ with
$(Y,\Delta_m-P)$ sub-klt.  Resolve a first-integral map for an algebraic
foliation $\mathcal F$ together with that pair, by $u:Y'\to Y$.
Write its crepant boundary as $B_m^+-B_m^-$, with disjoint effective
parts.  Then
\[
 K_{Y'}+B_m^+\sim_{\mathbb Q}u^*A/m+B_m^-,
 \qquad 0\leq u_*B_m^-\leq P.
\]
The positive boundary is klt.  A general fibre of the resolved
first-integral map meets the isomorphism locus of $u$; hence $u^*A$
restricts to a big divisor there.  A divisible positive multiple of
$K_{Y'}+B_m^+$ consequently has a nonzero section on that fibre.
Druel's relative pseudoeffectivity theorem, in the form recalled by Ou,
applies to the effective boundary and its log canonical general fibre
\cite[Proposition~4.1]{Druel2017},
\cite[Proposition~3.1]{Ou2023}. It yields pseudoeffectivity of
$K_{\mathcal F'}+B_m^+$.  Push down, add $P-u_*B_m^-\geq0$, and let $m$
tend to infinity.  This gives $K_{\mathcal F}-K_Y+P$ pseudoeffective.
The metric already gives $-K_Y+P$ pseudoeffective, so
Lemma~\ref{g:bt:slope-principle} applies.  This is the signed finite-volume
adaptation of Ou's tangent-slope argument; its relative positivity input
is due to Druel.
\end{proof}

\subsection{An unweighted integral and an annihilating movable class}
\label{g:bt:unweighted}

When the integral contains no boundary factor, orbifold cotangent
semipositivity gives a different estimate.  Its hypothesis concerns
$K_Y+J$, which need not be supported on an allowed-pole divisor.

\begin{proposition}[Unweighted movable-slope estimate]
\label{g:bt:unweighted-slope}
Let $Y$ be smooth projective of dimension $d>0$, and let $J$ have a
semipositive metric with $e^{-\psi}\in L^1_{\mathrm{loc}}$.
Suppose $K_Y+J$ is pseudoeffective.  For every movable class $\alpha$
and every line subsheaf $M\subseteq(\Omega_Y^1)^{\otimes h}$, $h>0$,
\begin{equation}\label{g:bt:orbifold-bound}
 c_1(M)\cdot\alpha\leq
           h d^{h-1}(K_Y+J)\cdot\alpha.
\end{equation}
Thus $\alpha\cdot(K_Y+J)=0$ gives $c_1(M)\cdot\alpha\leq0$.
For $h=0$ the latter inequality holds without the annihilation premise.
\end{proposition}

\begin{proof}
Use Lemma~\ref{g:bt:weighted-approximation} with $P=0$ to obtain klt
$\Delta_m\sim_{\mathbb Q}J+A/m$.  On a log resolution $u:Y'\to Y$
truncate the negative coefficients of the crepant boundary.  This gives
an effective simple normal crossing boundary $\Delta'_m$, with
coefficients less than one, and
\[
 K_{Y'}+\Delta'_m=u^*(K_Y+\Delta_m)+E_m,\qquad
 E_m\geq0\quad\hbox{$u$-exceptional}.
\]
The left side is pseudoeffective.  The numerical pullback $\alpha'$,
defined by $\alpha'\cdot D=\alpha\cdot u_*D$, is movable and annihilates
exceptional divisors.

On an adapted finite cover $v:\widehat Y\to Y'$, the orbifold cotangent
sheaf $\mathcal E$ contains $v^*\Omega^1_{Y'}$ and satisfies
$c_1(\mathcal E)=v^*(K_{Y'}+\Delta'_m)$.  In a boundary chart its
generators are pullbacks of $dx/x^\delta$, where $\delta$ is the boundary
coefficient and the cover clears its denominator.  Campana--P\u aun
tensor-quotient semipositivity gives nonnegative $v^*\alpha'$-degree to
torsion-free quotients of $\mathcal E^{\otimes h}$
\cite[Theorem~1.3]{CampanaPaun2019}.
The line $v^*u^*M$ injects through ordinary cotangent pullback.
Saturating increases its determinant by an effective divisor.  The
quotient has nonnegative degree; if its rank is zero, the same inequality
follows directly from the effective saturation divisor.  Since
$c_1(\mathcal E^{\otimes h})=h d^{h-1}c_1(\mathcal E)$, cancellation of
the finite covering degree gives
\[
 c_1(M)\cdot\alpha\leq
    h d^{h-1}(K_Y+J+A/m)\cdot\alpha.
\]
Let $m$ increase to obtain \eqref{g:bt:orbifold-bound}.
\end{proof}

For $J=-K_Y+\Theta$ with $\Theta\geq0$, cone duality gives
$h d^{h-1}\Theta-c_1(M)$ pseudoeffective.
For $J=-K_Y+kN+R$, the proposition instead applies to a movable class
annihilating $N$ and $R$.  Both situations occur below.

\subsection{Weighted transfer with arbitrary exceptional poles}
\label{g:bt:exceptional}

Let $X$ be smooth projective, let $L=-K_X$ carry a smooth semipositive
metric, and let a torus $T$ act with the natural linearization.  Average
the metric over a compact real form.  For a $T$-invariant rational map to
a smooth projective $Y$, let
\[
                 X\xleftarrow{\pi}Z\xrightarrow{g}Y
\]
be a smooth equivariant resolution, with geometrically integral generic
fibre $F$ over $K=\mathbb C(Y)$.  We distinguish the following strong
rank condition: there is a nonzero invariant section $s$ of
$a\pi^*L|_F$, $a>0$, and
\begin{equation}\label{g:bt:all-exceptional-rank}
 \dim_K H^0\bigl(F,(b\pi^*L+D)|_F\bigr)^T\leq1
\end{equation}
for every integer $b\geq0$ and every effective invariant
$\pi$-exceptional integral divisor $D$.
This condition persists on higher equivariant resolutions: push down
sections and test regularity at codimension one.  Exceptional divisors
for an additional resolution remain exceptional on the generic fibre.

Flattening by modification, followed by finite normalization, provides
a normal projective model $p:W\to X$, $f:W\to Y$ with no prime divisor
mapping into codimension at least two in $Y$ \cite{RaynaudGruson1971}.
Here and below the base function field is unchanged.
Regard $s$ as a rational section on $W$.  For every base prime $Q$, set
\[
 e_E=\operatorname{ord}_E f^*Q,\qquad
 n_Q=\min_{E\mapsto Q}\frac{\operatorname{ord}_E s}{e_E},
 \qquad N=\sum_Q n_Q Q.
\]
The sums have finite support.  Let $B$ be the reduced sum of $Q$ for
which every component above $Q$ is exceptional over $X$.  Away from $B$
define $n'_Q$ by the same minimum restricted to nonexceptional components;
on $B$ put $n'_Q=n_Q$.  Set $R=\sum_Q(n'_Q-n_Q)Q\geq0$.
Taking a power of $s$, make $N,R$ integral.  Then
\begin{equation}\label{g:bt:normalization}
 D_s=\operatorname{div}_W(s)-f^*N\geq0,\qquad
       f^*\mathcal O_Y(N)\longrightarrow a p^*L
       \quad\hbox{is regular}.
\end{equation}
Normality and the divisorial image property justify regularity.

\begin{lemma}[Extension of the normalized maximum]
\label{g:bt:bounded-maximum}
In the normalization \eqref{g:bt:normalization}, suppose the negative
logarithms of the adjoint moments are psh for all sufficiently large
exponents on a smooth connected-fibre open set.
Then the maximum norm gives a semipositive metric with locally bounded
weights on $N$.
\end{lemma}

\begin{proof}
In a local frame of $N$, write $\sigma$ for the regular normalized
section.  On the smooth open set, the moments have limit
\[
 -m^{-1}\log\int_{Z_y}|\sigma|^{2m}
                \frac{\pi^*d\mu}{g^*d\lambda_Y}
                     \longrightarrow-\log\max_{Z_y}|\sigma|^2.
\]
The fibre measure is finite and has dense positive locus.
After normalization to total mass one, its power means increase to
the maximum.  Proper smoothness makes that maximum continuous.
Integration over a fixed patch on which both factors are nonzero
gives local upper bounds for the rescaled psh weights.  The limit is
therefore psh.  No positivity of the zeroth moment is needed.

Regularity of $\sigma$ on the proper model gives a local lower bound
for the limiting weight up to the boundary.
Above each base prime a component attaining the vertical minimum has
$\sigma\ne0$ at its general point.  In transverse coordinates the map
is $y_1=x_1^{e_E}$, with the other base coordinates free.  A neighbourhood
there meets every nearby fibre and bounds its maximum away from zero.
This gives the upper bound needed for psh extension at general points
of every omitted divisor.  Psh extension across the remaining
codimension-two set then applies; local upper boundedness there follows
from transverse discs with boundary off the removed set.  The lower
bound persists.  The extended weights are locally bounded and obey the
frame transformation rule for $N$.
\end{proof}

\begin{proposition}[Weighted exceptional correction]
\label{g:bt:weighted-transfer}
Under \eqref{g:bt:all-exceptional-rank}, $N$ has a semipositive metric
with locally bounded weights.  For some integer $c\geq0$,
$-K_Y+cB$ has a semipositive metric satisfying
$|s_B|^{2c}e^{-\psi}\in L^1_{\mathrm{loc}}$.
A nonzero section of any positive divisible multiple of
$j(N+R)+lB$, with $j>0$ and $l\geq0$, gives a nonzero invariant section
of a positive multiple of $L$ on $X$.
\end{proposition}

\begin{proof}
If $Y$ is a point, all the base divisors vanish and their metrics are
positive constants. The section $s$ is then regular on $Z$ and pushes
down to the required section on $X$. Assume henceforth that $\dim Y>0$.

For every $m\geq0$, \eqref{g:bt:all-exceptional-rank} applied to
$D=K_{Z/X}$ gives invariant rank one for
$g_*(K_{Z/Y}+(ma+1)\pi^*L)$.  Its generator is the product of $s^m$,
the invariant Jacobian section, and a local inverse canonical frame of
$Y$.  In a local frame of $N$, let $\sigma$ be the normalized section
from \eqref{g:bt:normalization}.  The generator has squared norm
\[
 I_m(y)=\int_{Z_y}|\sigma|^{2m}\,
                         \frac{\pi^*d\mu}{g^*d\lambda_Y},
\]
where $d\mu$ is the smooth volume defined by the metric of $L$.
Lemma~\ref{g:moments} applies: Berndtsson positivity is used on the
smooth base-change locus and the invariant line is a holomorphic
orthogonal summand.  Lemma~\ref{g:bt:bounded-maximum} then gives the
bounded metric of $N$.

For the second metric use $\rho=I_0$.  Its negative logarithm is psh on
the smooth locus, and $\rho\,d\lambda_Y=g_*(\pi^*d\mu)$ has finite mass.
At a general point of a prime $E\mapsto Q$, choose transverse coordinates
with $y_1=x_1^{e_E}$ and the other base coordinates free.
If $a_E=\operatorname{ord}_E K_{Z/X}$, integration on a small patch gives
\[
                 \rho(y)\geq C|y_1|^{2((1+a_E)/e_E-1)}.
\]
Outside $B$ choose a nonexceptional component, so $a_E=0$.
On the finitely many components of $B$ choose $c$ large enough that
\[
                     \psi=-\log\rho+c\log|s_B|^2
\]
is locally bounded above at their general points.  The same bound is
already available at every other omitted divisor.  Psh extension across
these divisors, then across the remaining set of codimension at least
two, gives the metric on $-K_Y+cB$.  Its weighted integral is exactly
the original finite pushforward mass.

Finally, represent a section of a positive divisible multiple of
$j(N+R)+lB$ by a rational function $v$ on $Y$. After incorporating
that multiple in $j,l$, its divisor satisfies
\[
             \operatorname{div}(v)+j(N+R)+lB\geq0.
\]
Return the rational section $s^jf^*v$. At a nonexceptional component
$E$ above $Q\not\subset B$, its order is at least
$j(h_E-e_En'_Q)\geq0$, where $h_E=\operatorname{ord}_Es$.
There are no nonexceptional components above $B$, and
equidimensionality excludes vertical primes over higher-codimension
sets. Horizontal orders are nonnegative by the generic regularity of
$s$. Thus every pole is exceptional over $X$ and disappears on
pushdown. All factors are invariant.
\end{proof}

\subsection{Unweighted toroidal transfer}
\label{g:bt:toroidal}

The preceding correction uses only a divisorially equidimensional model.
A toroidal model gives a finer rational correction and an unweighted
integral.  The rank hypothesis in the next proposition includes zero.

\begin{proposition}[Integrable toroidal correction]
\label{g:bt:toroidal-transfer}
Use the smooth metric and action above.  Let $p:W\to X$ be birational
with $W$ normal, and let $f:W\to Y$ be projective, equidimensional and
toroidal for strict toroidal embeddings, with $Y$ smooth.  Include the
exceptional locus of $p$ in the source boundary.  Suppose that an
invariant rational section $s$ of $aL$, $a>0$, is regular on the generic
fibre and that on a smooth equivariant resolved diagram
\begin{equation}\label{g:bt:zeroth-rank}
 \dim_K H^0\bigl(F,(K_{Z/X}+ma\pi^*L)|_F\bigr)^T=1
                           \qquad(m\geq0).
\end{equation}
Normalize $s$ by vertical minima as in \eqref{g:bt:normalization}.
For a boundary prime $Q$ put
\[
 A_E=1+\operatorname{ord}_E K_{W/X},\qquad
 c_Q=\min_{E\mapsto Q}\frac{A_E}{e_E},\qquad
 \Theta=\sum_Q\max\{c_Q-1,0\}Q.
\]
Then $N$ has a semipositive metric with locally bounded weights, and
$-K_Y+\Theta$ has a semipositive metric with locally integrable squared
frame norms.  The pullback $f^*\Theta$ is exceptional over $X$.
\end{proposition}

\begin{proof}
For a point base, $N=\Theta=0$ and both metrics are positive constants.
We may therefore assume $\dim Y>0$.

Lemma~\ref{g:moments} supplies the psh moment weights, and
Lemma~\ref{g:bt:bounded-maximum} gives the metric of $N$.
The zeroth moment is the pushed-forward smooth volume density
$\rho$.  It can be computed on a toroidal resolution of $W$, even when
that model is not equivariant: change of variables identifies its
density with the one on an equivariant resolution.
Put $\theta_Q=\max\{c_Q-1,0\}$ and use
\[
                      \psi=-\log\rho+
                                   \sum_Q\theta_Q\log|s_Q|^2.
\]
At a minimizing component the transverse calculation in the preceding
proof gives $\rho\geq C|s_Q|^{2(c_Q-1)}$.  This is the local upper
bound needed to extend $\psi$ at general divisor points.
Off the boundary the morphism is smooth and gives a positive lower
bound.  Extend over the remaining codimension-two set by psh removability.

The correction $\Theta$ is the divisor $A_0$ in
Lemma~\ref{g:volume-subklt}. To compare the two definitions, it remains
only to check primes outside the base boundary. A vertical prime above
such a prime cannot belong to the source boundary of a toroidal
morphism. Since that boundary contains the exceptional locus, the
prime is nonexceptional over $X$; the morphism is smooth at its general
point. Its Jacobian order is zero and its pullback multiplicity is one,
so the corresponding coefficient of $A_0$ is zero.

The rank-free discrepancy lemma now gives both the exceptional support
of $f^*\Theta$ and, on every smooth model $Z'\to W$ with maps
$\pi':Z'\to X$ and $g':Z'\to Y$,
\[
       \operatorname{coeff}_F(K_{Z'/X}-(g')^*\Theta)>-1
                         \qquad\text{for every prime }F\subset Z'.
\]
On a simultaneous normal-crossing resolution, these are precisely the
exponents of the corrected source measure in the projection formula
\[
 e^{-\psi}\,d\lambda_Y
   =g'_*\left(\prod_Q|s_Q\circ g'|^{-2\theta_Q}
                                  (\pi')^*d\mu\right).
\]
They make that measure locally integrable. Proper pushforward proves
the required unweighted integrability on $Y$, including at boundary
intersections. Curvature here came from the invariant zeroth moment;
the discrepancy argument used no adjoint-rank assumption.
\end{proof}

No horizontal strict-curvature patch was assumed here.
The output is the integrability used in
Proposition~\ref{g:bt:unweighted-slope}.  The separate toroidal-volume
theorem includes the additional patch needed by its effectivity
applications.

\subsection{The large-moment construction for eventual rank one}
\label{g:bt:eventual}

We now construct the donor metric from a positive moment. The hypothesis
will initially specify adjoint rank only at sufficiently large degrees.
Let $F_X=\mathbb C(X)$, let $K\subseteq F_X^T$ be finitely generated and
relatively algebraically closed in $F_X$, and let $s$ be a nonzero
invariant rational section of $dL$, $d>0$.
For a prime divisorial valuation $E$ over $X$, write
$a_E=\operatorname{ord}_E K_{X'/X}$ on a smooth model carrying $E$.
Assume
\begin{equation}\label{g:bt:horizontal-regular}
 \operatorname{ord}_E s\geq0
       \quad\hbox{whenever $\operatorname{ord}_E|_{K^*}=0$},
\end{equation}
and, for every sufficiently large integer $m$,
\begin{equation}\label{g:bt:eventual-rank}
 \left\{t\in\operatorname{Rat}(mdL)^T:
   \begin{array}{l}
    \operatorname{ord}_E t+a_E\geq0\\[-2pt]
    \text{for all $E$ trivial on $K$}
   \end{array}\right\}=K s^m.
\end{equation}
Here $\operatorname{Rat}(mdL)$ is the one-dimensional $F_X$-space of
rational sections, with zero included in the displayed set.
Both conditions concern all divisorial valuations.  On one smooth
resolution of the map to a model of $K$, they can be tested at
horizontal primes: the divisor inequalities pull to further resolutions
because the new relative canonical divisor is effective.

Together these two conditions imply rank one at every nonnegative
multiple of $d$, including zero. Indeed, if $t$ has degree $jd$, with
$j\geq0$ an integer, and satisfies
the displayed adjoint inequalities, choose $m\geq j$ in the eventual
range. Horizontal regularity of $s$ makes $ts^{m-j}$ satisfy the same
inequalities at degree $md$. Equation~\eqref{g:bt:eventual-rank} gives
$ts^{m-j}\in Ks^m$, hence $t\in Ks^j$; the reverse inclusion follows
from regularity of $s$ and $a_E\geq0$. The distinction of the following
construction is its metric: it uses one positive moment and retains
the corresponding multiple of the normalized section line.

Choose a model $p:W\to X$, $f:W\to Y$ with $Y$ smooth projective,
$W$ normal and $\mathbb Q$-Gorenstein klt, and every vertical prime
mapping onto a base prime.  Such a model is obtained by toroidalization
and subdivision.  Refine target cones by source-cone images; on the
induced source subdivision each ray maps to a ray or to zero.
The reduced boundary of the initial smooth toroidal source defines a
positive slicing function on its cone complex, linear on each cone and
positive away from the origin. Its pullback remains
positive on the induced subdivision, including on horizontal rays.
Abramovich--Rojas \cite[Corollary~0.11]{AbramovichRojas1998} then gives
a coherent simplicial subdivision without new rays. Since the target
cones are smooth, every new source cone still maps to a target face;
equidimensionality is preserved. The resulting toroidal source has
quotient singularities, hence is klt. These are birational modifications,
with no field extension \cite{ADK2013,AK2000,KKMSD}.

For each base prime $Q$, put
\[
 h_E=\operatorname{ord}_E s,\qquad e_E=\operatorname{ord}_E f^*Q,\qquad
 n_Q=\min_{E\mapsto Q}h_E/e_E,\qquad N=\sum_Qn_Q Q,
\]
and $D=\operatorname{div}_W(s)-f^*N\geq0$.  Define
\[
 r_Q=\min_{\substack{E\mapsto Q\\D_E=0}}a_E/e_E,\qquad
                       R=\sum_Qr_Q Q\geq0.
\]
The minimum is over a nonempty finite set.  Only finitely many
coefficients of $N,R$ are nonzero.

\begin{proposition}[An integrable large moment]
\label{g:bt:large-moment}
There is an integer $k$ in the range of \eqref{g:bt:eventual-rank} such
that
\[
            \Xi=K_{W/X}+kD-f^*R\geq0.
\]
The rational line $N$ has a semipositive metric with locally bounded
weights, and $J=-K_Y+kN+R$ has a semipositive metric with
$\mathcal I(\psi)=\mathcal O_Y$.
A nonzero section of a positive divisible multiple of $N+R/k$ returns
to a nonzero invariant section of a positive multiple of $L$ on $X$.
\end{proposition}

\begin{proof}
At a component with $D_E=0$, the definition of $r_Q$ gives
$a_E-e_Er_Q\geq0$.  At the finitely many other relevant components,
increase $k$ to absorb $e_Er_Q$ by $kD_E$.  At horizontal components
there is no subtraction.  This proves $\Xi\geq0$, and over every $Q$
some component has $\Xi_E=0$.

On a smooth equivariant model over $Y$ define
\[
 I_m\,d\lambda_Y=g_*\bigl(|\pi^*s|^{2m}\pi^*d\mu\bigr).
\]
For all sufficiently large $m$, \eqref{g:bt:eventual-rank} identifies
the invariant adjoint line generated by the Jacobian times $s^m$.
Smooth direct-image positivity and orthogonal compact averaging imply
that $-\log I_m$ is psh on a common smooth open.
The densities are smooth there, so the base-change locus for an
individual $m$ can be removed by continuity.  High moments give
\[
 I_m^{1/m}\longrightarrow\beta(y)=\max_{Z_y}|\pi^*s|^2.
\]
For this limit one may divide by $I_0$ as a positive density on the
smooth open; no curvature assertion about $I_0$ is used.
The weights of the required lines are
\begin{equation}\label{g:bt:large-weights}
 \begin{split}
 \varphi_N&=-\log\beta+\sum_Q n_Q\log|s_Q|^2,\\
 \psi&=-\log I_k+\sum_Q(kn_Q+r_Q)\log|s_Q|^2.
 \end{split}
\end{equation}
The first extends with locally bounded weights by regularity of the
normalized section and a minimizing component above every divisor,
by Lemma~\ref{g:bt:bounded-maximum}.

For the second, choose a component above $Q$ with $\Xi_E=0$.
In coordinates $y_1=x_1^{e_E}$ at its general point, the corrected
source density has exponent
\[
          a_E+kh_E-e_E(kn_Q+r_Q)=0.
\]
Its pushforward is bounded below by $C|y_1|^{-2(1-1/e_E)}$.
Thus $\psi$ is bounded above at general points of every omitted
divisor and extends psh there, then across codimension two.
On a log resolution $v:\widehat W\to W$ the exponents in the corrected
source density are the coefficients of
\[
 K_{\widehat W/X}+kv^*D-v^*f^*R
                   =K_{\widehat W/W}+v^*\Xi.
\]
They are all greater than $-1$, since $W$ is klt and $\Xi$ is effective
and $\mathbb Q$-Cartier.  Normal crossing integration and properness
give $e^{-\psi}\in L^1_{\mathrm{loc}}$ on the entire base.

Finally $D-f^*R/k$ is nonnegative at each component nonexceptional
over $X$, since $a_E=0$ there and $\Xi\geq0$.  Pulling back a section
of $N+R/k$ and multiplying by a power of $s$ leaves only exceptional
poles, removed on pushdown.
\end{proof}

The resulting donor line is $J=-K_Y+kN+R$. Its $kN$ term records the
positive moment used to construct it. Although the hypotheses also
imply rank one at zero, this proof obtains integrability from $I_k$;
it does not produce the same metric on $-K_Y+R$.

\subsection{Signed anti-log-canonical classes}
\label{g:bt:signed}

The third slope input starts from a nef class rather than a weighted
integral.  It will be applied to the signed discriminant produced by the
toroidal density calculation.  Its negative part measures exactly the
possible loss in the cotangent estimate.

\begin{proposition}[Signed-boundary cotangent estimate]
\label{g:bt:signed-slope}
Let $Y$ be smooth projective and let $B$ be a rational divisor with simple
normal crossing support and coefficients less than one.  Write
$B=B^+-B^-$ in disjoint positive and negative parts, and suppose that
$-(K_Y+B)$ is nef.  Then
\[
                         hB^- -c_1(M)
                         \quad\hbox{is pseudoeffective}
\]
for every line subsheaf $M\subseteq(\Omega_Y^1)^{\otimes h}$, $h\geq0$.
If $\dim Y>0$ and $\alpha\ne0$ is a movable class with
$\alpha\cdot B^-=0$, this gives
$\mu_{\alpha,\min}(T_Y)\geq0$ and $c_1(M)\cdot\alpha\leq0$.
\end{proposition}

\begin{proof}
The class $-K_Y+B^-=-(K_Y+B)+B^+$ is pseudoeffective.
For an algebraically integrable foliation $\mathcal F$ of positive rank
and corank, resolve its rational first-integral map together with $B$ by
$u:Y'\to Y$, and write
\[
             K_{Y'}+C=u^*(K_Y+B),\qquad C=C^+-C^-.
\]
Every coefficient of $C$ is less than one: the positive part of the
original simple normal crossing boundary is klt, and decreasing its
coefficients can only increase discrepancies.  In particular
$(Y',C^+)$ is klt.  For an ample Cartier $H$ on $Y'$ and rational
$\delta>0$, choose a general effective rational divisor
\[
        A_\delta\sim_{\mathbb Q}
                    -u^*(K_Y+B)+\delta H
\]
such that $(Y',C^++A_\delta)$ is klt.  The class is ample, so a sufficiently
divisible multiple supplies such a member by Bertini, also on a general
fibre of the resolved map.  We have
\[
        K_{Y'}+C^++A_\delta\sim_{\mathbb Q} C^-+\delta H.
\]
It is big on that fibre and hence has a nonzero section after taking a
positive divisible multiple.  Apply Druel's relative pseudoeffectivity
theorem \cite[Proposition~4.1]{Druel2017}, in the form of
\cite[Proposition~3.1]{Ou2023}, to conclude that
$K_{\mathcal F'}+C^++A_\delta$ is pseudoeffective.
Its pushforward is numerically
\[
                 K_{\mathcal F}-K_Y+B^-+\delta u_*H.
\]
Letting $\delta$ decrease to zero proves
$K_{\mathcal F}-K_Y+B^-$ pseudoeffective.
Lemma~\ref{g:bt:slope-principle} now proves the assertion for every
movable class, and therefore in the pseudoeffective cone.
\end{proof}

The distinction between the three slope inputs is now explicit.
Proposition~\ref{g:bt:weighted-slope} retains the integral of a specified
boundary section.  Proposition~\ref{g:bt:unweighted-slope} uses an
unweighted integral and the class $K_Y+J$.
Proposition~\ref{g:bt:signed-slope} uses nefness of an anti-log-canonical
class with a signed boundary.  Each supplies the cotangent inequality
needed for the corresponding contraction; none asserts a global
strict-curvature lower bound.

\section{Choosing a smaller base}\label{g:reductions}
We now prove allowed-pole nonvanishing using the two metrics of
Section~\ref{g:analytic} and the weighted cotangent estimate of
Proposition~\ref{g:bt:weighted-slope}. The first two proofs apply when
there are no positive-degree holomorphic forms; the third gives a
different choice of base in the rationally connected case. All three
start from effective divisors with a fixed cotangent twist and must
produce both a horizontal section and every required adjoint rank.

\subsection{The common section conclusion}\label{rb:reductions}

\begin{theorem}[Allowed-pole nonvanishing]\label{g:allowed}\label{rb:allowed}
Let $V$ be a smooth connected projective variety over $\mathbb C$ with
$H^0(V,\Omega_V^j)=0$ for every $j>0$.  Let $D\geq0$ be integral and
suppose $-K_V+D$ carries a semipositive singular metric whose measure
$|s_D|^2$ has finite mass.  If a line bundle $L$ carries a semipositive
metric with bounded weights, then
\[
                    H^0(V,aL+cD)\ne0
\]
for some integers $a>0$ and $c\geq0$.  The same statement holds for a
rational line bundle $L$, with $a$ clearing its denominator.
\end{theorem}

In the three proofs below, whenever a dominant rational map from $V$
with connected general fibers is selected, use the model diagram of
Section~\ref{g:models}, with smooth projective base $Y$.
Thus $\bar p:\bar V\to V$ is birational,
$\bar f:\bar V\to Y$ is surjective and equidimensional with $\bar V$
normal, and a smooth resolution $W\to\bar V$ has composite maps
$p:W\to V$ and $f:W\to Y$. We write $J=K_{W/V}$.

\subsection{Maximal systems dominated by the target and boundary}
\label{rb:dominated}

\begin{proof}[First proof of Theorem~\ref{rb:allowed}]
Clear the denominator of $L$.  We induct on $\dim V$; a point is
immediate.  Call a nonempty linear system in a line bundle $P$
admissible if
\[
               aL+cD-P\quad\hbox{is pseudoeffective}
\]
for some integers $a>0,c\geq0$.  The constant system is admissible,
and products of admissible systems are admissible.  Choose one whose
image dimension is maximal.  If it is generically finite, then $P$ is
big.  Consequently $aL+cD$ is big, which supplies a section after a
positive multiple.  Otherwise, resolve its map and take the relative
algebraic closure of its image field to obtain $f:W\to Y$ with
connected general fibres and $\dim Y<\dim V$.

Every ratio of sections in every admissible system belongs to
$\mathbb C(Y)$.  Indeed adjoining such ratios cannot increase image
dimension, since products of the systems remain admissible.
The ratios are therefore algebraic over the chosen image field and
belong to its relative algebraic closure.  This argument also covers
subsystems and pencils.

Lemma~\ref{rb:determinant} gives $N_i\sim M+m_iL$, while
Proposition~\ref{rb:cotangent} gives $-M+kD$ pseudoeffective for some
$k\geq0$.  Systems in all positive combinations of the $N_i$ are
therefore admissible.  For $i<j<l$ the two effective divisors
\[
 (m_l-m_j)N_i+(m_j-m_i)N_l,
 \qquad (m_l-m_i)N_j
\]
are linearly equivalent, so the rational function relating them lies
in $\mathbb C(Y)$.  After pullback, their horizontal parts are equal.
It follows that the horizontal multiplicity at every prime is affine
in $m_i$.  Fixing $i<j$ and letting $l$ tend to infinity shows that its
slope is nonnegative. Thus $N_j-N_i$ defines a rational section $\rho$ of
$(m_j-m_i)L$ with no horizontal poles on $W$, and hence on $\bar V$.
If $Y$ is a point, it has no poles and already proves the theorem.

For a positive-dimensional base we verify the rank-one condition.
Let $P$ denote the original maximal system and let $A_0$ be the
pullback to $Y$ of its image hyperplane class.  It is big and
$p^*P-f^*A_0$ has an effective representative.  Given a sufficiently
ample base divisor $H$, a positive multiple of $A_0$ dominates $H$
up to an effective divisor.  Therefore $p_*f^*H$ is dominated in the
pseudoeffective order by a positive multiple of $L$ plus a nonnegative
multiple of $D$.

If $J+p^*D+mp^*L$ has two independent sections on a generic fibre,
spread them out and clear their vertical poles using $f^*H$ with $H$
sufficiently ample.  Their pushdowns are sections of
$D+mL+p_*f^*H$; the exceptional Jacobian $J$ contributes zero under
pushdown.  This system is admissible, but its ratio is nonconstant on
the generic fibre, contrary to the maximality property.  Hence its
generic dimension is at most one for every $m\geq0$.  For the multiples
of the exponent just constructed, $s_{J+p^*D}\rho^l$ supplies a nonzero
generic section, including $l=0$.  Lemmas~\ref{rb:rank-one} and
\ref{rb:transfer} apply.

The base has no positive-degree holomorphic forms: pullback by the
dominant map is injective, and forms on the smooth birational model
identify with forms on $V$.  The descended data thus satisfy exactly
the hypotheses of the theorem in lower dimension.  Induction gives a
section with poles on $D_Y$, and Lemma~\ref{rb:transfer} carries it to
a section with poles only on $D$.  An integral multiple of $D$ clears
those finitely many poles.
\end{proof}

\subsection{A supporting class and maximal combinations}
\label{rb:zero-degree}

Here is a different choice of the dimension-reducing map.  It uses
only effective combinations of the determinant divisors and the
permitted boundary, instead of all pseudoeffectively dominated
systems.

\begin{proof}[Second proof of Theorem~\ref{rb:allowed}]
Again induct on dimension and clear denominators.  If $L+D$ is big,
a positive multiple has a section.  Otherwise $L+D$ lies on the
boundary of the pseudoeffective cone.  Cone duality gives a nonzero
movable class $\alpha$ such that
\[
                       L\cdot\alpha=D\cdot\alpha=0.
\]
The zero-degree part of Proposition~\ref{rb:cotangent} and the determinant
construction give $N_i\sim M+m_iL$ with $M\cdot\alpha\leq0$.
Effectivity forces $N_i\cdot\alpha=M\cdot\alpha=0$.
No nonnegative integral combination of $D,N_1,N_2,\ldots$ is big.

Choose such a combination $G$ of maximal Iitaka dimension $d<\dim V$.
Take a multiple realizing dimension $d$, resolve its map, and take the
connected-fibre base $Y$.  Every ratio from the linear system of any
allowed combination lies in $\mathbb C(Y)$, since adding systems and
multiplying sections cannot increase $d$.  If $H$ is any base divisor,
then
\begin{equation}\label{rb:ample-domination}
                      ap^*G-f^*H\sim E_H\geq0
\end{equation}
for some sufficiently large $a$: the resolved system contains the
pullback of the image hyperplane class, whose transform on $Y$ is big.

We need the following stronger form of the ratio assertion.  A rational
function whose horizontal poles on $W$ are bounded by $p^*G'+E$, where
$G'$ is an allowed combination and $E$ is effective and $p$-exceptional,
belongs to $\mathbb C(Y)$.  Its remaining poles are vertical.  Choose
an effective sufficiently ample $H$ whose pullback bounds them; a
base divisor can contain the images of all these finitely many poles,
including images of codimension at least two, to the necessary orders.
By \eqref{rb:ample-domination}, the function is a ratio of sections of
$p^*(G'+aG)+E$.  Normality gives
$p_*\mathcal O_W(E)=\mathcal O_V$, so the exceptional term adds no
sections.  The original ratio assertion proves the claim.

The three-index relations among the $N_i$ now imply affine horizontal
multiplicities exactly as in the first proof.  Let $\rho$ have divisor
$N_j-N_i$ and exponent $k=m_j-m_i>0$.  It has no horizontal poles,
and its horizontal zeros are bounded by $p^*N_j$.  On a generic fibre
$s_{J+p^*D}\rho^l$ is a nonzero section of
$J+p^*D+lkp^*L$.  The ratio of any other section to this one has
horizontal poles bounded by
\[
                        J+p^*D+l p^*N_j.
\]
The stronger ratio assertion makes that ratio a base function.  Thus
the generic section space is one-dimensional for every $l\geq0$.
If $d=0$, $\rho$ is already regular.  Otherwise
Lemmas~\ref{rb:rank-one} and \ref{rb:transfer} give the two required
metrics and the controlled boundary on the smaller base.  Absence of
holomorphic forms descends, so dimension induction completes the
proof.
\end{proof}

\subsection{Maximal fibre spaces on a rationally connected variety}
\label{rb:max-fibration}

For rationally connected varieties one can choose a maximal fibre
space by two numerical and generic-section properties.  We give this
construction because its rank-one conclusion follows by enlarging a
field, rather than by fixing one class of systems in advance.

\begin{proposition}\label{rb:rc-method}
The conclusion of Theorem~\ref{rb:allowed} on a rationally connected
$V$ follows by dimension induction through rationally connected
bases, using fibre spaces that satisfy the following two conditions:
for an ample divisor $H$ on the base,
\[
 aL+cD-p_*f^*H\ \hbox{is pseudoeffective for some }a>0,c\geq0,
\]
and a positive multiple of $p^*L$ has a nonzero section on the generic
fibre without additional $D$-poles.
\end{proposition}

\begin{proof}
Suppose no positive multiple of $L$ has a section with poles only on
$D$.  In particular no $aL+cD$ with $a>0,c\geq0$ is big.
Rational connectedness gives absence of holomorphic forms, so take
the determinant divisors $N_i\sim M+m_iL$ and the inequality
$-M+kD$ pseudoeffective.  Let $K$ be the relative algebraic closure in
$\mathbb C(V)$ of the field generated by functions with divisors
\[
 (m_l-m_j)N_i+(m_j-m_i)N_l-(m_l-m_i)N_j,
                           \qquad i<j<l.
\]
These functions exist because the displayed divisor class is zero.
A finite selection gives the same transcendence degree, and the
relative algebraic closure is a finite extension of a finitely
generated subfield.  Thus $K$ defines a smooth projective base after
resolution.

Each of those functions is represented by a pencil whose total
bundle has class $uL+vM$, with $u,v>0$.  Finitely many such pencils give
a generically finite map from the base to their combined image; the
sum of the pulled-back hyperplane classes is therefore big on the
base.  The resolved pencils are bounded above by their total bundles.
Pushing down and using $-M+kD$ pseudoeffective proves the first
condition of the proposition.  A large multiple of the big base
class dominates any chosen ample divisor, so the same condition holds
with an ample divisor.

On the generic fibre all the displayed principal divisors vanish.
The effective restrictions of the $N_i$ therefore have affine
multiplicities in $m_i$.  Their slope is nonnegative by the unbounded
sequence, so $N_j-N_i$ supplies the second condition.  A point base
would already give a regular section, contrary to our supposition.
A base of full dimension would make the pushdown of an ample base
class big, hence would make some $aL+cD$ big, also impossible.
We have found a base of dimension strictly between zero and $\dim V$.

Choose, among fibre spaces with these two properties, one with maximal
base dimension.  The properties survive higher models.  For the first,
pullback of an ample class to a birational base model is big and a
large multiple dominates every fixed divisor; pulling and pushing
preserve effectivity.  The second is the regularity of a fixed rational
section over the generic fibre, which survives resolution.

Let $\rho$ be the supplied generic section of $kp^*L$.  If
$J+p^*D+mkp^*L$ had two independent generic sections, their ratio would
be transcendental over the relatively algebraically closed base field.
Adjoining it and taking relative algebraic closure enlarges the base
dimension.  Clear the two sections' vertical poles by a base twist
$f^*H'$.  Their ratio pencil has class
\[
                    J+p^*D+mkp^*L+f^*H'.
\]
Add an ample system from the old base.  The product map supplies a big
class on the enlarged base.  On pushdown $J$ vanishes, and the first
property for the old map bounds the remaining class by a multiple of
$L$ plus a multiple of $D$.  Thus the first property holds for the
enlarged map.  Restricting $\rho$ to its new generic fibre gives the
second property and stays nonzero at the generic point.  This
contradicts maximality.

The spaces in Lemma~\ref{rb:rank-one} consequently have dimension at
most one and have a nonzero element $s_{J+p^*D}\rho^m$, also at $m=0$.
Transfer gives the required data on the smaller base, which is
rationally connected because it is a smooth projective model of a
rational image of $V$.  Induction gives its permitted section and
transfer gives one on $V$, the required contradiction.
\end{proof}

\section{Logarithmic orbit integrals and invariant sections}
\label{g:orbits}\label{rb:torus}

The invariant problem uses the natural action on $-K_V=\det T_V$.
Twisting this action by a character changes the invariant subspace and
is not part of the statement.
Both proofs below construct the two weights required by
Lemma~\ref{g:transfer}: a pushforward-volume weight and a fiber-supremum
weight. The first obtains their plurisubharmonicity from invariant
adjoint rank one; the second obtains it by logarithmic orbit integration.
The second calculation is independent of invariant direct-image
curvature, while its final section construction still uses the
allowed-pole induction already proved.

\begin{theorem}\label{rb:invariant}
Let $V$ be a smooth connected projective variety over $\mathbb C$ with
$H^0(V,\Omega_V^j)=0$ for all $j>0$.  Suppose $-K_V$ has a smooth
Hermitian metric of semipositive curvature.  For any algebraic torus
$T$ acting on $V$, there is an integer $m>0$ such that
\[
                         H^0(V,-mK_V)^T\ne0
\]
for the natural anticanonical action.
\end{theorem}

\subsection{The rank-one quotient construction}

\begin{proof}
Average local weights, with their line-bundle transition rule, over the
compact form of $T$ using the natural linearization.  This gives a
smooth semipositive metric invariant under that compact group.
If the action is trivial, Theorem~\ref{rb:allowed} applies with
$D=0$ and $A=L=-K_V$.  Otherwise take the rational quotient with field
$\mathbb C(V)^T$ \cite{Rosenlicht1956,BellGhiocaReichstein2017}, and its connected-fibre smooth
projective model $Y$.  The invariant field is relatively algebraically
closed in $\mathbb C(V)$: an element algebraic over it has a finite
orbit, and the connected torus acts trivially on that orbit.  General
fibres contain a dense orbit.

Choose the model diagram equivariantly, by modifying the trivially
acted-on base, normalizing the dominating graph, and taking an
equivariant resolution.  Write $p:W\to V$ and $f:W\to Y$ for the smooth resolution maps,
$J=K_{W/V}$, and $\bar V$ for the normal equidimensional intermediate
model of Section~\ref{g:models}. Let $r$ be the orbit dimension.  A wedge of
$r$ independent infinitesimal action fields is an invariant regular
relative anticanonical section on the generic fibre; commutativity of
$T$ gives invariance.  A rational frame of $-K_Y$ and the canonical
Jacobian section $s_J$ put it in
\[
 p^*(-K_V)=-K_{W/Y}+f^*(-K_Y)+J.
\]
It has no horizontal poles. Regard this section of $p^*(-K_V)$
birationally as a rational section $\rho$ on $V$.  The Jacobian section
is invariant for the induced canonical actions.  The ratio of any two
invariant sections of $J+mp^*(-K_V)$ is a rational invariant function,
hence belongs to the base field.  The sections $s_J(p^*\rho)^m$ supply the
nonzero element at every $m\geq0$, including $m=0$.  Thus the invariant
rank-one condition holds.

If $Y$ is a point, the full wedge of infinitesimal fields on $V$ already
gives a global invariant anticanonical section.  Otherwise the invariant
case of Lemma~\ref{rb:rank-one} and Lemma~\ref{rb:transfer} gives a
bounded semipositive base bundle and a finite-volume donor metric.
The base has no holomorphic forms by dominant pullback, and
Theorem~\ref{rb:allowed} gives a section with permitted base poles.
Its transferred section is invariant, and its only permitted poles on
$V$ lie in $D=0$.  It is therefore regular.
\end{proof}

\subsection{A direct integral argument on logarithmic orbits}

We next establish the two psh functions needed for transfer by an
integral calculation, independently of the invariant direct-image
curvature argument.  The calculation uses the variance estimate of
Brascamp--Lieb \cite{BrascampLieb1976}; we include the needed proof on a
convex ball to make the limiting procedure explicit.
The integral calculation is a complex Pr\'ekopa argument, related to
Berndtsson's formulations \cite{Berndtsson1998,Berndtsson2006}; the
supremum assertion is the corresponding minimum principle of Kiselman
\cite{Kiselman1978}. The tube-domain formulation in
\cite[Theorem~1.1]{DengWangZhangZhou2018} gives a direct comparison for
that second assertion. We prove both statements here because the
almost-everywhere finite integrals and the exhaustion limits are
precisely the data used by the transfer construction.

\begin{lemma}[Integral and supremum weights]\label{g:prekopa}\label{rb:prekopa}
Let $r\geq0$ be an integer, let $U$ be a complex manifold, and let $\Phi(z,x)$ be a smooth function
on $U\times\mathbb R^r$ such that $\Phi(z,\operatorname{Re}w)$ is
plurisubharmonic on $U\times\mathbb C^r$.  Suppose
\[
                     0<\int_{\mathbb R^r}e^{-\Phi(z,x)}\,dx<\infty
\]
for almost every $z$.  Then the almost-everywhere functions
\[
 -\log\int_{\mathbb R^r}e^{-\Phi(z,x)}\,dx,
 \qquad
 -\log\sup_{x\in\mathbb R^r}e^{-\Phi(z,x)}
\]
have psh representatives, whenever the second expression is taken on
the locus where its supremum is finite and positive.  Under the stated
integrability assumption that supremum is finite and positive for
almost every $z$.
\end{lemma}

\begin{proof}
For $r=0$, give $\mathbb R^0$ its unit mass. Both expressions equal $\Phi(z)$, so the assertion is immediate. Assume $r>0$.
Plurisubharmonicity implies convexity of $x\mapsto\Phi(z,x)$.
An integrable positive smooth log-concave function on $\mathbb R^r$
is bounded.  To see this, choose a fixed ball on which its value is at
least $c>0$.  If values $a_j$ at points $x_j$ tended to infinity,
log-concavity on the midpoint ball between $x_j$ and that fixed ball
would give a lower bound $(ca_j)^{1/2}$ on a ball of a fixed positive
volume.  Its integral would tend to infinity, a contradiction.
This proves the last assertion.

Fix a bounded real ball $D$ and first replace $\Phi$ by
$\Phi+\epsilon|x|^2$, so the Hessian in $x$ is positive definite. Until $\epsilon$ is sent to zero,
write $\Phi$ for this regularized weight. Restrict the base variable to a complex disc.  Let $\mathbb E$ denote
expectation for the probability density proportional to $e^{-\Phi}$
on $D$.  Differentiating under the integral gives
\begin{equation}\label{rb:variance}
 \frac{\partial^2}{\partial z\partial\bar z}
          \left(-\log\int_D e^{-\Phi}\,dx\right)
   =\mathbb E\Phi_{z\bar z}
       -\mathbb E|\Phi_z-\mathbb E\Phi_z|^2.
\end{equation}
For a smooth real weight $b$ whose Hessian $b_{xx}$ is positive definite
on $\bar D$, put
\[\mathbb E_b F=\frac{\int_D F e^{-b}dx}{\int_D e^{-b}dx}.\]
The Brascamp--Lieb bound is
\begin{equation}\label{rb:bl}
 \mathbb E_b|v-\mathbb E_bv|^2
       \leq\mathbb E_b\bigl(b_{xx}^{-1}(\nabla v,\nabla\bar v)\bigr).
\end{equation}
Here it is enough to prove it for smooth real mean-zero $v$ and then
add the estimates for real and imaginary parts. The bounded-domain
form and its Neumann identity are also proved in
\cite[Theorems~1.1 and~1.2(1), and Section~3.3]{KolesnikovMilman2017}.
The preceding addition of \(\epsilon|x|^2\) provides the stated positive
definiteness; the limit \(\epsilon\downarrow0\) is taken afterwards.

Solve the weighted Neumann problem
\[
             (-\Delta+\nabla b\cdot\nabla)u=v,
             \qquad \partial_{\mathbf n}u=0.
\]
The mean-zero condition is precisely its solvability condition on the
ball.  Weighted integration by parts and the differentiated equation
give
\[
 \int_Dv^2e^{-b}
  =\int_D|\nabla^2u|^2e^{-b}
     +\int_Db_{xx}(\nabla u,\nabla u)e^{-b}
     +\int_{\partial D}\mathrm{II}(\nabla u,\nabla u)e^{-b}.
\]
The boundary second fundamental form $\mathrm{II}$ is nonnegative
because $D$ is convex, and $\nabla u$ is tangential there.  Also
$\int_Dv^2e^{-b}=\int_D\nabla v\cdot\nabla u\,e^{-b}$.
Cauchy--Schwarz for the quadratic form $b_{xx}$ bounds the latter
integral by
\[
 \left(\int_Db_{xx}^{-1}(\nabla v,\nabla v)e^{-b}\right)^{1/2}
 \left(\int_Db_{xx}(\nabla u,\nabla u)e^{-b}\right)^{1/2}.
\]
Combining the two identities gives \eqref{rb:bl}, including the case
of zero variance by continuity.

Apply \eqref{rb:bl} with $b(x)=\Phi(z,x)$ and $v=\Phi_z$.  The complex Hessian of
$\Phi(z,\operatorname{Re}w)$ is positive semidefinite.  Its Schur
complement gives
\[
         \Phi_{z\bar z}\geq
         \Phi_{zx}(\Phi_{xx})^{-1}\Phi_{x\bar z}.
\]
The factors $1/2$ in complex vertical derivatives cancel against the
factor $1/4$ in the vertical Hessian.  Equations~\eqref{rb:variance}
and \eqref{rb:bl} prove plurisubharmonicity of
$-\log\int_D e^{-\Phi}dx$.  Decreasing $\epsilon$ to zero preserves
this conclusion.  Exhausting $\mathbb R^r$ by balls gives a decreasing
sequence of psh functions; its almost-everywhere finite limit is the
first required representative.

For the supremum, use the high-power integral argument of
\cite[p.~1635, following Theorem~1.2]{Berndtsson2006}: take the
probability-normalized Lebesgue measure on $D$ and replace $\Phi$ by
$l\Phi$, with integers $l\geq1$. The functions
\[
       -\frac1l\log\left(\frac1{|D|}\int_D e^{-l\Phi}dx\right)
\]
are psh and decrease with $l$, by monotonicity of $L^l$ norms.  Their
finite limit is $-\log\sup_D e^{-\Phi}$.  Increasing the balls again
gives a decreasing psh limit, finite almost everywhere by the first
paragraph.  This is the second representative.
\end{proof}

\begin{proof}[Orbit-integral proof of Theorem~\ref{rb:invariant}]
If the torus acts trivially, apply Theorem~\ref{g:allowed} with
$D=0$ and $A=L=-K_V$. Assume the action is nontrivial.
Give $L=-K_V$ the same compact-torus-invariant smooth semipositive
metric and let $\mu$ be its smooth positive volume. Use the maps
$p:W\to V$, $f:W\to Y$ and intermediate model $\bar V$ recalled
in the first proof.  Divide out the
ineffective kernel of the torus action.  The effective torus has
trivial generic stabilizer.  Indeed a subgroup fixing the generic
point would act trivially on the function field and hence on $V$;
subgroups of a split torus are determined by their character lattices,
so the geometric generic stabilizer is such a fixed subgroup.
The generic orbit over the rational quotient field is consequently
a torsor under a split torus.  Hilbert's Theorem~90 trivializes it.
Spreading out gives a dense orbit family $T\times Y^0$ over a nonempty
base open set.  These quotient and triviality statements use
\cite{Rosenlicht1956,BellGhiocaReichstein2017} for the quotient and
\cite[Corollary~3.47 and Proposition~3.50]{MilneAlgebraicGroups2017}
for split-torus torsor triviality.

On this product take logarithmic coordinates
$t_j=e^{x_j+i\theta_j}$ on $T$, with $x\in\mathbb R^r$ and
$\theta\in(\mathbb R/2\pi\mathbb Z)^r$.  Wedge the logarithmic action fields
with a base coordinate anticanonical frame to obtain a frame $\sigma$
of $L$ along the product.  Its weight
\[
                  \Phi(y,x)=-\log|\sigma|^2
\]
is smooth and psh in $(y,x+i\theta)$, and is independent of $\theta$
by compact invariance.  In one coordinate,
\[
 \frac{i}{2}\frac{dt_j}{t_j}\wedge
             \frac{d\bar t_j}{\bar t_j}=dx_j\wedge d\theta_j.
\]
The top form dual to the logarithmic anticanonical frame therefore
identifies the volume with a fixed positive constant times
$e^{-\Phi}\,dx\,d\theta\,d\lambda_y$. Angular integration contributes
a positive constant independent of $y$.
The complement of the orbit family is algebraic and has measure zero.
Consequently the pushforward density of $\mu$ is
\[
                  g_0(y)=C\int_{\mathbb R^r}e^{-\Phi(y,x)}dx
\]
almost everywhere.  It is finite almost everywhere by Fubini.

Choose a nonzero rational base anticanonical frame.  It differs from
the coordinate frame by a base factor and gives an invariant rational
section $\rho$ of $L$.  For almost every general fibre, the supremum of
$|\rho|^2$ on its dense orbit is finite by Lemma~\ref{rb:prekopa}.
It follows that $\rho$ has no horizontal pole on the normal intermediate model $\bar V$:
a pole along a horizontal prime would make its norm unbounded near
that prime on a general compactified fibre, and the dense orbit
approaches its general points.  The smooth metric is comparable to a
positive smooth frame there, so such unboundedness cannot be hidden
in the norm. Normality therefore makes the section regular on the
complete generic fiber of \(\bar V\to Y\), and its pullback to \(W\)
remains regular.

The supremum on the dense orbit equals the essential supremum for
transfer, since the rational section is now regular on the generic
fibre and the metric is continuous.  Lemma~\ref{rb:prekopa} proves
both $-\log g_0$ and $-\log S_\rho$ psh on the quotient open set;
nonzero base frame factors contribute only pluriharmonic terms.
If the quotient is a point, the same no-horizontal-pole argument
already says that $\rho$ is a regular invariant section. Assume now
$\dim Y>0$. Apply Lemma~\ref{g:transfer} with
$D=0$, $A=L=-K_V$, $k=1$, and this rational section $\rho$.
The preceding calculation supplies its two psh weights. Vertical
normalization, properness, and the bounded input metric give the local
upper bound for the normalized section norm; a component attaining
the vertical minimum gives the complementary lower bound. Thus the
transfer produces a bounded metric on its line $B$ and a finite-volume
metric on the different line $-K_Y+D_Y$.
Theorem~\ref{g:allowed} applies on $Y$, which inherits the absence of
positive-degree holomorphic forms by dominant pullback. The output has
the explicit form $(\psi\circ f)(p^*\rho)^m$ on $W$ for a rational
base function $\psi$. Its birational pushdown to $V$ is therefore invariant. Since the initial boundary was
$D=0$, the support conclusion of transfer removes all divisorial poles
on $V$. This completes the orbit-integral proof.
\end{proof}

\subsection{The rationally connected variant}
\label{rb:rc-invariant}
If $V$ is rationally connected, the first proof also closes using only
Proposition~\ref{rb:rc-method} as its nonvanishing input.  Here is the
exact substitution.  For the trivial action, apply that proposition
with $A=L=-K_V$ and $D=0$.  For a positive-dimensional orbit, construct
$\rho$ and the invariant rank-one spaces exactly as above, retaining
$m=0$.  The rational quotient base $Y$ is rationally connected, since
it is a smooth projective model of a dominant rational image of $V$.
The two metric lemmas give $B$ with bounded semipositive weights and
$A_Y=-K_Y+P$ with finite distinguished volume, where $P$ is the
allowed-prime divisor.  Thus Proposition~\ref{rb:rc-method} applies on
$Y$ with target $B$ and boundary $P$.  Its section transfers to an
invariant section on $V$, and the support assertion gives no pole
because the original boundary is zero.  The point quotient is handled
by the full wedge of action fields, as in the first proof.  Consequently
the maximal-fibre-space argument supplies invariant anticanonical
nonvanishing on every rationally connected smooth projective $V$ with
a smooth semipositive anticanonical metric.

\section{Section-ratio fields and horizontal regularity}
\label{g:cf:fields}

The contraction in Theorem~\ref{g:cf:contraction-theorem} uses a
maximal linear system. There is an
intrinsic description of its base field: it consists of the rational
functions whose zeros and poles have degree zero against the supporting
movable class. This identification lets us compare two ways of producing
the new generic section. Affine interpolation uses the precise relation
between three twists. A second argument uses finite horizontal support
and the ordering of nonnegative integer vectors. We first separate the
regularity and uniqueness conditions that both arguments must preserve.

\subsection{Regularity on the original variety and on a resolution}
\label{g:cf:conditions}

Let $X$ be smooth connected projective, let an algebraic torus $T$ act on
$X$, and let $L$ be a $T$-linearized line bundle. Fix a relatively
algebraically closed subfield
$\mathbb C\subseteq K\subseteq\mathbb C(X)^T$. A divisorial valuation is
\emph{horizontal over $K$} when it is zero on $K^*$; on a proper model
mapping to a projective model of $K$, this means that its center dominates
the base. All orders of sections on higher models refer to pullbacks of
their line bundles.

Suppose $s$ is an invariant rational section of $aL$, where $a>0$, whose
pullback is regular on the complete generic fiber of a resolution of the
map defined by $K$. This is a condition on every horizontal divisorial
valuation, including exceptional ones. It is unchanged by further
resolution, because a regular section pulls back regularly; the converse
follows by pushforward and normality. The uniqueness condition below,
in contrast, tests the other rational sections only at prime divisors of
$X$ itself.

\begin{lemma}[Two equivalent uniqueness conditions]
\label{g:cf:uniqueness}
With these data, the following statements are equivalent.
\begin{enumerate}
\item For every integer $j\geq0$ and every nonzero invariant rational
section $u$ of $jaL$ with no poles at the $K$-horizontal primes of $X$,
one has $u/s^j\in K$.
\item For every integer $b\geq0$, the ratio of any two nonzero invariant
rational sections of $bL$ having no poles at the $K$-horizontal primes of
$X$ belongs to $K$.
\end{enumerate}
Both statements include degree zero.
\end{lemma}

\begin{proof}
The second assertion implies the first by comparing $u$ with $s^j$;
the regularity hypothesis on $s$ guarantees that this comparison is
permitted. Conversely, let $u_1,u_2$ have degree $b$. The first assertion
applied to $u_i^a$ gives $u_i^a/s^b\in K$. Hence
$(u_1/u_2)^a\in K$, and relative algebraic closedness puts $u_1/u_2$ in
$K$. This argument also applies when $b=0$.
\end{proof}

These equivalent formulations do not identify regularity at original
primes with regularity at all upstairs primes. The latter is the stronger
regularity requirement on the chosen section $s$. In particular, a
uniqueness assertion that tests only sections already regular at every
upstairs horizontal prime is not, without another argument, the
uniqueness assertion in Lemma~\ref{g:cf:uniqueness}.

For $L=-K_X$, these data imply the adjoint rank hypotheses of
Theorem~\ref{g:cf:contraction-theorem}. Indeed, on a smooth resolution
$\pi:Z\to X$, $g:Z\to Y$, put $R=K_{Z/X}$. Dividing an invariant
section of $R+ja\pi^*L$ over the generic point of $Y$ by the Jacobian
section gives a rational section of $jaL$ without poles at horizontal
primes of $X$: their strict transforms have coefficient zero in $R$.
Lemma~\ref{g:cf:uniqueness} makes it a $K$-multiple of $s^j$.
Conversely $s^j s_R$ is present. Thus the invariant rank is one for
every $j\geq0$. This implication uses the original-prime test; adjoint
rank alone does not license replacing it by a test on arbitrary rational
sections with exceptional poles.

\subsection{The null-prime field is the maximal-system field}

\begin{lemma}\label{g:cf:null-field}
Let $Y$ be a smooth connected projective variety of dimension $d>0$ and
let $\alpha\ne0$ be a movable numerical curve class. Call a prime divisor
$P$ \emph{null} if $\alpha\cdot P=0$, and set
\[
 K_\alpha=\{0\}\cup\{u\in\mathbb C(Y)^*:
          \operatorname{Supp}(\operatorname{div}u)
          \text{ consists of null primes}\}.
\]
Then $K_\alpha$ is relatively algebraically closed in $\mathbb C(Y)$
and has transcendence degree less than $d$. It is the field of the Stein
base of any linear system of degree zero against $\alpha$ having maximal
image dimension among all such systems. Every prime of $Y$ vertical over
$K_\alpha$ is null. If $\operatorname{Pic}(Y)$ is finitely generated,
only finitely many null primes are horizontal over $K_\alpha$.
\end{lemma}

\begin{proof}
Products and inverses preserve the support condition. The poles of a
nonzero sum remain on null primes. Since a principal divisor has degree
zero, its effective zero divisor then has degree zero as well; movability
makes each of its components null. Thus $K_\alpha$ is a field.

Equivalently, every valuation at a non-null prime is trivial on
$K_\alpha$. Such a valuation remains trivial on an element algebraic over
that field. An element of positive valuation cannot satisfy a polynomial
with nonzero constant term over the field, since that term would uniquely
have minimal valuation. Applying the same argument to the inverse excludes
negative valuation. This proves relative algebraic closedness.

If $u_1,\ldots,u_d\in K_\alpha$ were algebraically independent, the
sum $P$ of their polar divisors would be null. The monomials of total
degree at most $m$ in the $u_i$ are linearly independent sections of
$mP$, so $P$ would be big. A nonzero functional nonnegative on the
pseudoeffective cone is strictly positive on its interior. Thus
$\alpha\cdot P>0$, a contradiction.

Every ratio in a zero-degree system belongs to $K_\alpha$, since the
zero divisors of its sections are effective of degree zero. Conversely,
if $u\in K_\alpha^*$ and $P=(\operatorname{div}u)^-$, then $1$ and
$u$ are sections of $P$ and $\alpha\cdot P=0$. Choose a maximal-image
zero-degree system and write $K_0$ for the field of its Stein base.
The product-system argument of Lemma~\ref{g:cf:maximal-field} puts
all zero-degree section ratios in $K_0$, hence $K_\alpha\subseteq K_0$.
The generating ratios of the chosen system lie in $K_\alpha$; its
relative algebraic closedness gives the reverse inclusion. This proves
the claimed equality, rather than just a comparison of dimensions.

A non-null prime is horizontal over $K_\alpha$ by definition. Finally,
if there were infinitely many horizontal null primes, finite generation
of $\operatorname{Pic}(Y)$ would give a nonzero integral relation among
finitely many of their classes (multiply to remove torsion). The resulting
nonzero divisor would be principal. Its defining function would lie in
$K_\alpha$ but have nonzero order at one of these horizontal primes,
which is impossible.
\end{proof}

The hypotheses for finiteness hold on every smooth base dominated by the
variety in Theorem~\ref{g:cf:contraction-theorem}: pullback and birational
invariance give $H^0(Y,\Omega_Y^1)=0$, hence $\operatorname{Pic}^0(Y)=0$.
Such subfields are finitely generated: choose a finite transcendence basis;
the relative algebraic closure of its rational function field inside the
finitely generated field $\mathbb C(Y)$ is finite. Consequently the field
has a smooth projective model. The same argument will apply to the fields
constructed below.

\subsection{Minima on all components and on original components}
\label{g:cf:minima}

We now retain the fixed-section data of
Section~\ref{g:cf:conditions}, with $L=-K_X$ smoothly semipositive and
$H^0(X,\Omega_X^q)=0$ for $q>0$. Assume also either of the equivalent
uniqueness conditions in Lemma~\ref{g:cf:uniqueness}. Let $Y$ be a smooth projective model
of $K$. Prepare a normal graph
\[
 X\xleftarrow{p}W\xrightarrow{f}Y
\]
whose fibers have dimension at most $\dim X-\dim Y$.
This preparation can be made without toroidalization. On the flat locus
of an initial graph, the fibers define a map to the projective Hilbert
scheme of $X$. Resolve that map on the base and pull back the universal
family. Its fixed Hilbert polynomial bounds the dimension of every fiber.
The main graph transform lies in this family, and its normalization is
finite over the transform. This gives the asserted bound. The graph and
its normalization are equivariant; take an equivariant
smooth resolution $\rho:Z\to W$ afterwards
\cite{RaynaudGruson1971}, and write $\pi=p\rho$ and $g_Z=f\rho$ for
its composite maps to $X$ and $Y$. In particular every vertical
prime of $W$ dominates a prime of $Y$.

For a prime $P\subset Y$, let $Q$ range over the components above it on
$W$, and put $e_Q=\operatorname{ord}_Q f^*P$. Define
\[
 d_P=\min_{Q\mapsto P}\frac{\operatorname{ord}_Q(s)}{e_Q},\qquad
 d_P^*=\min_{\substack{Q\mapsto P\\Q\text{ nonexceptional over }X}}
                       \frac{\operatorname{ord}_Q(s)}{e_Q}.
\]
A prime $P$ is \emph{omitted} if all components above it are exceptional
over $X$; set $d_P^*=d_P$ in this case. There are finitely many omitted
primes. The divisors
\[
 D=\sum_Pd_PP,\qquad D^*=\sum_Pd_P^*P
\]
have finite support, $D^*\geq D$, and
\[
 \operatorname{div}_W(s)-f^*D\geq0.
\]
The last divisor is rational Cartier, so it pulls back effectively to
$Z$. The first minimum controls the metric because some component has
zero normalized section order. The second controls sections on $X$
because a pole on an exceptional divisor disappears on pushdown.

\begin{lemma}\label{g:cf:finite-volume-base}
The rational divisor $D$ has a semipositive metric with locally bounded
weights. There is an effective integral divisor $P_0$, supported on the
omitted primes, and a psh metric weight $\phi$ on $-K_Y+P_0$ such that
$|s_{P_0}|^2e^{-\phi}$ is locally integrable.
\end{lemma}

\begin{proof}
If $Y$ is a point, then $D=0$; take $P_0=0$ and positive Hermitian
norms on the two trivial lines. Their weights are bounded and their
adjoint measure is finite. We may therefore assume $\dim Y>0$.
Apply Lemmas~\ref{g:moments} and~\ref{g:transfer} with original
variety $X$, allowed-pole divisor zero, $A=L=-K_X$, and both metrics equal
to the compact-averaged smooth metric on $L$. The section is $s$ of degree
$a$. It has no horizontal poles on the normal equidimensional graph.
The rank computation following Lemma~\ref{g:cf:uniqueness} verifies the
smooth invariant adjoint rank condition for every $j\geq0$. Thus the
two required psh functions are supplied by Lemma~\ref{g:moments}.

The full-component normalization in Lemma~\ref{g:transfer} is exactly
$D$; a denominator-clearing multiple of $D$ receives its bounded metric,
which scales back to the rational divisor. Since the allowed-pole divisor
on $X$ is zero, the invisible primes in that lemma are precisely our
omitted primes. Its correction is therefore an effective integral $P_0$
with the asserted support. The distinguished measure on $-K_Y+P_0$
equals the pushforward of the smooth anticanonical volume almost
everywhere, so its local integral is finite, including at boundary
intersections. This proves both assertions on the flattened graph,
without requiring toroidalization.
\end{proof}

\subsection{A supporting class and finite horizontal support}

\begin{lemma}\label{g:cf:supporting-class}
Assume the data of Section~\ref{g:cf:minima} and suppose no positive power of \(L\) has a nonzero
invariant global section. Then there is a nonzero movable class \(\alpha\)
on \(Y\) satisfying
\[
 \alpha\cdot D=\alpha\cdot D^*=0,
 \qquad \alpha\cdot\Gamma=0
             \quad\text{for every omitted prime }\Gamma.
\]
In particular, \(d^*_\Gamma=d_\Gamma\) whenever
\(\alpha\cdot\Gamma>0\).
\end{lemma}

\begin{proof}
The bounded psh metric makes \(D\) pseudoeffective.  Alternatively,
apply Nadel vanishing and regularity to its divisible multiples:
local boundedness makes their multiplier ideals trivial, and
\(mD+A\) is effective for fixed sufficiently positive \(A\).

The divisor \(D^*\) cannot be rationally linearly equivalent to an
effective rational divisor modulo the span of omitted primes.
Otherwise, after clearing denominators, there would be \(N>0\)
and \(g\in K^*\) with
\[
           N d^*_\Gamma+\operatorname{ord}_\Gamma(g)\geq0
                 \quad\text{at every nonomitted prime}.
\]
Then \(s^Ng\) would have no pole along any prime divisor on \(X\).
For horizontal divisors this follows from the generic regularity of $s$.
For a vertical divisor, its nonexceptional strict transform on
\(W\) dominates a nonomitted prime of \(Y\), and the displayed
inequality and the minimum defining \(d^*_\Gamma\) give the result.
Thus \(s^Ng\) would be a nonzero invariant global section.

Let \(\mathcal S\subset N^1(Y)_{\mathbb R}\) be the span of omitted primes,
and project the cone of effective divisor classes to
\(N^1(Y)_{\mathbb R}/\mathcal S\).  Denote this convex cone by \(C\).
Its closure has nonempty interior, since the ample cone already has
nonempty interior before projection.  The image of \([D^*]\)
belongs to \(\overline C\), by pseudoeffectivity of \(D\) and
effectivity of \(D^*-D\).  It cannot lie in its interior.  Indeed,
the interior of a convex cone agrees with that of its closure and
is contained in the cone itself.  Membership of this rational point
in \(C\) would express it as a finite nonnegative real combination
of rational effective classes.  The corresponding finite rational
linear system then has a nonnegative rational solution.  Since
numerical and rational linear equivalence agree on \(Y\), that
would give the rational effectivity just excluded.

A supporting functional at this boundary point is nonzero,
nonnegative on \(C\), and zero on the image of \([D^*]\).
Pull it back to \(N^1(Y)_{\mathbb R}\).  It is nonnegative on
effective classes and annihilates \(\mathcal S\); by cone duality
\cite{BDPP2013}, it is a nonzero movable class \(\alpha\).
Pseudoeffectivity of \(D\) and effectivity of \(D^*-D\) give
\(0\leq\alpha\cdot D\leq\alpha\cdot D^*=0\).  The effective
difference therefore has zero degree; each of its components has
zero degree.  This proves the final assertion as well.
\end{proof}


For clarity, the numerical-to-linear step in this proof uses
$\operatorname{Pic}^0(Y)=0$: the kernel of the map from the N\'eron--Severi
group to numerical divisor classes is torsion, so rational numerical
relations among the indicated divisors are rational linear relations.
The movable class just constructed annihilates $P_0$ in
Lemma~\ref{g:cf:finite-volume-base}. Apply
Proposition~\ref{g:bt:weighted-slope} with that effective correction, its
integrable distinguished section, and $\alpha\cdot P_0=0$. Every
cotangent line has nonpositive $\alpha$-degree.
Lemma~\ref{g:determinant}, applied to a Cartier multiple of $D$,
gives a fixed cotangent determinant and effective integral divisors
\begin{equation}\label{g:cf:null-divisors}
 N_m\sim M+mD,\qquad m\longrightarrow\infty,\qquad
 \alpha\cdot N_m=0,
\end{equation}
for one fixed line bundle $M$. All indices clear the denominators of $D$.

The field $K_\alpha$ from Lemma~\ref{g:cf:null-field} is the same smaller
field that maximal zero-degree systems give. We now use a different
section-extraction argument on that field, keeping its separate input:
finitely many horizontal null primes. This is the relative counterpart
of the finite-support comparison in \cite[Lemma~2.1]{LP2018}; the
uniqueness verification after extraction is an additional step here.

\begin{proposition}[Finite-support extraction]\label{g:cf:dickson}
Under the data of Section~\ref{g:cf:minima}, assume no positive multiple of
$L$ has an invariant global section. There are a positive integer $a'$ and
an invariant rational section $s'$ of $a'L$ satisfying both the generic
regularity condition and either uniqueness condition of
Lemma~\ref{g:cf:uniqueness} over $K_\alpha$. Moreover
$\operatorname{trdeg}K_\alpha<\operatorname{trdeg}K$.
\end{proposition}

\begin{proof}
Choose one higher equivariant resolution \(\widehat Z\to Z\) on
which the rational map for \(K_\alpha\) is also a morphism.  Pull back
the divisors \(N_m\).  Their parts horizontal for \(K_\alpha\) have
support in a fixed finite set of prime divisors on \(\widehat Z\).
Indeed, a horizontal component mapping onto a prime of \(Y\)
must lie over one of the finitely many horizontal null primes in
Lemma~\ref{g:cf:null-field}.  There are only finitely many components
over those primes.  A prime whose image in \(Y\) has codimension
at least two must be exceptional over \(W\), by the fibre dimension
bound, and there are only finitely many such divisors on this fixed
resolution.

List the horizontal multiplicities of the pulled-back \(N_m\) as
vectors in \(\mathbb Z_{\geq0}^r\), using this fixed finite support.
Order the indices \(m\) increasingly.  Dickson's lemma, in its sequence
form (see \cite[Theorem~1.1(1) and Lemma~3.2]{AichingerAichinger2018}), gives
\(m_2>m_1\) such that the earlier vector is componentwise at most
the later one.  For completeness, the needed sequence form follows
by induction on \(r\): a sequence of nonnegative integers either
has an infinitely repeated value or has a strictly increasing
subsequence; apply this successively to each coordinate.  This
also covers empty horizontal support.  Hence
\(N_{m_2}-N_{m_1}\), pulled back to \(\widehat Z\), is effective
on every \(K_\alpha\)-horizontal divisor.

Put \(b=m_2-m_1>0\).  From \eqref{g:cf:null-divisors}, choose
\(g_0\in K^*\) satisfying
\[
          N_{m_2}-N_{m_1}=bD+\operatorname{div}_Y(g_0).
\]
Set
\begin{equation}\label{g:cf:new-section}
                         a'=ba,\qquad s'=s^bg_0.
\end{equation}
The section \(s'\) is invariant.  On \(\widehat Z\), its divisor
is the pullback of \(b(\operatorname{div}_Z(s)-g_Z^*D)\) plus the pullback of
\(N_{m_2}-N_{m_1}\).  It is therefore nonnegative along all
\(K_\alpha\)-horizontal divisors.  The finitely many remaining poles
lie over proper closed subsets of the new base, so a dense open
avoids them.  This proves generic regularity of the new section.

It remains to verify the ratio requirement, which is essential for
another iteration.  Let \(k\geq0\), and let \(t\) be an invariant
nonzero rational section of \(ka'L\) without poles along prime
divisors on \(X\) horizontal for \(K_\alpha\).  Since \(K_\alpha\subseteq K\),
every \(K\)-horizontal valuation is \(K_\alpha\)-horizontal.  The old
fixed-section uniqueness condition therefore gives
\[
                             g=\frac{t}{s^{kb}}\in K.
\]
Let \(\Gamma\) be a non-null prime of \(Y\).  It is nonomitted,
and \(d^*_\Gamma=d_\Gamma\) by
Lemma~\ref{g:cf:supporting-class}.  Each nonexceptional component
\(Q\) over it on \(W\) is horizontal for \(K_\alpha\): its restriction
to \(K_\alpha\) is \(e_Q\operatorname{ord}_\Gamma=0\).  Its valuation
is the valuation of a prime divisor on \(X\).  Consequently
\[
  e_Q\operatorname{ord}_\Gamma(g)+kb\operatorname{ord}_Q(s)
                         =\operatorname{ord}_Q(t)\geq0.
\]
Taking the minimum over these nonexceptional components proves
\begin{equation}\label{g:cf:ratio-inequality}
                   \operatorname{ord}_\Gamma(g)+kb\,d_\Gamma
                   \geq0.
\end{equation}
The rational divisor \(\operatorname{div}_Y(g)+kbD\) has total
\(\alpha\)-degree zero.  Its coefficients at every non-null prime
are nonnegative by \eqref{g:cf:ratio-inequality}, and the remaining
primes have degree zero.  Each of those nonnegative coefficients
must therefore be zero.  Moreover,
\(\operatorname{div}_Y(g_0)+bD=N_{m_2}-N_{m_1}\) is supported on
null primes.  Subtracting \(k\) times this identity shows that
\[
       \frac{t}{(s')^k}=\frac{g}{g_0^k}\in K_\alpha.
\]
This includes \(k=0\), so fixed-section uniqueness is preserved in
its full stated range.


Strict transcendence-degree decrease follows from
Lemma~\ref{g:cf:null-field}. The regularity just proved holds on the
complete generic fiber and hence survives further modification. This
completes the new data, including the original-prime uniqueness test.
\end{proof}

For the same supporting movable class, affine interpolation in
Lemma~\ref{g:cf:horizontal-slope} and
Proposition~\ref{g:cf:dickson} therefore reach the same field by different
arguments. The first proves that the horizontal multiplicities are affine
functions of the twist. The second only needs finite horizontal support
and chooses a comparable pair of multiplicity vectors. Neither method
replaces the subsequent uniqueness verification.

\subsection{A toric construction of the initial section}

The rational torus quotient can also be initialized using the geometry of
its complete generic fiber. This gives a second construction of existence;
its essential input is toric generation and a weight argument, rather than
the wedge of fundamental vector fields.

\begin{lemma}\label{g:cf:toric-initial-section}
Let an algebraic torus $T$ act on a smooth connected projective variety
$X$, and suppose that $L=-K_X$ is nef. On an equivariant smooth projective
resolution $\pi:Z\to X$, $g:Z\to Y$ of the rational quotient
$\mathbb C(Y)=\mathbb C(X)^T$, some positive power of $\pi^*L$ has a
nonzero invariant section on the complete generic fiber, for its natural
linearization.
\end{lemma}

\begin{proof}
The invariant field is relatively algebraically closed in $\mathbb C(X)$:
an element algebraic over it has a finite orbit, and the divisible abelian
group $T(\mathbb C)$ has no nontrivial finite quotient. Rosenlicht's
quotient theorem gives a dense orbit on the smooth connected general
fiber $F$ of $g$ \cite{Rosenlicht1956}. The stabilizer of that orbit fixes
its closure, so $F$ is toric for an effective quotient of $T$. The nef
bundle $\pi^*L|_F$ has a basepoint-free positive multiple by the toric
basepoint-freeness theorem \cite[Theorem~6.3.12]{CLS2011}.

For any one-parameter subgroup $\lambda$, the Bia\l{}ynicki-Birula
decomposition of $X$ has an open attracting stratum as $t\to0$, whose
fixed component has nonnegative tangent weights, and an open attracting
stratum as $t\to\infty$, whose tangent weights are nonpositive
\cite{BB1973}. Choose $F$ very general so that its image meets both
strata for every $\lambda$; there are only countably many integral
cocharacters. The corresponding orbit limits on the proper $F$ give
fixed points with the two weak signs for the fiber weights of $\pi^*L$,
since $L=\det T_X$ has its natural action. Evaluation of a basepoint-free
multiple at either point is equivariant and surjective. Thus its section
space has weights of both weak signs for every $\lambda$.

It follows that zero is in the convex hull of the torus characters in
this section space: otherwise a strict separating functional could be
chosen rational and then integral. Choose a rational convex combination
with sum zero and clear denominators. The product of the corresponding
nonzero eigensections is a nonzero invariant section of a positive power;
nonzeroness uses integrality of $F$. This is the invariant-sections
construction used in the toric setting of \cite[Section~3]{Muller2025}.

For each positive integer $m$, cohomology and base change and torus
weight splitting identify the invariant section space of $m\pi^*L$ on
fibers over a dense open with a vector bundle of its generic rank.
Choose the very general fiber also in all these countably many opens.
The invariant section just found therefore implies a nonzero invariant
generic rank for its particular exponent. This yields the required
section over $\mathbb C(Y)$. The empty-dimensional fiber and the trivial
torus are included in the argument.
\end{proof}

\begin{remark}
The same weight argument can be made directly on the geometric generic
fiber. Its generic point avoids the complements of the two attracting
strata for every chosen one-parameter subgroup, and properness supplies
the orbit limits. Toric generation and separation of the character
polytope again produce a weight-zero product. Invariants for a split
torus commute with field extension, so this section descends from the
algebraic closure of \(\mathbb C(Y)\). This gives a geometric-generic
form of the preceding construction, without choosing a very general
closed fiber.
\end{remark}

\section{Evaluation metrics and fields with controlled poles}
\label{g:field-variants}

The previous section extracted a generic section from divisors of zero
movable degree. We now vary the choice of the smaller base. Evaluation
in an adjoint direct image supplies a transverse slope estimate.
We apply it to systems generated by four fixed divisor classes,
checking that their maximality still restores uniqueness after
contraction. A separate construction then uses a field of functions
with controlled poles to encode the divisors that may be absorbed
on the original variety.

\subsection{Evaluation metrics and cotangent slopes}
\label{g:cf:evaluation-subsection}

There is a second way to obtain the slope inequality needed for
contraction.  It uses evaluation in an adjoint direct image over a
\emph{further} rational quotient of the base.  This gives a metric on
the determinant of its transverse tangent sheaf, with a correction
supported only on the omitted base primes.

\begin{proposition}\label{g:cf:evaluation-slope}
Let \(X,Y\) be smooth connected projective varieties, let
\(L=-K_X\) have a smooth semipositive metric, and let an algebraic
torus \(T\) act on \(X\).  Suppose that \(f:X\dashrightarrow Y\) is
dominant and invariant and that every invariant rational function on
\(X\) having no poles at the \(f\)-horizontal prime divisors of
\(X\) belongs to \(\mathbb C(Y)\).  Choose a normal projective graph
model
\[
 X\xleftarrow{p}W\xrightarrow{\varphi}Y
\]
on which every vertical prime maps onto a prime divisor of \(Y\).
Let \(N\) be the reduced sum of the primes of \(Y\) above which all
prime divisors of \(W\) are exceptional over \(X\).

For every dominant rational map \(g:Y\dashrightarrow U\) to a smooth
projective variety, let \(\mathcal F\subset T_Y\) be the saturated
sheaf tangent to its general fibers.  There is an integer \(c\geq0\)
such that
\[
 \det(T_Y/\mathcal F)+cN
 \quad\text{is pseudoeffective}.
\]
Consequently, if \(\alpha\ne0\) is a movable curve class on \(Y\)
with \(\alpha\cdot N=0\), then every line subsheaf
\(A\subset(\Omega_Y^1)^{\otimes j}\), \(j\geq0\), satisfies
\(\alpha\cdot c_1(A)\leq0\).
\end{proposition}

The uniqueness assumption in this statement is precisely the
degree-zero instance of Lemma~\ref{g:cf:uniqueness}.  In particular, the
evaluation argument does not require a nef anticanonical class on
\(Y\).

\begin{proof}
Average the metric on \(L\) over the compact form of \(T\), so that
it is invariant.  When \(U\) is a point the determinant assertion is
immediate.  Otherwise replace \(U\) birationally, without changing
\(\mathcal F\), so that every \(g\)-vertical prime \(P\subset Y\)
maps onto a prime \(Q\subset U\).  To arrange this, flatten the
dominating graph and resolve the base \cite{RaynaudGruson1971}.
Each $g$-vertical prime of $Y$ has a strict transform on the
resulting equidimensional source. Properness defines the rational
lift from $Y$ at its generic point, so the resulting map sends
that prime onto a base prime. Put
\(e_P=\operatorname{ord}_P(g^*Q)\).  At the general point of \(P\),
the map has coordinates
\[
 (x,y_2,\ldots,y_{\dim Y})
 \longmapsto (x^{e_P},y_2,\ldots,y_{\dim U}).
\]
The determinant of the differential, after quotienting by its
saturated kernel, therefore has order \(e_P-1\).  Along a horizontal
prime it has order zero.  Thus, in codimension one and hence as line
bundles on \(Y\),
\begin{equation}\label{g:cf:evaluation-determinant}
 \det(T_Y/\mathcal F)=-g^*K_U-R_g,
 \qquad R_g=\sum_{P\text{ vertical}}(e_P-1)P.
\end{equation}
The pullback here is defined on the domain of the rational map and
extended across its codimension-two complement.  The sum is finite,
since its coefficients are the divisorial orders of one determinant
map.

Take a smooth projective equivariant resolution dominating \(W\)
and resolving both rational maps, with morphisms
\[
 \pi:Z\longrightarrow X,\qquad q:Z\longrightarrow Y,
 \qquad r:Z\longrightarrow U.
\]
Write \(e_\pi\) for the invariant Jacobian section of
\(K_Z+\pi^*L\).  On a dense open of \(U\), the sheaf
\[
 \mathcal E=\bigl(r_*(K_{Z/U}+\pi^*L)\bigr)^T
\]
is locally free, commutes with base change, and has its adjoint
\(L^2\) metric.  Berndtsson's theorem gives Griffiths semipositivity
\cite{Berndtsson2009}.  The same is true of the invariant summand:
the weight-zero projection is holomorphic and orthogonal, because
the compact torus preserves the coefficient metric and the fiber
integral.  Its metric is therefore the quotient metric.  This
argument concerns a smooth coefficient metric on the smooth-family
open.

Choose a rational frame for \(\mathcal E\) and a rational frame
\(\eta\) of \(K_U\).  Its sections on \(Z\) have the form
\begin{equation}\label{g:cf:evaluation-coefficients}
 (q^*\phi_i)e_\pi/(r^*\eta),\qquad
 \phi_i\in\mathbb C(Y).
\end{equation}
Indeed, dividing first by the Jacobian and the inverse base frame
gives invariant rational functions on \(X\).  These have no poles
at the primes dominating \(Y\): such primes also dominate \(U\),
and the Jacobian is a unit at their strict transforms.  The assumed
degree-zero uniqueness places the functions in \(\mathbb C(Y)\).

After shrinking dense opens, evaluation of the coefficients in
\eqref{g:cf:evaluation-coefficients} defines a holomorphic functional
on \(g^*\mathcal E\).  For \(v\in\mathcal E_{g(y)}\), denote its
coefficient at \(y\) by \(\phi_v(y)\).  In the frame
\(g^*\eta^{-1}\), set
\begin{equation}\label{g:cf:evaluation-weight}
 \psi(y)=\log\sup_{\|v\|_{g(y)}=1}|\phi_v(y)|^2.
\end{equation}
This is psh, being the log squared norm of a holomorphic section of
the dual of a Griffiths-semipositive bundle.  Replacing \(\eta\) by
\(a\eta\) adds \(\log|g^*a|^2\) to \(\psi\), exactly the
transformation law for a weight on \(-g^*K_U\).  The section
\(e_\pi/(r^*\eta)\) has coefficient \(1\), so this weight is not
identically \(-\infty\).

It remains to account for ramification and for the Jacobian zeros
above the boundary.  Fix a prime \(P\subset Y\), choose a prime
\(D\subset Z\) mapping onto it, and put
\[
 l=\operatorname{ord}_D(q^*P),\qquad
 d_\pi=\operatorname{ord}_D(e_\pi).
\]
If \(P\not\subset N\), choose \(D\) nonexceptional over \(X\),
so that \(d_\pi=0\).  Suppose first that \(P\) maps onto a prime
of \(U\).  With \(j=le_P\), suitable coordinates around a general
point of \(D\) express \(r\) as
\[
 (\tau,z_2,\ldots,z_n)\longmapsto
 (\tau^j,z_2,\ldots,z_h),
 \qquad h=\dim U,\quad D=(\tau=0).
\]
Take \(\eta\) regular and nonvanishing on this base patch.  If a
fiber section \(v\) has coefficient \(a_v\) in a regular frame of
\(K_{Z/U}+\pi^*L\), its adjoint fiber norm contains the density
\[
 \frac{|a_v|^2}{|j\tau^{j-1}|^2}
 |dz_{h+1}\cdots dz_n|^2
\]
up to positive bounded factors with bounded inverse.  The
denominator comes from dividing the total coordinate top form by
the pulled-back base form.  For \(\|v\|=1\), holomorphic submean
on a smaller fiber polydisc gives
\( |a_v|\leq C|\tau|^{j-1}\).
When the fiber dimension is zero, its point contribution to the
norm gives the same bound.  Dividing by the Jacobian coefficient,
which is \(\tau^{d_\pi}\) times a unit, yields
\begin{equation}\label{g:cf:evaluation-boundary}
 |\phi_v(q(z))|\leq C'|\tau(z)|^{j-1-d_\pi}.
\end{equation}
If \(P\) dominates \(U\), then \(r\) is submersive at a general
point of \(D\).  Submean gives instead
\( |\phi_v(q(z))|\leq C'|\tau(z)|^{-d_\pi}\).
In both cases the constants are uniform for unit-norm sections.

Let \(x\) define \(P\) near a general point.  On the chosen patch,
\(q^*x\) is a unit times \(\tau^l\).  Generic submersivity along
\(D\to P\), together with this transverse power, shows that the
patch maps onto a neighborhood of that point of \(P\).  We can
therefore apply the preceding estimates over each nearby \(y\),
with \(|\tau|\) comparable to \(|x(y)|^{1/l}\).  They give
\begin{equation}\label{g:cf:evaluation-exponents}
 \psi\leq O(1)+
 \begin{cases}
 \displaystyle\frac{le_P-1-d_\pi}{l}\log|x|^2,
     &P\text{ vertical over }U,\\[4pt]
 \displaystyle-\frac{d_\pi}{l}\log|x|^2,
     &P\text{ horizontal over }U.
 \end{cases}
\end{equation}
For \(P\not\subset N\), subtracting the weight of \(R_g\) leaves
a locally bounded-above function: in the vertical case the
remaining coefficient is \(1-1/l\geq0\), and in the horizontal
case it is zero.  At a component of \(N\), adding
\(c\log|x|^2\) makes the remaining coefficient nonnegative once
\[
 c\geq (1+d_\pi)/l-1\quad\text{in the vertical case},
 \qquad c\geq d_\pi/l\quad\text{in the horizontal case}.
\]
There are finitely many omitted primes, so one nonnegative integer
\(c\) works for all of them.  The adjusted weights are psh off the
divisors, since the logarithms there are pluriharmonic.  Their upper
bounds extend them across the general points of every boundary
divisor.  The remaining exceptional sets have codimension at least
two; psh removability extends the weights there as well.  The frame
changes identify the resulting metric as one on
\(-g^*K_U-R_g+cN\).  This proves pseudoeffectivity and hence
\begin{equation}\label{g:cf:evaluation-quotient-degree}
 \alpha\cdot c_1(T_Y/\mathcal F)\geq0
 \qquad(\alpha\text{ movable},\ \alpha\cdot N=0).
\end{equation}

We finish with the movable-slope argument of Ou, using the tensor
slope theorem of Greb--Kebekus--Peternell and the algebraicity
criterion of Campana--P\u{a}un
\cite{GKP2016,CampanaPaun2019,Ou2023,LMPTX2023}.
Suppose that a line in \((\Omega_Y^1)^{\otimes j}\) has positive
\(\alpha\)-degree.  For \(j=0\) this is impossible, since a line
subsheaf of \(\mathcal O_Y\) has an effective negative.  For
\(j>0\), tensor slopes and duality imply
\(\mu_{\alpha,\min}(T_Y)<0\).  Applying
\eqref{g:cf:evaluation-quotient-degree} to the identity of \(Y\)
gives \(\alpha\cdot c_1(T_Y)\geq0\).  Thus the positive-slope
part \(\mathcal F_+\) of its Harder--Narasimhan filtration is
nonzero and proper.  Its minimal slope is positive, whereas
\(T_Y/\mathcal F_+\) has nonpositive maximal slope and strictly
negative degree.

The bracket map
\(\mathcal F_+\otimes\mathcal F_+\to T_Y/\mathcal F_+\)
vanishes by the slope gap, so \(\mathcal F_+\) is a foliation.
The positive-minimal-slope criterion makes it algebraically
integrable.  Its rational quotient contradicts
\eqref{g:cf:evaluation-quotient-degree}.  These uses of slope theory
also apply to a boundary movable class.  Equivalently, perturb
toward a rational interior movable class: the strict positive
minimal slope and the bracket slope gap persist.  The tensor
inequality at the original class follows by taking limits from
the interior, using
\[
 \mu_{\alpha+\epsilon\beta,\min}(\mathcal G)
 \geq\mu_{\alpha,\min}(\mathcal G)
       +\epsilon\mu_{\beta,\min}(\mathcal G)
\]
and the corresponding upper bound for maximal slopes.  This
proves the cotangent assertion.
\end{proof}

\subsection{Restricting the contraction to a four-generator semigroup}

The evaluation slope of Proposition~\ref{g:cf:evaluation-slope} can be
used with a smaller collection of linear systems than all null-degree
systems. We record the precise collection because maximality within it
still has to control the poles appearing in the uniqueness test.

\begin{proposition}\label{g:cf:finite-semigroup}
Keep the fixed-section data of Section~\ref{g:cf:minima} and assume
$H^0(Y,\Omega_Y^q)=0$ for $q>0$. After taking a power of $s$, let $D,D^*$
be integral, put $G=D^*-D\geq0$, and let $N$ be the reduced omitted-prime
divisor. Suppose $D$ has a bounded semipositive metric and that every
cotangent line has nonpositive degree against every movable class
annihilating $N$. If there is no invariant global section of a positive
multiple of $L$, a smaller base preserving an invariant generic section
of a positive multiple of $L$ and the original-prime uniqueness condition
is obtained from one of the systems
\[
 |mD+iM+kG+lN|,\qquad m,i,k,l\in\mathbb Z_{\geq0},
\]
where $M$ is a single cotangent determinant line and
$H^0(Y,M+m_jD)\ne0$ for unbounded positive $m_j$.
\end{proposition}

\begin{proof}
The divisor $D+G+N=D^*+N$ cannot be big. A section of a positive
multiple, represented by $\gamma\in\mathbb C(Y)^*$, would satisfy
$\operatorname{ord}_P\gamma+j d_P^*\geq0$ at every nonomitted prime.
The minimum defining $d_P^*$ then makes $s^j\gamma$ regular at every
vertical prime of $X$; its horizontal regularity is already known.
This would be an invariant global section. Since $D$ is pseudoeffective
and $G,N$ are effective, cone duality supplies $\alpha\ne0$ movable with
\[
 \alpha\cdot D=\alpha\cdot G=\alpha\cdot N=0.
\]
The assumed cotangent estimate and the determinant proof in
Lemma~\ref{g:cf:effective-sequence} produce $M$ and effective
$C_j\sim M+m_jD$; effectivity forces $\alpha\cdot M=0$.
Thus every displayed system has degree zero, and none gives a generically
finite map on $Y$.

Choose one nonempty displayed system $|\Lambda|$ of maximal image
dimension, and let $Y'$ be a smooth model of its Stein field. Products
remain in the displayed collection. Exactly the product-map argument in
Lemma~\ref{g:cf:maximal-field} therefore shows that ratios in any one
of these systems lie in $\mathbb C(Y')$. Moreover $\dim Y'<\dim Y$.
The interpolation of Lemma~\ref{g:cf:horizontal-slope} uses only products
of the sections of $M+m_jD$, so it applies within this collection. For
$m_1>m_0$, the ratio of those sections gives a rational section of
$(m_1-m_0)D$ regular on the complete generic fiber over $Y'$. Write $\delta=m_1-m_0$ and let $u$ be the normalized section of
$aL-f^*D$. If $\xi$ is the generic section of $\delta D$ just obtained,
then $u^\delta f^*\xi$ is the new invariant generic section of
$a\delta L$.

We check all-degree original-prime uniqueness directly. Let $w_1,w_2$
be invariant rational sections of $bL$, $b\geq0$, regular at the primes
of $X$ horizontal over $Y'$. Old uniqueness gives
\[
 w_i^a=s^b f^*\gamma_i,\qquad \gamma_i\in\mathbb C(Y)^*.
\]
At a nonomitted prime $P$ of $Y$ horizontal over $Y'$, every
nonexceptional component above it is horizontal over $Y'$ too. Testing
regularity there and taking the nonexceptional minimum gives
\[
 \operatorname{ord}_P\gamma_i+b(d_P+\operatorname{coeff}_P G)\geq0.
\]
Hence, for one sufficiently large $l$, the $\gamma_i$ are rational
sections of $bD+bG+lN$ with poles only at primes vertical over $Y'$.

These finitely many poles can be cleared by the same section of some
$n\Lambda$; if there are no poles, use $n=0$ and the constant section. Resolve $|\Lambda|$. Each such prime maps into a proper
closed subset of its image: a valuation trivial on the image field is
trivial on its algebraic closure, so verticality persists for the image
field. Choose a projective hypersurface containing those images but not
the whole system image. Substituting the defining sections of $\Lambda$
gives a nonzero section vanishing at each pole prime; its sufficiently
high power gives the required orders. The fixed divisor of the resolved
system only adds vanishing. Multiplication by this common section places
both $\gamma_i$ in one system of the displayed semigroup. Maximality
puts their ratio in $\mathbb C(Y')$. Thus
$(w_1/w_2)^a\in\mathbb C(Y')$, and relative algebraic closedness inside
$\mathbb C(X)$ gives $w_1/w_2\in\mathbb C(Y')$.
Lemma~\ref{g:cf:uniqueness} now restores fixed-section uniqueness,
including degree zero. This proves the proposition.
\end{proof}

For the geometric data constructed from a rational quotient, the bounded
metric is supplied by Lemma~\ref{g:cf:finite-volume-base}, while
Proposition~\ref{g:cf:evaluation-slope} supplies exactly the cotangent
hypothesis. This verifies a complete contraction using the evaluation
metric, independently of the signed-volume slope argument. The four-generator
semigroup need not generate the whole null-prime field; its own
maximality and pole-clearing property are what the proof uses.

\subsection{A parameter field allowing poles at omitted primes}
\label{g:cf:parameter-subsection}

The next construction uses the signed-boundary slope inequality
of Proposition~\ref{g:bt:signed-slope}.  Its field is built directly
from rational functions with prescribed orders away from the
omitted primes.  Allowing those poles is useful because no divisor
of the original variety lies above an omitted prime.  The
effective divisors used for horizontal interpolation, however,
remain effective on the whole base.

Use fixed-section data \((f:X\dashrightarrow Y,k,s)\) satisfying
the regularity and uniqueness conditions of
Section~\ref{g:cf:conditions}, and a normal equidimensional model
\(X\xleftarrow pW\xrightarrow\varphi Y\).  For a prime
\(\Gamma\subset Y\), and primes \(A\subset W\) mapping onto it,
write
\[
 e_A=\operatorname{ord}_A(\varphi^*\Gamma),\qquad
 l_A=\operatorname{ord}_A(p^*s),\qquad
 b_\Gamma=\min_{A\to\Gamma}\frac{l_A}{e_A}.
\]
Let \(C_1,\ldots,C_v\) be the omitted primes, above which every
such \(A\) is exceptional over \(X\).  Off this finite list put
\[
 b^\sharp_\Gamma=
 \min_{\substack{A\to\Gamma\\A\text{ nonexceptional over }X}}
 \frac{l_A}{e_A},
\]
and give \(b^\sharp_\Gamma\) the value zero on the omitted primes.
Denote the resulting divisors by \(B,B^\sharp\).  Replace \(s\)
by a power, and scale \(k,B,B^\sharp\), so that both divisors are
integral.  We have
\begin{equation}\label{g:cf:parameter-model}
 D_s:=\operatorname{div}(p^*s)-\varphi^*B\geq0,
 \qquad b^\sharp_\Gamma\geq b_\Gamma
 \quad(\Gamma\ne C_i).
\end{equation}
Moreover, for \(l>0\) integral,
\begin{equation}\label{g:cf:parameter-lift}
 \operatorname{ord}_\Gamma v+l b^\sharp_\Gamma\geq0
 \quad(\Gamma\ne C_i)
 \quad\Longrightarrow\quad
 s^l f^*v\in H^0(X,klL)^T\setminus\{0\}
\end{equation}
for \(v\in\mathbb C(Y)^*\).  Indeed the inequalities test every
vertical prime of \(X\) through its nonexceptional transform on
\(W\); the original horizontal primes have nonnegative orders
by the fixed-section condition.  Normality of \(X\) then gives
regularity.

For the geometric application, assume again that $L=-K_X$ is smoothly
semipositive and $H^0(X,\Omega_X^q)=0$ for $q>0$. Pullback of forms
gives the same vanishing on $Y$, hence $\chi(\mathcal O_Y)=1$.
Prepare this graph toroidally as in Section~\ref{g:cf:preparation};
then define the minima on that prepared graph as above. The bounded metric and signed nef anticanonical class in
Lemma~\ref{g:cf:metric-conclusions}, the determinant construction in Lemma~\ref{g:determinant}, and
Proposition~\ref{g:bt:signed-slope} give an integral divisor \(R\)
and an unbounded set \(\mathcal M\subset\mathbb Z_{>0}\) with
\begin{equation}\label{g:cf:parameter-effective}
 [R]\in\overline{\operatorname{Psef}(Y)
                     +\sum_{i=1}^v\mathbb R[C_i]},
 \qquad N_m=mB-R+\operatorname{div}(q_m)\geq0
 \quad(m\in\mathcal M).
\end{equation}
Here \(q_m\in\mathbb C(Y)^*\), and the effectivity of \(N_m\)
includes the omitted primes. To make the sign of \(R\) explicit,
Lemma~\ref{g:determinant} gives a line \(M\) in a cotangent tensor
power of order \(h\) and effective divisors
\[
 N_m\sim mB+M,\qquad R=-M.
\]
The signed boundary \(B_{\mathrm{sgn}}\) furnished by
Lemma~\ref{g:cf:metric-conclusions} has negative part supported
on the \(C_i\). Proposition~\ref{g:bt:signed-slope} gives
\[
 [R]+h[B_{\mathrm{sgn}}^-]\in\operatorname{Psef}(Y),
\]
so \([R]\) lies in
\(\operatorname{Psef}(Y)+\sum_i\mathbb R[C_i]\), and in particular
in its closure as required in \eqref{g:cf:parameter-effective}.
The following argument depends only on
the explicit data \eqref{g:cf:parameter-model}--
\eqref{g:cf:parameter-effective} and fixed-section uniqueness.

\begin{proposition}\label{g:cf:parameter-contraction}
Assume fixed-section data $(f:X\dashrightarrow Y,k,s)$ as in
Section~\ref{g:cf:conditions}. Let $B,B^\sharp$ be the integral
vertical-minimum divisors of the preceding normal equidimensional graph,
and let $C_i$ be its omitted primes. Assume that an integral divisor $R$
and an unbounded sequence of effective divisors satisfy
\eqref{g:cf:parameter-effective}.
Then either a positive multiple of \(L\) has a
nonzero invariant section on \(X\), or there is a dominant rational
map \(Y\dashrightarrow Y'\), with \(\dim Y'<\dim Y\), such that
its composite with \(f\) admits fixed-section data satisfying
the regularity and uniqueness conditions of
Section~\ref{g:cf:conditions}.  Its base field is relatively
algebraically closed in \(\mathbb C(X)\).
\end{proposition}

\begin{proof}
Fix \(a\in\mathcal M\).  For \(t\in\mathbb Q\), \(t\geq a\),
and \(h\in\mathbb Z_{>0}\), define a vector space of rational
functions by
\begin{equation}\label{g:cf:parameter-spaces}
 V_{t,h}=\{0\}\cup
 \left\{v\in\mathbb C(Y)^*:
 \operatorname{div}(v)+h(tB^\sharp-R)\geq0
 \text{ away from the }C_i\right\}.
\end{equation}
Rational coefficients cause no ambiguity: this is an inequality
of rational divisors.  These spaces need not be finite-dimensional.
Their multiplication law is
\begin{equation}\label{g:cf:parameter-product}
 V_{t_1,h_1}V_{t_2,h_2}
 \subset
 V_{(h_1t_1+h_2t_2)/(h_1+h_2),\,h_1+h_2}.
\end{equation}
By \eqref{g:cf:parameter-effective} and
\(B^\sharp\geq B\) away from the omitted primes,
\(q_m\in V_{m,1}\).  Every rational \(t\geq a\) is a rational
convex combination of \(a\) and some \(m\in\mathcal M\) with
\(m\geq t\).  Products of powers of \(q_a,q_m\) therefore show
that \(V_{t,h}\ne0\) for some \(h>0\).

At a fixed \(t\), let \(F_t\) be the set of ratios \(v/w\) with
\(v,w\in V_{t,h}\), \(w\ne0\), allowing all positive \(h\).
It is a field: multiplication and cross-multiplication for addition
use \eqref{g:cf:parameter-product}, and inverses interchange
numerator and denominator.  Its relative algebraic closure in
\(\mathbb C(Y)\) is independent of \(t>a\).  To see this, take
\(a<t_1<t_2\) and a ratio at \(t_1\).  Pad equal powers of its
numerator and denominator with a common power of \(q_m\), where
\(m>t_2\).  Choose the powers so that their weighted-average
parameter is \(t_2\).  A positive power of the original ratio
then belongs to \(F_{t_2}\).  Conversely, pad a ratio at \(t_2\)
with \(q_a\) to reach \(t_1\).  Clearing the rational weights
makes all the powers integral.  Each field is therefore contained
in the relative algebraic closure of the other.

Write \(K\) for the common relatively closed field, and choose a
smooth projective model \(Y'\) with \(\mathbb C(Y')=K\).  This
field is finitely generated: choose a transcendence basis of
\(K/\mathbb C\); the algebraic closure of the resulting rational
function field inside the finitely generated field \(\mathbb C(Y)\)
is finite over it.  Since \(\mathbb C(Y)\) is relatively
algebraically closed in \(\mathbb C(X)\), the same is true of
\(K\).  We now distinguish whether this field has lost dimension.

\emph{The full-dimensional case.}
Suppose \(\operatorname{trdeg}_{\mathbb C}K=d=\dim Y\).
At one fixed rational \(t>a\), finitely many ratios in \(F_t\)
are algebraically independent of cardinality \(d\).  Use common
products and powers to place their numerator and denominator
sections in one degree.  A sufficiently large rational divisor
\(C\) supported on the \(C_i\) compensates their finitely many
omitted-prime poles.  After clearing denominators, these sections
belong to a multiple of \(tB^\sharp-R+C\) and give a rational map
with \(d\)-dimensional image.  Hence
\begin{equation}\label{g:cf:parameter-big}
 tB^\sharp-R+C\quad\text{is big}.
\end{equation}

The closure in the first condition of
\eqref{g:cf:parameter-effective} is sufficient to remove \(R\)
from this assertion.  Choose pseudoeffective classes \(P_j\) and
real divisors \(E_j\) supported on the \(C_i\) such that
\(P_j+[E_j]\to[R]\).  Then
\[
 [tB^\sharp+C-E_j]-P_j
 \longrightarrow [tB^\sharp-R+C].
\]
The classes on the left are big for large \(j\), by openness of
the big cone.  Adding \(P_j\) preserves bigness, so
\(tB^\sharp+C-E_j\) is big.  Rational approximation of the
coefficients on the finite support yields a rational divisor
\(C'\), supported on that same list, with \(tB^\sharp+C'\) big.
Consequently, for some positive integer \(h\) clearing all
denominators and some \(v\in\mathbb C(Y)^*\),
\[
 \operatorname{div}(v)+h(tB^\sharp+C')\geq0.
\]
Away from the omitted primes this is precisely the hypothesis of
\eqref{g:cf:parameter-lift}, with the positive integer \(l=ht\).
It produces the asserted global invariant section.

\emph{The smaller-base case.}
Suppose now that \(\dim Y'<d\).  Choose \(b\in\mathcal M\) with
\(b>a\), write \(\delta=b-a\), and set
\begin{equation}\label{g:cf:parameter-new-section}
 f':X\dashrightarrow Y',\qquad
 k'=k\delta,\qquad
 s'=s^\delta f^*(q_b/q_a),
\end{equation}
where \(f'\) is the composite rational map.  Both the map and
section are invariant.  We verify separately that the section is
regular over the generic point and that it restores uniqueness.

For every \(m\in\mathcal M\) with \(m>b\), the numerator and
denominator in
\begin{equation}\label{g:cf:parameter-interpolation}
 \frac{q_m^{\delta}q_a^{m-b}}{q_b^{m-a}}
\end{equation}
belong to \(V_{b,m-a}\), since
\(\delta m+(m-b)a=b(m-a)\).  Their ratio therefore lies in
\(K\).  On a common smooth resolution mapping to \(W,Y,Y'\),
let \(n_m\) be the coefficient of the pullback of \(N_m\) at a
prime horizontal over \(Y'\).  Its valuation is trivial on
\(K\), so \eqref{g:cf:parameter-interpolation} implies
\[
 \delta(n_m-n_a)=(m-a)(n_b-n_a).
\]
Every \(n_m\geq0\), and the indices \(m\) are unbounded.
It follows that \(n_b-n_a\geq0\).  The divisor of the pullback
of \(s'\) is
\[
 \delta D_s+\varphi^*(N_b-N_a),
\]
with pullbacks to the common resolution understood.  Its
horizontal coefficients are nonnegative by
\eqref{g:cf:parameter-model} and the preceding inequality.
Thus \(s'\) has no poles over the generic point of \(Y'\),
including at exceptional horizontal primes.

To prove uniqueness, take \(r\geq0\) and a nonzero invariant
rational section \(\tau\) of \(rk'L\) having no poles on the
prime divisors of \(X\) horizontal over \(Y'\).  These include
the primes horizontal over \(Y\).  The old fixed-section
condition writes
\[
 \tau=s^{\delta r}f^*u,\qquad
 (s')^r=s^{\delta r}f^*u_s,
 \qquad u_s=(q_b/q_a)^r,
\]
with \(u,u_s\in\mathbb C(Y)^*\).  At every non-omitted prime
\(\Gamma\) whose valuation is trivial on \(K\), the
nonexceptional primes above it are horizontal over \(Y'\).
Regularity of both sections there gives
\begin{equation}\label{g:cf:parameter-horizontal-tests}
 \operatorname{ord}_\Gamma(u)+\delta r b^\sharp_\Gamma\geq0,
 \qquad
 \operatorname{ord}_\Gamma(u_s)+\delta r b^\sharp_\Gamma\geq0.
\end{equation}
Only finitely many non-omitted primes fail either inequality.
Every such prime has valuation nontrivial on \(K\).

We can compensate all these deficits with a common element of
one \(V_{b,h}\).  For a non-omitted prime \(\Gamma\), define
the excess of \(v\in V_{b,h}\setminus\{0\}\) by
\[
 \epsilon_\Gamma(v)=
 \operatorname{ord}_\Gamma(v)
   +h(bb^\sharp_\Gamma-\operatorname{coeff}_\Gamma R)
 \geq0.
\]
If \(\operatorname{ord}_\Gamma\) is nontrivial on \(K\), it is
nontrivial on \(F_b\): a valuation trivial on a field remains
trivial on an algebraic extension.  Some ratio \(v_1/v_2\), with
\(v_1,v_2\in V_{b,h_\Gamma}\), therefore has nonzero order.
The two nonnegative excesses are different, so at least one is
strictly positive.  Choose that element.  Products of sufficiently
large powers of these chosen elements, one for each deficit
prime, give \(v\in V_{b,h}\setminus\{0\}\) whose excess
compensates every deficit of both \(u\) and \(u_s\).
At all other non-omitted primes, the excesses remain nonnegative.
If there are no deficits, any nonzero element of \(V_{b,h}\)
suffices.  Thus
\[
 uv,\ u_s v\in V_{\,b+\delta r/h,\,h}.
\]
The parameter is still greater than \(a\), also when \(r=0\).
It follows that
\[
 \tau/(s')^r=u/u_s\in F_{b+\delta r/h}\subset K.
\]
This is the required fixed-section uniqueness for the smaller
base.

These verifications also retain the adjoint rank conditions:
on a smooth resolution, multiplying powers of \(s'\) by the
Jacobian gives a nonzero invariant adjoint section in every
nonnegative multiple of $k'$.  Dividing any other such section by the
Jacobian gives a rational section regular at the original
horizontal primes; the uniqueness just proved makes it a
multiple over \(K\) of that generator.  Thus the contraction
returns precisely the section and rank data needed for the next
step, including degree zero.
\end{proof}

\section{Contractions with different boundary data}
\label{g:ba}

Theorem~\ref{g:cf:contraction-theorem} keeps a generic section and its
adjoint rank. The boundary constructions of Section~\ref{g:bt} also accommodate
three other useful forms of control: arbitrary exceptional twists,
eventual discrepancy inequalities, and pole tests on the original
variety.  We verify these forms separately.  This identifies the exact
support condition that must survive each dimension decrease.

Throughout this section $X$ is smooth connected projective,
$H^0(X,\Omega_X^q)=0$ for $q>0$, and $L=-K_X$ has a smooth semipositive
metric.  A torus $T$ acts with its natural linearization on $L$.
Every smooth projective rational image $Y$ has no positive-degree
holomorphic forms, so $\chi(\mathcal O_Y)=1$ and
$\operatorname{Pic}(Y)\otimes\mathbb Q$ is finite dimensional.
We use geometrically integral generic fibres throughout.

Two of the pole formulations below express the same uniqueness
condition. Given the distinguished generic section, the bound for
every exceptional twist in \eqref{g:bt:all-exceptional-rank} is equivalent
to original-prime uniqueness in Lemma~\ref{g:cf:uniqueness}.
Indeed, finitely many invariant rational sections regular at the
original horizontal primes acquire only exceptional horizontal poles
on a fixed equivariant resolution. One common effective invariant
exceptional divisor clears those poles. Conversely an exceptional
twist changes no order at an original prime. The constructions below
remain different: they use, respectively, a weighted metric and Iitaka
rigidity, or a signed boundary and the null-prime field.

\subsection{Rigid Iitaka fibres and arbitrary exceptional twists}
\label{g:ba:iitaka}

The following elementary form of Iitaka rigidity is useful because the
base may need further resolution after its Iitaka map is chosen.

\begin{lemma}\label{g:ba:rigidity}
Let $Q\geq0$ be an integral divisor on a smooth projective variety $Y$,
and let $u:Y'\to Y$, $h:Y'\to Y_2$ resolve its Iitaka fibration, with
connected generic fibre $S$.  For every effective $u$-exceptional
integral divisor $A$ and every positive integer $m$,
\[
             h^0\bigl(S,m(u^*Q+A)|_S\bigr)=1.
\]
\end{lemma}

\begin{proof}
Choose the Iitaka map from a multiple of $Q$ attaining the maximal image
dimension and take its Stein factorization.  A multiple of $u^*Q$
dominates the pullback of a big divisor on $Y_2$; after taking a higher
multiple and replacing that divisor by an ample divisor plus an
effective divisor, it dominates a pulled-back ample divisor.
The generic fibre is geometrically integral.
Adding an effective exceptional divisor changes no global space of
sections, by normality, and hence no Iitaka dimension.
If a multiple of $(u^*Q+A)|_S$ had two independent sections, extend
them rationally over $Y'$ and clear their vertical poles with a
sufficiently high pulled-back ample twist.  The preceding domination
absorbs that twist into a multiple of $u^*Q+A$.
Their ratio remains nonconstant on $S$.  Combined with ratios defining
the base, it increases the transcendence degree of a global section-ratio
field beyond $\kappa(Q)$, a contradiction.
The canonical effective divisor supplies the nonzero section.
\end{proof}

\begin{proposition}[Preserving arbitrary exceptional twists]
\label{g:ba:exceptional-contraction}
Let $\pi:Z\to X$, $g:Z\to Y$ be a smooth equivariant resolution
of a dominant $T$-invariant rational map to a smooth projective $Y$,
with $\dim Y>0$ and geometrically integral generic fibre
$F$ over $K=\mathbb C(Y)$. Suppose that
$H^0(F,a\pi^*L|_F)^T\ne0$ for some $a>0$, and that
\[
 \dim_K H^0(F,(b\pi^*L+A)|_F)^T\leq1
\]
for every integer $b\geq0$ and every effective invariant
$\pi$-exceptional integral divisor $A$.
Then either a positive multiple of $L$ has a nonzero invariant section
on $X$, or a dominant rational map $Y\dashrightarrow Y_2$ to a smooth
projective base of smaller dimension preserves both the generic section
and this bound for every exceptional twist on a smooth equivariant
resolution of the composite map.
\end{proposition}

\begin{proof}
Use the divisors $N,R,B$ of Proposition~\ref{g:bt:weighted-transfer}.
If any $j(N+R)+lB$, $j>0$, $l\geq0$, is rationally effective, that
proposition gives the section on $X$.  Assume otherwise.
The weighted cotangent bound gives $hcB-c_1(M)$ pseudoeffective for a
line $M$ in a cotangent tensor power of order $h$.  Consequently
\[
                    M+j(N+R)+B
                    \quad\hbox{is not big}\qquad(j>0);
\]
adding the pseudoeffective class $hcB-M$ to a big class would contradict
the preceding noneffectivity.

The bounded metric of $N$, hard Lefschetz and the nonzero constant term
of $\chi(K_Y+mN)$ give, by the determinant construction
of Lemma~\ref{g:determinant}, a fixed cotangent line $M$ and
effective $D_j\sim M+jN$ for unbounded $j>0$.
Fix two indices $i<b$ and put $Q=D_b+bR+B$.  This divisor is effective
and not big.  Resolve its Iitaka map by $u:Y'\to Y$, $h:Y'\to Y_2$.
Then $\dim Y_2<\dim Y$.

Restrict pullbacks to the generic fibre $S$ of $h$.
The interpolation equivalence \eqref{g:cf:interpolation}, applied at
the twist exponents $i<b<j$, gives two effective representatives of
the class $(j-i)D_{b,S}$.  Add the same effective complement
$(j-i)(Q_S-D_{b,S})$ to both.  Lemma~\ref{g:ba:rigidity} makes the
resulting representatives of $(j-i)Q_S$ equal, and cancellation gives
\[
           (j-i)D_{b,S}=(b-i)D_{j,S}+(j-b)D_{i,S}
                                                  \qquad(j>b).
\]
In particular all the later $D_{j,S}$ are supported in $D_{b,S}$.
Since $j$ is unbounded and $D_{j,S}\geq0$, the divisor
\[
                   V_S=\frac{D_{b,S}-D_{i,S}}{b-i}
\]
is effective and rationally equivalent to $N_S$.
Its support lies in $D_{b,S}$.  For every effective $u$-exceptional $A$
and integers $n,l\geq0$ this also gives
\begin{equation}\label{g:ba:exceptional-fibre-bound}
 h^0\bigl(S,n(N+R)_S+lB_S+A|_S\bigr)\leq1.
\end{equation}
To see this, represent $N_S$ by $V_S$.  The indicated effective
representative is supported in $(u^*Q+A)|_S$ and therefore dominated
by a multiple of it.  Two independent sections would remain independent
after taking powers and multiplying by a fixed section, contrary to
Lemma~\ref{g:ba:rigidity}.

On a common resolution over $Y'$, lift a denominator-clearing multiple
of $V_S$ through $f^*N\to ap^*L$.  It gives the new invariant generic
section. For the rank assertion, let $F_{\mathrm{exc}}$ be any effective
invariant divisor exceptional over $X$ on the common resolution.
A section with this twist over the new generic fibre is, by the old
rank condition, rationally
$\lambda s^n s_{F_{\mathrm{exc}}}$ with $\lambda\in\mathbb C(Y)$, first for exponents
divisible by $a$.  At a prime of $S$ not exceptional over $Y$ and not
in $B_S$, choose a nonexceptional component realizing the restricted
minimum $n'_Q=(N+R)_Q$.  The exceptional twist has coefficient zero
there.  Regularity gives
\[
                       \operatorname{ord}_Q\lambda
                                  \geq-n(N+R)_Q.
\]
All remaining poles occur on $B_S$ or on exceptional divisors of $u$.
For any proposed pair of sections, choose an effective $u$-exceptional
divisor $A_Y$ on $Y'$ and $l\geq0$ so that
$lB_S+A_Y|_S$ absorbs them. Equation~\eqref{g:ba:exceptional-fibre-bound}
then forces dependence.  Other exponents follow by taking powers to
clear divisibility; independent sections cannot become dependent under
a power on a geometrically integral fibre.
\end{proof}

\begin{proposition}[The unweighted rigid-fibre variant]
\label{g:ba:toroidal-contraction}
Let $\pi:Z\to X$, $g:Z\to Y$ be a smooth equivariant resolution
of a dominant $T$-invariant rational map to a smooth projective $Y$,
with $\dim Y>0$ and geometrically integral generic fibre
$F$ over $K=\mathbb C(Y)$. Suppose there is a nonzero invariant section
of $a\pi^*L|_F$, $a>0$, and
\[
 \dim_K H^0(F,(K_{Z/X}+ma\pi^*L)|_F)^T=1
                    \qquad(m\geq0).
\]
Then either a positive multiple of $L$ has a nonzero invariant section
on $X$, or a dominant rational map $Y\dashrightarrow Y_2$ to a smooth
projective base of smaller dimension preserves the generic section and
these Jacobian-twisted ranks, with a possibly different positive
integer $a$.
\end{proposition}

\begin{proof}
Take the prepared toroidal model of
Proposition~\ref{g:bt:toroidal-transfer}.
Let $B$ now be the reduced sum of primes $Q$ for which no component
attaining the minimum $n_Q$ is nonexceptional over $X$.
For small rational $\delta>0$, the normalized divisor $D_s$ absorbs
$\delta f^*B$ at every nonexceptional prime.
Thus effectivity of $N+\delta B$ supplies a section on $X$.
Otherwise this pseudoeffective divisor is not big, so a nonzero movable
class $\alpha$ annihilates both $N$ and $B$.
The correction $\Theta$ of Proposition~\ref{g:bt:toroidal-transfer} is
supported on $B$.  Proposition~\ref{g:bt:unweighted-slope} gives
$\alpha\cdot M\leq0$ for every cotangent line $M$.

The determinant argument produces effective
$D_j\sim M+jN+B$, $D_j\geq B$, for unbounded $j$; their degree against
$\alpha$ is zero.  Resolve the Iitaka map of one $D_b$ after fixing
an earlier index $i<b$.  The same interpolation identity as above puts
every later restricted divisor in the support of $D_{b,S}$.
Include the finitely many exceptional divisors of the chosen resolution
in an auxiliary effective divisor $A$ and apply
Lemma~\ref{g:ba:rigidity} to $u^*D_b+A$.
Choose comparable indices $j_1<j_2$ in the tail beyond $b$, using
the ascending-pair property in a finite product of $\mathbb N$.
Put $v=j_2-j_1>0$. Then
\[
 C=(D_{j_2}-D_{j_1})|_S\geq0,\qquad C\sim vN_S,
 \qquad \operatorname{Supp}C\subseteq\operatorname{Supp}D_{b,S}.
\]
Let $\zeta$ be the corresponding section of $vN_S$, and let
$\sigma$ denote the normalized regular section of
$a\pi^*L-g^*N$ on a common resolution. The section
$\sigma^v g^*\zeta$ gives a new invariant generic section $s'$
of degree $avL$.

We verify the new ranks at exactly these multiples. Let $q\geq0$
and let $t$ be an invariant section of $K_{Z'/X}+qav\pi'^*L$
on the new generic fibre, where $Z'$ is a common smooth resolution.
The old rank condition, applied over $K$, says that division by
the Jacobian section and $\sigma^{qv}$ leaves the pullback of a
rational section $r$ of $qvN_S$. At a prime of $S$ outside $B_S$
and the $u$-exceptional locus, a nonexceptional minimizing component
has both normalized section order and Jacobian order zero.
Regularity of $t$ therefore makes $r$ regular at that prime.

The rational function $r/\zeta^q$ has poles only in
$\operatorname{Supp}(u^*D_b+A)|_S$: the divisor $D_b$ contains $B$,
the zeros of $\zeta$ lie in $D_{b,S}$, and $A$ contains every
$u$-exceptional prime of the fixed resolution.
Every multiple of this effective divisor is rigid by
Lemma~\ref{g:ba:rigidity}, so the function is constant over
$\mathbb C(Y_2)$. Thus $t$ is a base multiple of the Jacobian
times $(s')^q$. This also treats $q=0$, when $\zeta^q=1$.
Those Jacobian products supply nonzero sections in every tested
degree, proving rank exactly one. The argument permits extra
resolutions without assuming that their exceptional centres lie
in the original base locus.
\end{proof}

\subsection{Contraction using a large adjoint moment}
\label{g:ba:eventual}

\begin{proposition}\label{g:ba:eventual-contraction}
Suppose $K,s,d$ satisfy
\eqref{g:bt:horizontal-regular}--\eqref{g:bt:eventual-rank}.
If $\operatorname{trdeg}_{\mathbb C}K>0$, either a positive multiple of
$L$ has a nonzero invariant section, or these same two conditions hold
for some $K'\subset K$ with strictly smaller transcendence degree and a
new invariant rational section of a positive multiple of $L$.
\end{proposition}

\begin{proof}
Use Proposition~\ref{g:bt:large-moment} and put $N^+=N+R/k$.
If $N^+$ is rationally effective, its section returns to $X$.
Otherwise $N^+$ is pseudoeffective and not big.  Choose a nonzero
movable $\alpha$ annihilating it.  Since $N,R$ are pseudoeffective,
\[
                         \alpha\cdot N=\alpha\cdot R=0.
\]
The unweighted metric of $J=-K_Y+kN+R$ and
Proposition~\ref{g:bt:unweighted-slope} give the cotangent bound for
$\alpha$.  The determinant construction supplies effective
$D_j\sim M+m_jN$ with $m_j$ increasing and $\alpha\cdot D_j=0$.

Let $K'$ be the field of functions with poles only on $\alpha$-null
primes.  Lemma~\ref{g:cf:null-field} gives finite generation, relative
algebraic closure, and strict transcendence-degree decrease.
On a resolution of the resulting map there are only finitely many
possible horizontal components of the $D_j$: the null horizontal primes
are finite because $\operatorname{Pic}(Y)\otimes\mathbb Q$ is finite
dimensional, and the resolution has only finitely many exceptional
primes.  The finite-support argument in the proof of Proposition~\ref{g:cf:dickson}
therefore gives $j>i$, $v=m_j-m_i>0$ and $c\in K^*$ with
\[
             D_j-D_i=vN+\operatorname{div}(c)
\]
horizontally effective over $K'$.
Set $s'=s^vc$.  Its divisor upstairs is the sum of the pullback of this
difference and $vD\geq0$.  Hence it is regular at every valuation
trivial on $K'$, proving \eqref{g:bt:horizontal-regular}.

For all sufficiently large $m$, a section $t$ satisfying the new
discrepancy inequalities satisfies the old ones, so
$t=s^{mv}u$ with $u\in K$.  If $Q$ has positive $\alpha$-degree,
$r_Q=0$.  A component attaining the two minima in the definition of
$R$ consequently has $a_E=0$ and $h_E=e_En_Q$.
Its valuation is trivial on $K'$, and the inequality for $t$ gives
\[
                       \operatorname{ord}_Q u+mv n_Q\geq0.
\]
The null support of $D_j-D_i$ gives
$\operatorname{ord}_Q c+v n_Q=0$.
Thus $u/c^m$ has no pole on any positive-degree prime and belongs to
$K'$.  We obtain $t\in K'(s')^m$.
The opposite inclusion follows from the regularity of $s'$ and the
nonnegative discrepancies over smooth $X$.
\end{proof}

\subsection{Original horizontal primes and the signed boundary}
\label{g:ba:signed}

The next variant tests poles only on primes of $X$ itself.
It therefore permits arbitrary poles at exceptional valuations in the
rank test, while the distinguished section remains regular on the
complete generic fibre.

\begin{proposition}\label{g:ba:original-primes}
Let $K=\mathbb C(Y)\subseteq\mathbb C(X)^T$ be relatively algebraically
closed, and let $s$ be an invariant rational section of $aL$, $a>0$, regular on
a resolved generic fibre over $K$.  Suppose that, for every $j\geq0$,
\begin{equation}\label{g:ba:original-rank}
 u\in\mathbb C(X)^T,\quad
 s^ju\text{ has no poles on primes of $X$ horizontal over $K$}
                  \quad\Longrightarrow\quad u\in K.
\end{equation}
Then either a positive multiple of $L$ has a nonzero invariant section,
or a contraction to a smaller base preserves these two conditions.
\end{proposition}

\begin{proof}
On a prepared equidimensional toroidal model $p:W\to X$, $f:W\to Y$,
normalize $s$ by vertical minima:
\[
 D_s=\operatorname{div}(s)-f^*N\geq0.
\]
The hypothesis is the fixed-section uniqueness condition of
Lemma~\ref{g:cf:uniqueness}. It gives invariant adjoint rank one
in every nonnegative multiple of $a$: divide by the Jacobian and
test the nonexceptional horizontal primes. Hence
Lemma~\ref{g:cf:metric-conclusions} supplies a bounded semipositive
metric on $N$ and a signed boundary $B_{\mathrm{sgn}}$ with simple
normal crossing support and coefficients less than one, such that
$-(K_Y+B_{\mathrm{sgn}})$ is nef.

We compare the two choices of permissible base poles before choosing
the contraction. Let $S_0$ be the reduced sum of omitted primes,
above which every component is exceptional over $X$. Define $N'$
by taking the minimum only over nonexceptional components off $S_0$,
and retaining the coefficient of $N$ on $S_0$. Put $G=N'-N\geq0$.
Let $B_{\min}$ be the reduced sum of primes at which no minimizing
component is nonexceptional. Then
\[
 \operatorname{Supp}B_{\min}
   =\operatorname{Supp}S_0\ \sqcup\ \operatorname{Supp}G,
 \qquad
 \operatorname{Supp}B_{\mathrm{sgn}}^-
   \subseteq\operatorname{Supp}S_0.
\]
Indeed, away from the omitted primes, failure to attain the full
minimum on a nonexceptional component is exactly the strict inequality
$n'_Q>n_Q$.

There are two effectivity tests. If $N'+S_0$ is rationally effective,
its section returns to $X$ by the nonexceptional minimum. Alternatively,
for sufficiently small rational $\epsilon>0$,
$D_s-\epsilon f^*B_{\min}$ has no negative coefficient at any
nonexceptional prime; effectivity of $N+\epsilon B_{\min}$ therefore
also returns a section. If neither test succeeds, the chosen
pseudoeffective class is not big. Cone duality gives a nonzero movable
class $\alpha$. Either choice yields the same vanishing conditions:
\[
 \alpha\cdot N=\alpha\cdot G=\alpha\cdot S_0=0,
 \quad\text{equivalently}\quad
 \alpha\cdot N=\alpha\cdot B_{\min}=0.
\]
In particular $\alpha\cdot B_{\mathrm{sgn}}^-=0$, so
Proposition~\ref{g:bt:signed-slope} supplies the cotangent bound.
The determinant construction gives effective
$D_j\sim M+jN$ at unbounded divisible indices, with
$\alpha\cdot D_j=0$.

Let $K'$ be the null-prime field of Lemma~\ref{g:cf:null-field};
it has smaller transcendence degree. On a common resolution, the
affine interpolation of Lemma~\ref{g:cf:horizontal-slope} gives
$j>i$ and
\[
 D_j-D_i=vN+\operatorname{div}(c),\qquad v=j-i>0,\quad c\in K^*,
\]
with effective horizontal part over $K'$. One may instead use
the finite-support extraction of Proposition~\ref{g:cf:dickson}
to choose a comparable pair of horizontal multiplicity vectors.
Both choices give the same required property. Since the divisor
of $s'=s^vc$ is $vD_s+f^*(D_j-D_i)$, this section is regular on
the new complete generic fibre, including at exceptional horizontal
primes.

For uniqueness, suppose that $t=(s')^kg$ has no poles at original
primes horizontal over $K'$, where $k\geq0$ and
$g\in\mathbb C(X)^T$. Since these include the $K$-horizontal primes,
the old test applied to $t=s^{vk}c^kg$ gives $c^kg\in K$ and
hence $g\in K$. If a base prime $Q$ has positive $\alpha$-degree,
then $Q$ is not omitted, $n'_Q=n_Q$, and its valuation is trivial
on $K'$. Its nonexceptional components are therefore original
primes horizontal over $K'$. Testing $t$ on all those components
gives
\[
 \operatorname{ord}_Q g+
       k(vn'_Q+\operatorname{ord}_Q c)\geq0.
\]
The null support of $D_j-D_i$ gives
$vn_Q+\operatorname{ord}_Q c=0$. Thus $g$ has no pole at any
positive-degree prime. Its pole divisor is null, and the field
criterion in Lemma~\ref{g:cf:null-field} gives $g\in K'$.
This includes $k=0$ and proves the required uniqueness.

This proof uses the whole bad-minimum set through $G$ and $S_0$.
Annihilating only $B_{\mathrm{sgn}}^-$ would not provide the
nonexceptional minimizing components needed in the last pole test.
\end{proof}

\subsection{A toric construction of the initial generic section}
\label{g:ba:toric-start}

The rational torus quotient also admits a geometric construction of its
first invariant section, separate from wedging fundamental vector
fields.  This explains why a generated line on a generic orbit
compactification suffices, without assuming semiampleness on $X$.

\begin{proposition}\label{g:ba:toric-section}
For the rational quotient $K=\mathbb C(X)^T$, a positive multiple of
$L$ has a nonzero invariant section on the smooth resolved generic
fibre.  The data satisfy \eqref{g:ba:original-rank}.
\end{proposition}

\begin{proof}
Smooth semipositivity implies nefness, so
Lemma~\ref{g:cf:toric-initial-section} supplies the invariant section.
Ratios of invariant rational tensors of the same power of \(L\) lie in
\(\mathbb C(X)^T=K\), which proves \eqref{g:ba:original-rank}.
\end{proof}

The preceding contractions preserve their stated data and strictly
decrease a nonnegative base dimension.  At a point base the generic
section is global and pushes to $X$.  They therefore give complete
alternative reductions to an invariant section under their respective
initial conditions.

\section{Contraction by multiplicative cotangent systems}
\label{g:multiplicative}

The contraction in Section~\ref{g:cf:contraction} chooses systems of
zero degree for one supporting movable class. There is another useful
choice: close all cotangent-twisted systems under multiplication.
This avoids fixing that class, while retaining the same decisive
verification of generic adjoint rank after contraction. We state the
construction with its prepared data so that its difference from the
zero-degree method is explicit.

\begin{proposition}\label{g:multiplicative-contraction}
Let $X$ be smooth projective and rationally connected, let
$L=-K_X$ be smoothly semipositive, and let an algebraic torus $T$ act
with the natural linearization. Let
$\pi:Z\to X$, $g:Z\to Y$ resolve a dominant invariant rational map,
with $Y$ smooth projective of positive dimension and
$\mathbb C(Y)$ relatively algebraically closed in $\mathbb C(X)$.
Write $E=K_{Z/X}$. Suppose a nonzero invariant rational section $s$ of
$\ell\pi^*L$, $\ell>0$, and rational divisors $N$ on $Y$ and $D$ on
$Z$ satisfy
\[
 \ell^{-1}\operatorname{div}_Z(s)=g^*N+D,\qquad D\geq0.
\]
Assume that $N$ has a semipositive metric with locally bounded weights.
Suppose there is a rational simple normal crossings divisor
$\Delta=\Delta^+-\Delta^-$ on $Y$, with every coefficient of
$\Delta^+$ less than one, such that $-K_Y-\Delta$ is nef.
Let $B$ be a reduced effective divisor on $Y$ such that
\begin{enumerate}[label=(\roman*)]
\item $\operatorname{Supp}\Delta^-\subset\operatorname{Supp}B$;
\item $\pi_*(D-tg^*B)\geq0$ for every sufficiently small rational $t>0$;
\item above every base prime outside $B$, some prime dominating it has
coefficient zero in both $D$ and $E$.
\end{enumerate}
Finally suppose the invariant generic spaces of
$E+k\pi^*L$ have rank one for $k=0$ and every positive multiple of
one fixed integer.
Then either $H^0(X,rL)^T\ne0$ for some $r>0$, or a rational map
$Y\dashrightarrow Y'$ with $\dim Y'<\dim Y$ gives a composite
fibration retaining a nonzero invariant generic section in a positive
multiple of $L$ and invariant adjoint rank one at zero and in every
sufficiently divisible positive degree.
\end{proposition}

These prepared data are supplied by vertical minima and the toroidal
logarithmic density calculation of
Lemma~\ref{g:cf:metric-conclusions}; no strict-curvature patch is
required. The signed slope statement needed in the proof is
Proposition~\ref{g:bt:signed-slope}, proved for these signed boundaries.
In particular the argument does not apply an effective-boundary
positivity theorem directly to $\Delta$.

\begin{proof}
Suppose first that $N+tB$ is big for some sufficiently small rational
$t>0$. Choose a nonzero section of a divisible multiple. Pull it back,
use the corresponding power of $s$, and push its divisor to $X$.
Condition (ii) leaves an effective divisor. All rational coefficients
can be cleared by the power, so this is a regular nonzero invariant
section of a positive multiple of $L$. We may therefore assume that
$N+tB$ is not big for all small positive rational $t$.

For every line bundle $M$ admitting a nonzero map
$M\to(\Omega_Y^1)^{\otimes p}$,
Proposition~\ref{g:bt:signed-slope} gives
\begin{equation}\label{g:multiplicative-cotangent}
 p\Delta^- -M\quad\text{pseudoeffective}.
\end{equation}
The case $p=0$ follows directly from $M\to\mathcal O_Y$.
Choose an integer $b>0$ clearing all divisor denominators, divisible by
$\ell$, and compatible with the generic-rank hypothesis. If $s_{bN}$
denotes the canonical rational section of $bN$, the normalized section
\[
 \sigma=\frac{s^{b/\ell}}{g^*s_{bN}}
 \quad\text{of }b\pi^*L-g^*(bN)
 \quad\text{has divisor }bD\geq0.
\]
It is regular and invariant. Consider the
line bundles
\begin{equation}\label{g:multiplicative-collection}
 Q=kN+aB+M,
 \quad k\in b\mathbb Z_{\geq0},\quad a\in\mathbb Z_{\geq0},
 \quad M\longrightarrow(\Omega_Y^1)^{\otimes p},\quad p\geq0,
\end{equation}
where the last map is nonzero. This collection is closed under tensor
product: tensoring two nonzero line maps gives a nonzero map at the
generic point and hence a nonzero global map into the corresponding
tensor power.

No member of \eqref{g:multiplicative-collection} is big. Indeed,
\eqref{g:multiplicative-cotangent} would make
$kN+aB+p\Delta^-$ big if $Q$ were big. Choose $a'\geq a$ with
$a'B-aB-p\Delta^-\geq0$. Then $kN+a'B$ would be big. For any prescribed
small $t>0$, choose $\delta>0$ so that $\delta k<1$ and
$\delta a'\leq t$. Since $N$ and $B$ are pseudoeffective,
\[
 N+tB=\delta(kN+a'B)+(1-\delta k)N+(t-\delta a')B
\]
would be big, contrary to the assumption.

Choose a nonempty system in a member $Q_0$ of
\eqref{g:multiplicative-collection} whose image has maximal dimension.
The constant system shows that such systems exist. Resolve its map and
pass to the connected-fiber base to obtain
\[
 \eta:\widetilde Y\to Y,\qquad q:\widetilde Y\to Y'.
\]
We can choose $Y'$ normal projective; its smooth birational model may
be used later. Its dimension is strictly smaller than $\dim Y$, since
$Q_0$ is not big. A suitable power of the chosen system supplies an
ample line $A$ on the normal Stein base with
$\eta^*Q_0-q^*A$ effective, after replacing $Q_0$ by that power.
Every ratio from every system in
\eqref{g:multiplicative-collection} belongs to $\mathbb C(Y')$.
Otherwise the product with the chosen system would increase the image
dimension; relative algebraic closure removes the remaining algebraic
extension.

For the generic fiber $F$ of $q$ this implies
\begin{equation}\label{g:multiplicative-rigidity}
 h^0(F,\eta^*Q|_F)\leq1
 \quad\text{for every }Q\text{ in }
 \eqref{g:multiplicative-collection}.
\end{equation}
To check the assertion for generic sections, spread two putative
independent sections after a sufficiently ample base twist $q^*(nA)$.
The effective difference $\eta^*(nQ_0)-q^*(nA)$ puts them in
$\eta^*(Q+nQ_0)$. Their pushdowns are sections on the smooth $Y$,
and their ratio contradicts the preceding field statement.

Because $Y$ is a rational image of the rationally connected $X$, it
has no positive-degree holomorphic forms. Apply
Lemma~\ref{g:determinant} to $bN$. We obtain one cotangent line $M_0$
and effective divisors $D_i\sim M_0+m_iN$ for an unbounded increasing
sequence $m_i\in b\mathbb Z_{>0}$. Their restrictions to $F$ are
nonzero section divisors. By \eqref{g:multiplicative-rigidity}, the
products realizing any three-index linear relation are proportional
on $F$. Therefore each horizontal multiplicity is affine in $m_i$.
All these multiplicities are nonnegative at unbounded exponents, so
the slope is nonnegative. A difference $D_j-D_i$, $j>i$, consequently
gives a nonzero section of $(m_j-m_i)\eta^*N$ on $F$.

Write $\delta=m_j-m_i$, and choose $\gamma\in\mathbb C(Y)^*$ with
\[
 D_j-D_i=\delta N+\operatorname{div}_Y(\gamma).
\]
The integer $\delta$ is divisible by $b$ and hence by $\ell$.
Resolve the composite map from $X$ to $Y'$ by a smooth model $Z'$
dominating $Z$ and $\widetilde Y$; write $j:Z'\to Y'$ and continue
to write $g:Z'\to Y$, $\pi:Z'\to X$, with $D$ pulled back from $Z$.
The rational section
\[
 \tau=s^{\delta/\ell}g^*\gamma\quad\text{of }\delta\pi^*L
 \quad\text{satisfies}\quad
 \operatorname{div}(\tau)=\delta D+g^*(D_j-D_i).
\]
The first summand is effective, and the second is effective on the
complete generic fiber over $Y'$, as just proved. Thus $\tau$ is a
nonzero invariant generic section. The composite generic fiber is geometrically
connected, since both successive function-field inclusions are
relatively algebraically closed in characteristic zero.

It remains to restore the adjoint ranks on this larger fiber.
Suppose two invariant sections of $K_{Z'/X}+k\pi^*L$ are independent
over $\mathbb C(Y')$, for a sufficiently divisible $k\geq0$.
Clear their poles over $Y'$ by $j^*(nA)$, and pull back the effective
difference $\eta^*(nQ_0)-q^*(nA)$ to regard them as sections of
\[
 K_{Z'/X}+k\pi^*L+g^*(nQ_0).
\]
Over $\mathbb C(Y)$ the old rank is one. Divide the two sections by
the Jacobian section and $\sigma^{k/b}$, using the pullback of the
normalized section defined above. The quotients descend to rational
sections $v_1,v_2$ of $kN+nQ_0$ on $Y$.
For a prime $P$ outside $B$, condition (iii) supplies a prime $R$
above it whose strict transform has order zero in both $D$ and
$K_{Z'/X}$. If $e_R=\operatorname{ord}_R g^*P$, regularity of the
two sections upstairs gives
\[
 0\leq\operatorname{ord}_R(s_{K_{Z'/X}}\sigma^{k/b}g^*v_i)
      =e_R\operatorname{ord}_P(v_i),\qquad i=1,2.
\]
Thus neither $v_i$ has a pole outside $B$. A sufficiently large $a$ therefore makes
both sections regular in $kN+nQ_0+aB$, which again belongs to
\eqref{g:multiplicative-collection}. Their ratio lies in
$\mathbb C(Y')$, a contradiction.

The new generic section, multiplied by the Jacobian section, supplies
a nonzero adjoint section at every positive multiple of its exponent;
the Jacobian supplies one at zero. Combining these lower bounds with
the upper bound proves all required ranks after enlarging the fixed
divisibility integer. Finally replace the normal Stein base $Y'$ by a
smooth projective model and resolve the composite map. Its function
field is unchanged, and generic sections and Jacobian-twisted ranks
are unchanged by this birational replacement, as in
Section~\ref{g:cf:generic-conditions}. This completes the contraction.
\end{proof}

\section{Reduced fibers and descent by a finite flat norm}
\label{g:reduced-flat}

The limiting fiber norm in Lemma~\ref{g:moments} is defined first on
an open subset of the base.  To extend it as a bounded metric, we need a
line whose rational section has the same growth at the boundary.  This
section constructs that line after an alteration with reduced fibers,
then returns it to the original function field by a finite flat norm.
The alteration is used to control a section on every fiber; the final
base is birational to the original base.

We begin with an algebraic construction that explains the role of
reduced fibers.  All schemes in the next two lemmas are of finite type
over \(\mathbb C\).  A modification of an integral base means a proper
birational morphism, and all families over a modified base are pulled
back.

\begin{lemma}[Extending a generic line of sections]
\label{g:flat-section-line}
Let \(g:W\to T\) be a projective flat morphism with geometrically reduced
fibers, where \(T\) is a smooth integral projective variety.  Let \(P\)
be a line bundle on \(W\), and let
\(0\ne\tau\in H^0(W_{\eta},P_{\eta})\), where
\(\eta=\operatorname{Spec}\mathbb C(T)\).
After a smooth projective modification of \(T\), there are a line bundle
\(J\), a rational section \(\eta_J\) of \(J\), and a homomorphism
\[
                 a:g^*J\longrightarrow P
\]
such that \(a(g^*\eta_J)=\tau\) generically and
\(a|_{W_t}\ne0\) for every \(t\in T\).
\end{lemma}

\begin{proof}
Choose a sufficiently high
power \(A\) of a relatively ample line bundle so that
\[
 E_A=g_*A,\qquad F_A=g_*(A\otimes P)
\]
are vector bundles, their formation commutes with arbitrary base change,
and \(g^*E_A\to A\) is surjective.  Relative Serre vanishing and flatness
supply these assertions simultaneously.  Multiplication by \(\tau\)
defines a nonzero element of
\(\operatorname{Hom}(E_A,F_A)_{\eta}\).
Take the closure of the associated rational section of the projective
bundle of lines in \(\operatorname{Hom}(E_A,F_A)\), and resolve this
closure.  The resulting base modification carries a line subbundle
\[
                 J\hookrightarrow\operatorname{Hom}(E_A,F_A).
\]
In particular its map to each fiber of this Hom bundle is nonzero.

The induced map
\(g^*J\otimes g^*E_A\to A\otimes P\) factors through
\(g^*J\otimes A\).  Indeed, the kernel of
\(g^*E_A\to A\) is locally free.  Its composite with this map vanishes
over the generic base point, where the map is multiplication by
\(\tau\).  It therefore vanishes everywhere: flatness over the integral
base makes the generic part schematically dense, and the target is
locally free.  Dividing by \(A\) gives \(a:g^*J\to P\).
If its restriction to some fiber were zero, the corresponding map
\((E_A)_t\to(F_A)_t\) would be zero, contrary to the line subbundle
construction.  The prescribed generic equality defines \(\eta_J\).

\end{proof}

\begin{lemma}[The fiber norm has a positive lower bound]
\label{g:flat-supremum}
In the situation of Lemma~\ref{g:flat-section-line}, give \(J\) and
\(P\) continuous positive Hermitian metrics.  For
\(a:g^*J\to P\), put
\[
                  c(t)=\sup_{w\in W_t}|a(w)|^2.
\]
There are constants \(0<c_-\le c_+<\infty\) with
\(c_-\le c(t)\le c_+\) for all \(t\in T\).
\end{lemma}

\begin{proof}
Properness gives the upper bound.  On a reduced fiber, a nonzero section
of a line bundle is nonzero at a point in its smooth locus.  Choose such
a point over \(t\).  Flatness and smoothness of the fiber at that point
make \(g\) smooth there.  After shrinking, a local analytic section of
\(g\), together with continuity of the metric, supplies a positive lower
bound on a neighborhood of \(t\).  Finitely many such neighborhoods
cover the compact base.  Reducedness is used precisely here: a nonzero
section supported only in the nilpotent structure of a fiber would not
supply a point at which its norm is positive.
\end{proof}

\begin{remark}[Proper flat families]
The conclusion of Lemma~\ref{g:flat-section-line} also holds when
\(g\) is proper and flat without being projective. To prove this
variant, replace the Hom bundle construction by the complex computing
proper flat cohomology. The complex \(Rg_*P\) is perfect and commutes
with derived base change \cite[Tag~0DJT]{StacksProject}. On the smooth
projective base it has a global bounded vector bundle representative,
by the resolution property \cite[Tags~0F87 and 0F8F]{StacksProject}.
It can be chosen in degrees \(0,\ldots,N\): the derived fiber
cohomology has no negative degrees, so the negative part of a bounded
representative is successively removed by locally split injections and
their vector bundle cokernels. Write the resulting complex as
\[
              K^0\xrightarrow{d^0}K^1\longrightarrow\cdots.
\]
The vector \(\tau\) is a nonzero element of
\(\ker(d^0_{\eta})\subset K^0_{\eta}\). Modify the base to extend its
line to a subbundle \(J\subset K^0\). The composite \(J\to K^1\)
vanishes generically and hence everywhere. Thus
\(J\to\ker d^0=g_*P\) gives the required evaluation map.
On each fiber, the subbundle gives a nonzero vector in
\(\ker d^0_t=H^0(W_t,P_t)\); there is no incoming differential in
negative degree. This proves nonvanishing on every fiber. The
projective construction suffices for the transfer below.
\end{remark}

\subsection{The norm line on the original field}

We now use these lemmas to construct the two base metrics and the map
that will return sections to the source.  The anticanonical metric and
the bounded metric in the statement are separate data.

\begin{proposition}[Reduced-flat transfer]
\label{g:reduced-flat-transfer}
Let \(X,V,Y\) be smooth integral projective varieties, let
\(\pi:V\to X\) be birational, and let \(f:V\to Y\) be surjective with
connected fibers.  Let \(D\) be a reduced effective divisor on \(X\),
let \(c\ge0\) be an integer, and let \(L\) be a line bundle.
Suppose that \(L\) has a semipositive metric \(h_L\) with locally bounded
weights and that \(M=-K_X+cD\) has a semipositive singular metric
\(h_M\).  In local coordinates and a compatible frame of \(M\), write
\[
              d\mu=|d|^{2c}e^{-\varphi_M}\,dV_X,
\]
where \(d\) is an equation of \(D\), and assume \(\mu(X)<\infty\).
Put \(E=K_{V/X}\).  Assume that for some integer \(a>0\) there is a
nonzero rational section \(\tau\) of \(a\pi^*L\), regular over a dense
base open subset, and that
\begin{equation}
 \operatorname{rk} f_*\mathcal O_V\bigl(E+c\pi^*D+pa\pi^*L\bigr)=1
                \qquad\text{for every integer }p\ge0.
 \label{g:flat-adjoint-ranks}
\end{equation}
If \(\dim Y>0\), there are a smooth projective modification
\(Y'\to Y\), a line bundle \(B\) on \(Y'\), a reduced effective divisor
\(D'\), and integers \(d_0>0\), \(c'\ge0\) with the following properties.
The line \(B\) has a semipositive metric with locally bounded weights.
The line \(-K_{Y'}+c'D'\) has a semipositive metric inducing a finite
measure.  On a smooth resolution
\(\pi':V'\to X\), \(f':V'\to Y'\) of the main transform there is a
nonzero regular homomorphism
\begin{equation}
                  (f')^*B\longrightarrow d_0a(\pi')^*L,
 \label{g:flat-output-map}
\end{equation}
and
\begin{equation}
       \operatorname{Supp}\bigl(\pi'_* (f')^*D'\bigr)
                            \subseteq\operatorname{Supp}D.
 \label{g:flat-pole-support}
\end{equation}
Consequently, for every \(m>0\), a nonzero rational section of \(mB\)
with poles only on \(D'\) gives a nonzero rational section of
\(md_0aL\) with poles only on \(D\).

There is also a smooth invariant version.  In that version \(D=0\),
both metrics are smooth, an algebraic torus acts compatibly on \(X,V\)
and the line bundles, its action on \(Y\) is trivial, and \(\tau\) is
invariant.  A Zariski dense compact subgroup preserves the metrics.
Replace the ranks in \eqref{g:flat-adjoint-ranks} by the ranks of the
invariant parts, still for every \(p\ge0\).  Then the transferred section
on \(X\) is invariant and regular.  Anticanonical bundles in this version
carry their natural linearizations.
\end{proposition}

\begin{proof}
On a dense smooth base open, use coordinates \(y\) and write \(I_p(y)\)
for the density of
\[
              f_*\bigl(|\tau|_{(\pi^*h_L)^a}^{2p}\,\pi^*\mu\bigr).
\]
The pullback measure includes the Jacobian of \(\pi\).  The adjoint
section is \(s_Es_{c\pi^*D}\tau^p\); its fiberwise squared norm is
\(I_p\), after trivializing the canonical base factor.
Lemma~\ref{g:moments} applies to
\eqref{g:flat-adjoint-ranks}, including \(p=0\).  It gives psh
representatives of \(-\log I_0\) and
\begin{equation}
 u=-\log S,\qquad
 S(y)=\operatorname*{ess\,sup}_{V_y}|\tau|^2_{(\pi^*h_L)^a}.
 \label{g:flat-limit-norm}
\end{equation}
The essential supremum uses the fiber density of \(\pi^*\mu\).
For almost every smooth fiber that density is positive almost everywhere.
Since \(h_L\) has bounded weights, \(S\) is comparable to the ordinary
supremum computed with a fixed smooth reference metric.  The comparison
constants are uniform on the compact total space wherever the rational
section is defined. Fix this smooth reference metric on \(L\), and use
its pullbacks on every model below. In the invariant case the same lemma uses
the holomorphic orthogonal projection to the invariant direct summand
for the smooth coefficient metric.  Positivity of that quotient metric
is the additional input required in this case.

\emph{Constructing a line on an alteration.}
Apply weak semistable reduction of Abramovich--Karu to \(f\)
\cite[Theorem~0.3 and Definition~0.1, arXiv version~1]{AK2000}.  After a projective generically finite
alteration \(T\to Y\) and modifications, it gives a projective family
\(g:W\to T\) with geometrically reduced fibers, with \(T\) smooth
projective, and a morphism \(W\to V\) over \(Y\).  Its equidimensional
toroidal source is Cohen--Macaulay, so it is flat over the smooth base.
The source \(W\) need not be smooth.  We keep this flat family when
making further base modifications; resolving its source would not
preserve the flatness needed below.  The generic fiber of \(f\) is
geometrically integral, because it is smooth and connected in
characteristic zero, as required for this application of reduction.
The two families fit into the diagram
\[
 \begin{array}{ccc}
 W&\longrightarrow&V\\
 {\scriptstyle g}\downarrow&&\downarrow{\scriptstyle f}\\
 T&\longrightarrow&Y.
 \end{array}
\]

Let \(P\) be the pullback to \(W\) of \(a\pi^*L\).
Lemma~\ref{g:flat-section-line} gives \(J\), \(\eta_J\), and
\(g^*J\to P\) on a further modified base.  The family remains flat with
reduced fibers.  On the generic base open the map \(W\to V\) is
surjective on the fibers under consideration. Give \(J\) a continuous
positive reference metric and let \(c(t)\) be the supremum norm of the
map \(g^*J\to P\), measured in these reference metrics. Since this map
sends \(g^*\eta_J\) to \(\tau\), we have
\[
 \sup_{W_t}|\tau|_{\mathrm{ref}}^2
     =|\eta_J(t)|_{\mathrm{ref}}^2c(t),
 \qquad c_-\leq c(t)\leq c_+
\]
by Lemma~\ref{g:flat-supremum}. Surjectivity onto the original fiber
identifies the left side with \(\sup_{V_y}|\tau|_{\mathrm{ref}}^2\).
Comparison with the bounded metric \(h_L\) therefore gives
\begin{equation}
                  |\eta_J(t)|^2\asymp S(y),\qquad t\mapsto y,
 \label{g:flat-upstairs-comparison}
\end{equation}
for almost every parameter in that open.  Here and below comparison
means upper and lower positive constant bounds in local reference
metrics, with constants that remain bounded toward the omitted locus.
The fixed reference metrics are pulled back under all subsequent base
changes, and the same bounds on \(c(t)\) continue to hold.
No rational connectedness of \(T\) is asserted or needed.

\emph{Returning to the original function field.}
By flattening by modification, choose a smooth projective modification
\(Y'\to Y\) on which both the main transform \(V_0\to Y'\) of \(V\)
and the main transform \(q:T'\to Y'\) of \(T\) are flat
\cite{RaynaudGruson1971}.  The second map is finite: it is projective,
flat and has zero-dimensional fibers.  Its degree
\[
                  d_0=\operatorname{rk}q_*\mathcal O_{T'}
\]
is positive and constant.  The first map is equidimensional of the
original relative dimension.  Subsequent smooth modifications of the
base preserve these properties by flat base change.  The main
transforms map to \(T\) and \(V\), respectively, so \(J\) and its
rational section pull back to \(T'\).  The base \(Y'\) still has the
function field \(\mathbb C(Y)\).
The resulting diagram is
\[
 \begin{array}{ccccc}
 T'&\xrightarrow{\ q\ }&Y'&\longleftarrow&V_0\\
 \downarrow&&\downarrow&&\downarrow\\
 T&\longrightarrow&Y&\longleftarrow&V.
 \end{array}
\]
Here both maps to \(Y'\) are flat, and \(q\) is finite. The line
\(J\) lives on \(T'\); its norm will live on \(Y'\).

Define the actual norm line and its rational section by
\begin{equation}
 B=\operatorname{Nm}_q(J)
   :=\det(q_*J)\otimes\det(q_*\mathcal O_{T'})^{-1},
 \qquad \beta=\operatorname{Nm}_q(\eta_J).
 \label{g:flat-norm-definition}
\end{equation}
Both direct images in this formula are locally free of rank \(d_0\).
To describe the section, view multiplication by \(\eta_J\) as a rational
map \(q_*\mathcal O_{T'}\to q_*J\) and take its determinant.
Locally a line bundle can be trivialized near a finite inverse fiber;
changing the frame by a unit changes its norm frame by the determinant
of multiplication by that unit.  This proves the compatibility of
\eqref{g:flat-norm-definition} with changes of frame and arbitrary base
change.  The construction uses finite local freeness, and does not
require \(T'\) to be smooth or \(q\) to be unramified
\cite[Tags~0BUT, 0BCX and 0BD2]{StacksProject}.

Over the generically unramified locus the norm is the product over the
\(d_0\) sheets.  Applying \eqref{g:flat-upstairs-comparison} on those
sheets gives
\begin{equation}
                         |\beta(y)|^2\asymp S(y)^{d_0}.
 \label{g:flat-norm-comparison}
\end{equation}
These constants remain bounded at the branch locus.  Indeed, choose a
frame for \(J\) near the finite inverse fiber.  Its continuous reference
norm is bounded above and away from zero there, and its norm is a
nonvanishing frame of \(B\).  Products of the finitely many reference
norms are therefore comparable to any fixed reference norm on \(B\).
This also explains why the exponent is multiplied by \(d_0\).

Define a metric on \(B\), initially where \(\beta\) is a nonzero regular
frame, by \(|\beta|^2=S^{d_0}\).  If \(\beta=b e\) in a local
holomorphic frame \(e\), its weight is
\[
                    \psi=d_0u+\log|b|^2.
\]
It is psh on that open.  Equation~\eqref{g:flat-norm-comparison} bounds
\(\psi\) above and below locally even toward the omitted locus.  The
bounds hold for the psh representatives, by subaveraging and upper
semicontinuity.  Removable singularities across analytic sets therefore
extend \(\psi\) to a locally bounded psh weight on all of \(Y'\).

Resolve \(V_0\) to \(V'\), with maps \(\pi'\) and \(f'\) as in the
statement.  The rule
\begin{equation}
                          \beta\longmapsto\tau^{d_0}
 \label{g:flat-section-rule}
\end{equation}
defines a rational homomorphism in \eqref{g:flat-output-map} over the
original field. It extends regularly. Indeed, write
\(\beta=b e_B\) and \(\tau=t e_L^{\otimes a}\) in local holomorphic
frames, with the fixed reference norms bounded above and below. The
coefficient of this map is \(t^{d_0}/(b\circ f')\). On almost every
fiber in the good open, \eqref{g:flat-norm-comparison} and the definition
of the essential supremum give
\[
 \left|\frac{t^{d_0}}{b\circ f'}\right|^2
 \leq C\frac{S(f'(\,\cdot\,))^{d_0}}
                  {|\beta\circ f'|_{\mathrm{ref}}^2}
 \leq C'.
\]
The bound first holds almost everywhere on these fibers and then
everywhere on their good locus by continuity of the rational coefficient.
A rational coefficient with a pole along a prime divisor would be
unbounded on open subsets approaching its general point; those subsets
meet these fibers.  This excludes every divisorial pole, and normality
extends the map across codimension two.  In the invariant case the rule
\eqref{g:flat-section-rule} is invariant as a rational identity on \(X\).
The auxiliary alterations and final resolution need not be equivariant.

We have constructed the bounded base metric and its map back to \(X\).
On the new base, \(S\) denotes the same generic fiber supremum, whereas
\(I_0\) now denotes the density of \((f')_*(\pi')^*\mu\) in the chosen
coordinates on \(Y'\).  The latter transforms as a volume density,
including the base Jacobian.  Off the exceptional locus of the base
modification, this changes \(-\log I_0\) by a pluriharmonic coordinate
term and retains its plurisubharmonicity.
It remains to extend the volume metric and ensure that all its correction
primes disappear, or land in \(D\), upon return to \(X\).

\emph{The correction divisor and every component of its pullback.}
Among the boundary prime divisors of the good open in \(Y'\), let \(D'\)
be the reduced sum of those primes \(Q\) for which every prime of
\(V'\) dominating \(Q\) is either \(\pi'\)-exceptional or maps into
\(D\).  This definition first addresses components that dominate base
primes.  It gives \eqref{g:flat-pole-support} for \emph{all} components
of the total pullback because of the equidimensional model \(V_0\).
To see this, put \(n=\dim X\) and \(r=n-\dim Y'\).  If a prime
\(R\subset V'\) maps into a set \(A\subset Y'\) of codimension at
least two, then
\[
       \dim(V_0\times_{Y'}A)\le\dim A+r\le n-2.
\]
The image of \(R\) in \(V_0\), hence in \(X\), has dimension at most
\(n-2\), since \(V_0\to X\) is a morphism.  Thus \(R\) is
\(\pi'\)-exceptional.  Every other component of \((f')^*D'\) dominates
a component of \(D'\) and is covered by its definition.  This proves
\eqref{g:flat-pole-support}, including the divisors introduced by the
last resolution.

At a general point of a boundary prime \(Q\), choose a component
\(R\subset(f')^*Q\) dominating it, and write
\[
 k=\operatorname{ord}_R(f')^*Q>0,\qquad
 b=\operatorname{ord}_R\bigl(K_{V'/X}+c(\pi')^*D\bigr)\ge0.
\]
The psh weight of \(h_M\) is locally bounded above.  Hence the pullback
measure is bounded below by a positive constant times
\(|z_1|^{2b}\) times coordinate volume near a general point of \(R\).
Coordinates can be chosen there with \(t_1=z_1^k\), where \(Q=(t_1=0)\),
and with the other base coordinates among the remaining source
coordinates.  Integration in a fixed small fiber polydisc and change of
variables give
\begin{equation}
                    I_0(t)\ge C|t_1|^{2((b+1)/k-1)}
 \label{g:flat-volume-lower}
\end{equation}
almost everywhere near a general point of \(Q\).  If \(Q\not\subset D'\),
choose \(R\) nonexceptional over \(X\) and not mapping into \(D\).
Then \(b=0\); the exponent in \eqref{g:flat-volume-lower} is nonpositive,
so in particular \(I_0\) has a positive constant lower bound there.

There are only finitely many boundary primes.  Choose an integer
\(c'\ge0\) at least as large as each positive exponent
\((b+1)/k-1\) needed for a component of \(D'\).  If \(d'\) is an equation
of \(D'\), the psh function
\[
                     -\log I_0+c'\log|d'|^2
\]
is locally bounded above at the generic points of all boundary primes.
It extends across them, and then across the remaining analytic subset
of codimension at least two by the Hartogs extension property for psh
functions.  These are the weights of a semipositive metric on
\(-K_{Y'}+c'D'\).  Dividing its metric by the divisor metric of \(c'D'\)
recovers the measure with density \(I_0\).  It has finite mass because it
is the pushforward of the original finite measure.

Finally take a nonzero rational section of \(mB\) with poles on \(D'\),
pull it back, and apply the \(m\)-th power of
\eqref{g:flat-output-map}.  The only possible poles lie in
\((f')^*D'\).  Equation~\eqref{g:flat-pole-support} shows that its
pushdown has poles only on \(D\), and it is nonzero at the generic point.
In the smooth invariant case this rational section is invariant by
\eqref{g:flat-section-rule}; since \(D=0\), normality makes it a global
regular section on \(X\).
\end{proof}

\begin{remark}[Valuations of the norm]
\label{g:flat-norm-valuations}
The finite norm itself creates no polar prime away from the images of
polar primes upstairs.  At a prime \(Q\subset Y'\), pass to the discrete
valuation ring \(\mathcal O_{Y',Q}\) and normalize its finite generic
extension.  If \(v_i\) are the extension valuations and \(f_i\) their
residue degrees, then for a nonzero rational function \(z\),
\[
                \operatorname{ord}_Q\operatorname{Nm}(z)
                             =\sum_i f_i v_i(z).
\]
This is the determinant of multiplication by \(z\), computed by lengths
of lattices after clearing denominators.  Applied in local frames, it
gives the same formula for a normed rational section.  For the section
returned to \(X\), the stronger statement needed is
\eqref{g:flat-pole-support}: it also treats components lying over
codimension-two sets in the base.
\end{remark}

\begin{remark}[What is descended]
The determinant construction in \eqref{g:flat-norm-definition} descends
a line on a possibly ramified, nonsmooth finite flat cover and changes
\(a\) to \(d_0a\).  The finite-group quotient construction in
\cite[Subsection ``Descending the alteration and its section'']{CCohomologicalTransfer}
instead first kills stabilizer actions by
a tensor power and then uses a function-field norm on a normal target.
These are different model constructions.  They are also distinct from
the finite \emph{\'etale} section norm of a line already pulled back from
the original variety.  None of the three permits induction on the
altered base \(T\) here; the induction below takes place on \(Y'\),
whose function field remains \(\mathbb C(Y)\).
\end{remark}

\subsection{An approximation proof of the slope input}

For the rationally connected application of the reduced-flat
construction, the required slope inequality also has an approximation
proof.  Its distinctive step is to convert finite weighted volume into
a klt boundary on a resolution.  We give that step, rather than infer
it from the existence of the descended metric.

\begin{proposition}[Floor approximation for foliation determinants]
\label{g:reduced-flat-slope}
Let \(X\) be smooth projective, let \(D\) be a reduced effective divisor,
and let \(c\ge0\) be an integer. Suppose that \(M=-K_X+cD\) has a
semipositive singular metric with finite induced measure
\(|d|^{2c}e^{-\varphi}\,dV_X\). Every algebraically integrable saturated
foliation \(\mathcal F\subset T_X\) of positive rank and corank satisfies
\begin{equation}
                         K_{\mathcal F}+M
                          \quad\hbox{is pseudoeffective}.
 \label{g:flat-foliation-psef}
\end{equation}
\end{proposition}

\begin{proof}
Its leaves define a rational fibration.  Fix a sufficiently ample line
bundle \(H\), independent of the integer \(l\ge2\).  Nadel vanishing
and Castelnuovo--Mumford regularity make
\[
               \mathcal O_X(lM+H)\otimes\mathcal J(l\varphi)
\]
globally generated.  The local Bergman approximation inequality gives,
for its local base ideal generators \(g_j\),
\begin{equation}
                          e^{l\varphi}\le C_l\sum_j|g_j|^2.
 \label{g:flat-approximation}
\end{equation}
The fixed ample twist can be chosen by applying Nadel vanishing after
subtracting \(1,\ldots,\dim X\) times a fixed very ample line bundle.
The local inequality follows from point extension in the weighted
holomorphic \(L^2\) space; generators of the multiplier ideal represent
the singularities of that space.  These are the usual Bergman
approximation and global generation constructions of Demailly
\cite[Theorems~0.3 and 1.2]{Demailly2015}.

Let \(\rho:U\to X\) resolve the rational fibration and principalize this
base ideal, with its divisor \(R_l\), and let \(E_U=K_{U/X}\).  Take a
simultaneous log resolution with \(\rho^*D\).  Finite mass and
\eqref{g:flat-approximation} imply, along every prime on \(U\),
\begin{equation}
                         \lfloor R_l/l\rfloor
                                  \le E_U+c\rho^*D.
 \label{g:flat-floor-bound}
\end{equation}
Indeed the comparison makes
\(|s_{E_U+c\rho^*D}|^2|s_{R_l}|^{-2/l}\) locally integrable.
At a prime with respective orders \(b\) and \(r\), integrability says
\(b-r/l>-1\); since \(b\) is integral this is exactly
\(\lfloor r/l\rfloor\le b\).

Pull back the linear subsystem
\[
 H^0\bigl(X,\mathcal O_X(lM+H)\otimes\mathcal J(l\varphi)\bigr)
                  \subseteq H^0(X,lM+H)
\]
to \(U\), and remove its fixed divisor \(R_l\).  Let \(G_l\) be a
general member of the resulting basepoint-free system, and put
\[
           \Delta_l=R_l/l-\lfloor R_l/l\rfloor+G_l/l.
\]
Bertini and the log resolution make \(\Delta_l\) an effective klt
boundary, also klt on the general fiber of the resolved fibration.
Moreover,
\begin{equation}
 K_U+\Delta_l\sim_{\mathbb Q}
       E_U+c\rho^*D-\lfloor R_l/l\rfloor+\rho^*H/l.
 \label{g:flat-adjoint-approximation}
\end{equation}
By \eqref{g:flat-floor-bound}, a positive divisible multiple of the
right side restricts to an effective divisor on the general fiber.
One may take a general effective representative of \(H\), which does
not contain that fiber.

The relative pseudoeffectivity theorem for an algebraically integrable
foliation now gives
\(K_{\mathcal F_U}+\Delta_l\) pseudoeffective on \(U\), where
\(\mathcal F_U\) is the induced foliation
\cite[Proposition~3.1, arXiv version~4]{Ou2023}; its relative positivity
input is due to Druel \cite[Proposition~4.1]{Druel2017}.
The hypotheses used are the effective boundary, log canonicity of its
general-fiber restriction, and the general-fiber nonvanishing in
\eqref{g:flat-adjoint-approximation}.  Push down to \(X\), and add back
the effective divisor \(\rho_*\lfloor R_l/l\rfloor\).  Since foliation
determinants agree away from the exceptional locus, this proves that
\(K_{\mathcal F}+M+H/l\) is pseudoeffective.  Its limit as
\(l\to\infty\) proves \eqref{g:flat-foliation-psef}.

\end{proof}

Apply Lemma~\ref{g:bt:slope-principle} with \(P=cD\).
For every nonzero movable class \(\alpha\) on a positive-dimensional
\(X\) with \(D\cdot\alpha=0\), the preceding proposition gives
\(\mu_{\alpha,\min}(T_X)\ge0\). Movable tensor semistability then gives
nonpositive maximal slope for every cotangent tensor power and its
exterior powers. Thus every coherent subsheaf of
\((\Omega_X^1)^{\otimes k}\), or of \(\Omega_X^p\), has nonpositive
\(\alpha\)-degree. For degree zero this follows directly from the
inclusion in \(\mathcal O_X\). The floor construction above is an
alternative proof of the foliation input to the general weighted
estimate, Proposition~\ref{g:bt:weighted-slope}.

\subsection{An alternative proof of allowed-pole nonvanishing}
\label{g:rc-allowed-poles}

The reduced-flat transfer preserves the original base field.  This is
what permits rational connectedness to pass to the base on which we
induct.  The following proof records how the common field and
determinant constructions meet every rank hypothesis of
Proposition~\ref{g:reduced-flat-transfer}.

We prove the rationally connected case of Theorem~\ref{g:allowed}
using reduced-flat transfer. First reduce its arbitrary effective
integral boundary \(D_0\) to a multiple of a reduced divisor. Put
\(D=\operatorname{Supp}(D_0)_{\mathrm{red}}\), choose an integer
\(c\geq0\) with \(cD-D_0\geq0\), and tensor the metric on
\(-K_X+D_0\) with the canonical divisor metric of \(cD-D_0\).
The new metric on \(-K_X+cD\) is semipositive, and its distinguished
measure is unchanged: the squared divisor factor cancels the added
metric factor. A rational section of \(mL\), \(m>0\), with poles only
on \(D\) belongs to \(H^0(X,mL+rD_0)\) for some sufficiently large
\(r\geq0\).
The zero boundary is included by taking \(D=0\) and \(c=0\).

It therefore suffices to take \(X\) smooth projective and rationally
connected, \(D\) reduced effective, and \(L\) a line bundle that
has a semipositive metric with locally bounded weights. Suppose
\(-K_X+cD\) has a semipositive metric with finite induced measure for
an integer \(c\ge0\). We seek a nonzero section of \(mL+rD\) with
integers \(m>0\) and \(r\ge0\). A point is included.

\begin{proof}
Induct on \(n=\dim X\); the assertion for a point is immediate.
Suppose the conclusion fails.  The class \(L+D\) is pseudoeffective
and not big, so cone duality gives a nonzero movable class \(\alpha\)
with
\[
                       L\cdot\alpha=D\cdot\alpha=0
\]
\cite{BDPP2013}.  Proposition~\ref{g:reduced-flat-slope} and the following slope argument apply.
Rational connectedness gives \(H^i(X,\mathcal O_X)=0\) for \(i>0\), so
\(\chi(X,K_X)=(-1)^n\ne0\).  Riemann--Roch makes
\(\chi(X,K_X+mL)\) a nonzero polynomial.  Thus one cohomological degree
occurs with nonzero cohomology for an unbounded sequence of positive
\(m\).  Hard Lefschetz with multiplier ideals supplies twisted
holomorphic forms for that sequence; the multiplier ideals are trivial
because the metric weights on \(mL\) are bounded
\cite{DPS2001}.

Apply the determinant construction of Lemma~\ref{g:determinant}
to these forms.  Its fixed determinant maps nontrivially to a
cotangent tensor power, so its \(\alpha\)-degree is nonpositive by the slope consequence of
Proposition~\ref{g:reduced-flat-slope}.  We obtain a fixed line bundle
\(B_0\) and effective divisors
\begin{equation}
                 N_j\sim jL+B_0,\qquad \alpha\cdot N_j=0
 \label{g:flat-determinant-sequence}
\end{equation}
for an unbounded set of positive integers \(j\).

Let \(K_\alpha\subset\mathbb C(X)\) be the field of functions whose
polar divisors have \(\alpha\)-degree zero.  By
Lemma~\ref{g:cf:null-field}, it is relatively algebraically closed,
its transcendence degree is less than \(n\), and it is represented by
a rational fibration with connected generic fiber.  Resolve it to
\(\pi:V\to X\), \(f:V\to Y\), with \(Y\) smooth projective.
For a general sufficiently ample divisor \(H_Y\) on \(Y\),
\begin{equation}
                         \alpha\cdot\pi_*f^*H_Y=0.
 \label{g:flat-base-degree}
\end{equation}
Indeed a general ratio from \(|H_Y|\) belongs to \(K_\alpha\); the two
general pulled-back divisors have no common component, so its pole
divisor represents the indicated class.

The affine interpolation of Lemma~\ref{g:cf:horizontal-slope} applied
to \eqref{g:flat-determinant-sequence} produces an integer \(a>0\) and
a nonzero rational section \(\tau\) of \(a\pi^*L\) that is regular over
a dense base open.  Explicitly, the linear equivalence
\[
                (k-i)N_j\sim(k-j)N_i+(j-i)N_k
                    \qquad(i<j<k)
\]
has its defining function in \(K_\alpha\).  Consequently the
coefficients of \(\pi^*N_j\) at every horizontal prime are affine in
\(j\); nonnegativity at unbounded \(j\) makes each slope nonnegative.
A difference of two such divisors has no negative horizontal part and
gives \(\tau\).  If \(Y\) is a point, this already gives a global
section of a positive multiple of \(L\).

Suppose therefore that \(\dim Y>0\), and put \(E=K_{V/X}\).
For every integer \(p\ge0\),
\(s_Es_{c\pi^*D}\tau^p\) is a nonzero generic-fiber section of
\(E+c\pi^*D+pa\pi^*L\).  Any finite set of other generic-fiber
sections becomes global after a common sufficiently ample base twist.
The pushdowns of their effective divisors have \(\alpha\)-degree zero,
by \eqref{g:flat-base-degree}, \(\pi_*E=0\), and
\(\alpha L=\alpha D=0\).  Ratios of two such sections therefore
belong to \(K_\alpha\), proving
\eqref{g:flat-adjoint-ranks} for every \(p\), including zero.
This argument establishes generic ranks and requires no simultaneous
base-change assertion for infinitely many sheaves.

Proposition~\ref{g:reduced-flat-transfer} produces \(Y',B,D'\).
The variety \(Y'\) is rationally dominated by \(X\), so it is
rationally connected, and its dimension is smaller than \(n\).
Induction gives a section of \(mB+rD'\) for some \(m>0\), \(r\ge0\).
Transfer returns a nonzero rational section of \(md_0aL\) with poles
only on \(D\), the required contradiction.
\end{proof}

\subsection{An alternative proof for the invariant quotient}
\label{g:flat-invariant-application}

We now prove the rationally connected case of Theorem~\ref{rb:invariant}
by the same reduced-flat method. Let \(F\) be smooth projective and
rationally connected, with \(L=-K_F\) smoothly semipositive, and let an
algebraic torus \(T\) act with the natural anticanonical linearization.
Average a smooth semipositive metric over a compact real form of \(T\).
We seek \(H^0(F,mL)^T\ne0\) for some integer \(m>0\).

\begin{proof}
For the trivial action this is the argument of Section~\ref{g:rc-allowed-poles} with
\(D=0\).  Otherwise take a rational quotient and an equivariant
resolution \(\pi:V\to F\), \(f:V\to Y\).  Its invariant field is
\(\mathbb C(Y)\), and the generic fiber is an orbit closure.  In
particular it is connected.  On the smooth quotient locus, choose
\(\dim V-\dim Y\) generically independent vector fields from the torus
action and wedge them.  Commutativity makes their wedge an invariant
section of \(-K_{V/Y}\).  Tensor it with a rational anticanonical frame
from \(Y\), then multiply by the Jacobian section \(s_E\), where
\(E=K_{V/F}\).  This produces an invariant rational section
\(\tau\) of \(\pi^*L\), regular over a dense base open.

For every \(p\ge0\), the invariant generic-fiber adjoint space for
\(E+p\pi^*L\) contains \(s_E\tau^p\).  Two nonzero invariant sections
have an invariant ratio, hence a ratio in \(\mathbb C(Y)\).
Its invariant rank is exactly one, also at \(p=0\).
Apply the smooth invariant part of
Proposition~\ref{g:reduced-flat-transfer}, using the given smooth metric
for both \(L\) and \(-K_F\).  If \(Y\) is a point, \(\tau\) is already
a global invariant section.  Otherwise the resulting original-field
base \(Y'\) is rationally connected; apply
the argument of Section~\ref{g:rc-allowed-poles} there and transfer its section.
The pole-support statement and \(D=0\) make the resulting invariant
section regular on \(F\).
\end{proof}

\bibliographystyle{amsplain}
\bibliography{references}
\end{document}
